% 高一126_建青实验_期中.tex
%   —— 高一上数学"2026-2027 年建青实验高一上期中"（试卷版/答案版）
% 试卷版：xelatex 高一126_建青实验_期中.tex
% 答案版：xelatex "\def\isanswer{1}% 高一126_建青实验_期中.tex
%   —— 高一上数学"2026-2027 年建青实验高一上期中"（试卷版/答案版）
% 试卷版：xelatex 高一126_建青实验_期中.tex
% 答案版：xelatex "\def\isanswer{1}% 高一126_建青实验_期中.tex
%   —— 高一上数学"2026-2027 年建青实验高一上期中"（试卷版/答案版）
% 试卷版：xelatex 高一126_建青实验_期中.tex
% 答案版：xelatex "\def\isanswer{1}% 高一126_建青实验_期中.tex
%   —— 高一上数学"2026-2027 年建青实验高一上期中"（试卷版/答案版）
% 试卷版：xelatex 高一126_建青实验_期中.tex
% 答案版：xelatex "\def\isanswer{1}\input{高一126_建青实验_期中.tex}"
% 紧凑版：xelatex "\def\iscompact{1}\input{高一126_建青实验_期中.tex}"
%
% 原始材料：微信公众号"魔都数学加油站"文章，作者破晓，2026-11-11 18:35 发布；
% 连载"27学年高一126"，原文标题"27学年高一126——2026-2027年上海市建青实验高一上期中"；
% 原文链接：https://mp.weixin.qq.com/s/HpxpBJu4d_wOnUCIbZSQiA。
% 正文为 8 张 PNG 页（01、05、08 为 2560x3620，02、03、04、06、07 为 4762x6735），
% 题目下印有【解析】或【答案】，为讲评版。题目、解析均据原图转录；
% 卷面粉红色水印"魔都数学加油站"及公众号页脚未录入。全卷无插图。
%
% 卷首标题：原卷印作"2026-2027 年建青实验高一上期中"（公众号标题中的"上海市"
% 卷面标题行没有，本稿照卷面），年份区间按全稿惯例排 en dash；原卷未印副标题，
% 本卷为期中卷、跨集合与常用逻辑用语／不等式／函数三章，按题目内容标为
% "集合与常用逻辑用语、不等式、函数"（与 125 同为期中卷，口径相同）：
% 集合与逻辑见 Q5、Q8、Q9、Q11、Q13、Q14、Q16、Q19，不等式见 Q6、Q7、Q10、
% Q17、Q20、Q21，函数（含指数对数、二次函数零点）见 Q3、Q4、Q12、Q15、Q18。
%
% 与原卷的排版差异：
%   ① 原文从第 3 题起，填空题编号为 3—12、选择题为 13—16、解答题为 17—21，本稿保留原题号；
%   ② 原卷每题下直接印【解析】／【答案】，本稿按全稿惯例置于答案版；第 4、5 题只印
%      【答案】，另提 \ans{}；第 13—16 题答案均在【解析】末句"故选 …"，题干括号内留空、
%      不另提 \ans{}；第 3 题印一行【解析】，按 \solbegin+\step 排；第 17 题解析保留
%      原卷评分点（照 高一54／64 成例排 \dots\dots N 分）；
%   ③ 原卷未印副标题，本稿按"知识模块"口径补出（见上"卷首标题"）；
%   ④ 原卷填空横线按全稿惯例排为 \blankline[4.4em]；选择题括号排 \mbox{（\quad）}、
%      选项段前按惯例加 \keepwithprev；第 14 题原卷括号另起一行，本稿并入题干行末；
%      解答题不留作答区（原卷为讲评版，与 109—125 同）；
%   ⑤ 原卷数集排斜体 $R$，按全稿惯例改 \R；绝对值排 \lvert…\rvert，集合描述竖线排 \mid；
%      补集记号照原卷上划线排 \overline{B}；不等号按全稿惯例排 \leqslant／\geqslant；
%      子题标号原卷一律用全角括号（1）（2）（3），本稿照录（\step 行内亦同，照 高一121 成例）。
%
% 判别式字形：第 10 题解析 $\Delta=a^2-4a<0$ 原卷排 $\Delta$（等腰、高），
%   照录 \Delta（与 高一124／125 同题排 $\triangle$（等边、扁）不同，逐处照本卷字形）。
%
% 子集符号：第 8 题题干与解析两处 $S\subset P$ 均印单线、无下划线，照录 \subset；
%   全卷未出现带下划线的 \subseteq。
%
% 原卷存疑照录：
%   第 10 题题干作 $ax^2-ax+1>0$，解析却作 $ax^2+ax+1>0$（一次项符号不一致），照录、未统一；
%   第 11 题解析"此时 $M=\dfrac32,1,\dfrac12$，满足"未用集合花括号（题干 $M$ 为集合），照录；
%   第 11 题解析"因为它们是集合 $M$ 中的元素，所以存在互异性"（"存在互异性"疑为
%     "满足互异性"之误），又"当 $x_1\neq x_2\neq x_3$ 时"连写两个不等号，均照录；
%   第 16 题解析两处"惟如此才能保证"（用"惟"字）及"符合定义的参数 $m$ 的值一定大于等于 1
%     或小于等于 0"的措辞，照录；
%   第 17 题 (1) 解析"当且仅当 $a=1,b=-1$ 时"用半角逗号、第 21 题 (3) 解析
%     "当 $a=\sqrt2$，$b=1$ 时取等号"用全角逗号，均照录；
%   第 19 题解析由 $\dfrac{x+4}{x+2}>2$ 得 "$2-\dfrac{x+4}{x+2}=\dfrac{x}{x+2}<0$"
%     （以 2 减左式，与常见的"左式减 2"次序相反），照录；
%   第 15 题选项 A、B 作"该方程…"、C、D 作"方程组…"，照录；
%   第 20 题题干作"招收代理商"，(1) 却作"招代理商后"，照录。
%
% 行数复核：逐题数 \step 条数，合计 91 条（Q3：1；Q4：0；Q5：0；Q6：3；Q7：2；
% Q8：5；Q9：5；Q10：4；Q11：8；Q12：3；Q13：3；Q14：2；Q15：3；Q16：9；
% Q17：7；Q18：5；Q19：9；Q20：13；Q21：9），与原卷印刷行逐题相等
% （cases 块作一个印刷行计）。第 4、5 题只印【答案】，提 \ans{}，不计入 \step 账；
% 四道选择题答案均在解析末句，括号内留空、不另提 \ans{}。
%
\documentclass[a4paper,11pt]{article}
\usepackage{common}

\begin{document}

\examtitlest[3]{2026--2027 年建青实验高一上期中}{集合与常用逻辑用语、不等式、函数}

\bigsec{一、填空题}

\begin{enumerate}
\setcounter{enumi}{2}

\item 已知 $\log_2 x=3$，则 $x=$\blankline[4.4em].

\ifanswer
\solbegin
\step{$x=2^3=8$.}
\solend
\fi

\item 将 $\sqrt{a^2\sqrt a}$（其中 $a>0$）化为有理数指数幂的形式为\blankline[4.4em].

\ifanswer
\ans{$a^{\frac54}$}
\fi

\item 用反证法证明命题：“已知 $m$、$n\in\R$，若 $m+n\leqslant5$，则 $m\leqslant2$ 或 $n\leqslant3$” 时应假设\blankline[4.4em].

\ifanswer
\ans{$m>2$ 且 $n>3$}
\fi

\item 若 $x\in\R$，则 $\lvert x+1\rvert+\lvert x-1\rvert$ 的最小值为\blankline[4.4em].

\ifanswer
\solbegin
\step{$\lvert x+1\rvert+\lvert x-1\rvert\geqslant\lvert(x+1)-(x-1)\rvert=2$，}
\step{当且仅当 $(x+1)(x-1)\leqslant0$，即 $-1\leqslant x\leqslant1$ 时取等号，}
\step{所以 $\lvert x+1\rvert+\lvert x-1\rvert$ 的最小值为 2.}
\solend
\fi

\item 已知 $x>1$，则 $x+\dfrac1{x-1}$ 的最小值是\blankline[4.4em].

\ifanswer
\solbegin
\step{$x-1>0$，则 $x+\dfrac1{x-1}=(x-1)+\dfrac1{x-1}+1\geqslant2\sqrt{(x-1)\cdot\dfrac1{x-1}}+1=3$，}
\step{当且仅当 $x=2$ 时取等号，故 $x+\dfrac1{x-1}$ 的最小值为 3.}
\solend
\fi

\item 已知集合 $P=\{x\mid x^2+x-6=0\}$，$S=\{x\mid ax+1=0\}$ 且 $S\subset P$，则 $a$ 的值为\blankline[4.4em].

\ifanswer
\solbegin
\step{集合 $P=\{-3,2\}$，集合 $S$ 中至多有一个元素，}
\step{若集合 $S$ 为空集，即 $a=0$ 时，显然满足条件 $S\subset P$，故 $a=0$.}
\step{若集合 $S$ 非空集，即 $a\neq0$，此时 $S=\left\{-\dfrac1a\right\}$，}
\step{若 $-\dfrac1a=-3$，则 $a=\dfrac13$，若 $-\dfrac1a=2$，则 $a=-\dfrac12$；}
\step{故 $a$ 的取值为 $0$、$\dfrac13$、$-\dfrac12$.}
\solend
\fi

\item 已知 $p:\begin{cases}a>4\\b>4\end{cases}$，$q:\begin{cases}a+b>8\\ab>16\end{cases}$，则 $p$ 是 $q$ 的\blankline[4.4em]条件.

\ifanswer
\solbegin
\step{若 $a>4$，$b>4$，则 $a+b>4+4=8$，$ab>4\times4=16$，}
\step{所以 $p$ 成立一定能推出 $q$ 成立，充分性成立，}
\step{取 $a=2$，$b=9$，$a+b=11>8$，$ab=18>16$，满足 $q$；}
\step{但是 $a=2<4$，不满足 $p$，所以由 $q$ 不能推出 $p$，必要性不成立，}
\step{综上所述，$p$ 是 $q$ 的充分不必要条件.}
\solend
\fi

\item 若关于 $x$ 的不等式 $ax^2-ax+1>0$ 对一切实数 $x$ 恒成立，则实数 $a$ 的取值范围是\blankline[4.4em].

\ifanswer
\solbegin
\step{当 $a=0$ 时，不等式即 $1>0$，满足条件.}
\step{当 $a\neq0$ 时，要使不等式 $ax^2+ax+1>0$ 对一切 $x\in\R$ 恒成立，}
\step{需 $\begin{cases}a>0\\\Delta=a^2-4a<0\end{cases}$，解得 $0<a<4$.}
\step{综上，实数 $a$ 的取值范围是 $[0,4)$.}
\solend
\fi

\item 已知集合 $M=\{x\mid (x-a)(x^2-ax+a-1)=0\}$ 各元素之和等于 3，则实数 $a=$\blankline[4.4em].

\ifanswer
\solbegin
\step{方程 $(x-a)(x^2-ax+a-1)=0$，化简得 $(x-a)(x-1)[x-(a-1)]=0$，}
\step{解得 $x_1=a$，$x_2=1$，$x_3=a-1$，}
\step{因为它们是集合 $M$ 中的元素，所以存在互异性.}
\step{当 $x_1\neq x_2\neq x_3$ 时，即 $a+1+a-1=3$，解得 $a=\dfrac32$，此时 $M=\dfrac32,1,\dfrac12$，满足.}
\step{当 $a=1$ 时，$1+0\neq3$，不合题意；}
\step{当 $a=a-1$ 时，不符合题意；}
\step{当 $1=a-1$ 时，$a=2$，此时 $M=\{2,1\}$ 满足题意.}
\step{故 $a$ 的取值为 $2$ 或 $\dfrac32$.}
\solend
\fi

\item 已知 $a>b>1$ 且 $a^2-6a-2m+1=0$，$b^2-6b-2m+1=0$，则实数 $m$ 的取值范围是\blankline[4.4em].

\ifanswer
\solbegin
\step{由题意得 $a$、$b$ 是函数 $f(x)=x^2-6x-2m+1$ 两个大于 1 的不同零点，}
\step{所以 $\begin{cases}f(1)>0\\f(3)<0\end{cases}$，即 $\begin{cases}-4-2m>0\\-8-2m<0\end{cases}$，解得 $-4<m<-2$，}
\step{故 $m$ 的取值范围为 $(-4,-2)$.}
\solend
\fi

\end{enumerate}

\bigsec{二、选择题}

\begin{enumerate}
\setcounter{enumi}{12}

\item “$x<2$” 是 “$x^2<4$” 的\mbox{（\quad）}条件

\keepwithprev
A. 充分不必要\qquad\qquad B. 必要不充分

C. 充要\qquad\qquad\qquad\quad D. 既不充分也不必要

\ifanswer
\solbegin
\step{解 $x^2<4$，得 $-2<x<2$，}
\step{由 $x<2$ 得不出 $-2<x<2$，而由 $-2<x<2$ 能得出 $x<2$，}
\step{所以 “$x<2$” 是 “$x^2<4$” 的必要不充分条件. 故选 $B$.}
\solend
\fi

\item 已知集合 $A=\{x\mid x<-2\}$，$B=\{x\mid x\geqslant a\}$，且 $A\cup B=\R$，则实数 $a$ 的取值范围是\mbox{（\quad）}

\keepwithprev
A. $a\leqslant-2$\qquad\qquad B. $a<-2$\qquad\qquad C. $a\geqslant-2$\qquad\qquad D. $a>-2$

\ifanswer
\solbegin
\step{集合 $A=\{x\mid x<-2\}$，$B=\{x\mid x\geqslant a\}$，且 $A\cup B=\R$，}
\step{则实数 $a$ 的取值范围是 $a\leqslant-2$. 故选 $A$.}
\solend
\fi

\item 下列关于 $x,y$ 的方程组 $\begin{cases}y=kx+1\\y=2kx+3\end{cases}$ 的解的说法中正确的是\mbox{（\quad）}

\keepwithprev
A. 该方程一定有唯一解\qquad\qquad B. 该方程没有解

C. $k\neq0$ 时，方程组有唯一解\qquad\qquad D. $k=0$ 时，方程组有无数组解

\ifanswer
\solbegin
\step{由题意得 $kx+1=2kx+3$，即 $kx=-2$，}
\step{当 $k=0$ 时，$0=-2$ 不成立，方程组无解；}
\step{当 $k\neq0$ 时，$x=-\dfrac2k$，方程组有唯一解. 故选 $C$.}
\solend
\fi

\item 设非空集合 $S=\{x\mid m\leqslant x\leqslant n\}$ 满足：当 $x\in S$ 时，有 $x^2\in S$. 给出如下三个结论：①若 $m=1$，则 $S=\{1\}$；②若 $m=-\dfrac12$，则 $\dfrac14\leqslant n\leqslant1$；③若 $n=\dfrac12$，则 $-\dfrac{\sqrt2}{2}\leqslant m\leqslant0$. 其中正确命题的个数为\mbox{（\quad）}

\keepwithprev
A. $3$\qquad\qquad B. $2$\qquad\qquad C. $1$\qquad\qquad D. $0$

\ifanswer
\solbegin
\step{由定义设非空集合 $S=\{x\mid m\leqslant x\leqslant n\}$ 满足当 $x\in S$ 时，有 $x^2\in S$，}
\step{符合定义的参数 $m$ 的值一定大于等于 1 或小于等于 0，}
\step{惟如此才能保证 $m\in S$ 时，有 $m^2\in S$ 即 $m^2\geqslant m$，}
\step{符合条件的 $n$ 的值一定大于等于 0，小于等于 1，}
\step{惟如此才能保证 $n\in S$ 时，有 $n^2\in S$ 即 $n^2\leqslant n$，对各个命题进行判断：}
\step{对于① $m=1$，$m^2=1\in S$，故必有 $\begin{cases}n^2\leqslant n\\n\geqslant1\end{cases}$，得 $n=1$，$S=\{1\}$，}
\step{② $m=-\dfrac12$，$m^2=\dfrac14\in S$，则 $\begin{cases}n^2\leqslant n\\\dfrac14\leqslant n\end{cases}$，解得 $\dfrac14\leqslant n\leqslant1$；}
\step{对于③，若 $n=\dfrac12$，则 $\begin{cases}\dfrac12\geqslant m\\m^2\geqslant m\\\dfrac12\geqslant m^2\end{cases}$，解得 $-\dfrac{\sqrt2}{2}\leqslant m\leqslant0$，}
\step{所以正确命题有 3 个. 故选 $A$.}
\solend
\fi

\end{enumerate}

\bigsec{三、解答题}

\begin{enumerate}
\setcounter{enumi}{16}

\item 已知 $a$、$b$ 是实数.

（1）求证：$a^2+b^2\geqslant2a-2b-2$，并指出等号成立的条件；

（2）若 $ab=1$，求 $a^2+4b^2$ 的最小值.

\ifanswer
\solbegin
\step{（1）因为 $a^2+b^2-(2a-2b-2)=a^2+b^2-2a+2b+2$\dots\dots2 分}
\step{$=(a-1)^2+(b+1)^2\geqslant0$\dots\dots4 分}
\step{所以 $a^2+b^2\geqslant2a-2b-2$.}
\step{当且仅当 $a=1,b=-1$ 时，不等式中等号成立.\dots\dots7 分}
\step{（2）$a^2+4b^2=a^2+(2b)^2\geqslant2\cdot a\cdot(2b)=4ab=4$，\dots\dots5 分}
\step{当且仅当 $a=2b$，即 $\begin{cases}a=\sqrt2\\b=\dfrac{\sqrt2}{2}\end{cases}$ 或 $\begin{cases}a=-\sqrt2\\b=-\dfrac{\sqrt2}{2}\end{cases}$ 时，不等式中等号成立，}
\step{所以 $a^2+4b^2$ 的最小值为 4.\dots\dots7 分}
\solend
\fi

\item 已知关于 $x$ 的方程 $x^2+(m-3)x+2=0$，其中 $m\in\R$，有两个实数根 $\alpha$、$\beta$.

（1）若 $m=8$，求 $\dfrac1\alpha+\dfrac1\beta$ 的值；

（2）若 $\alpha^2+\beta^2=\alpha\beta+3$，求实数 $m$ 的值.

\ifanswer
\solbegin
\step{（1）若 $m=8$，则关于 $x$ 的方程为 $x^2+5x+2=0$，}
\step{则 $\alpha+\beta=-5$，$\alpha\beta=2$，所以 $\dfrac1\alpha+\dfrac1\beta=\dfrac{\alpha+\beta}{\alpha\beta}=-\dfrac52$；}
\step{（2）由题意得 $\alpha+\beta=-(m-3)=3-m$，$\alpha\beta=2$，}
\step{又 $\alpha^2+\beta^2=\alpha\beta+3$，则 $(\alpha+\beta)^2-2\alpha\beta=\alpha\beta+3$，}
\step{即 $(\alpha+\beta)^2=3\alpha\beta+3$，所以 $(3-m)^2=3\times2+3=9$，解得 $m=0$ 或 $6$.}
\solend
\fi

\item 设集合 $A=\{x\mid\lvert x-a\rvert<2\}$，$B=\left\{x\mid\dfrac{x+4}{x+2}>2\right\}$.

（1）若 $a=1$，试用区间表示集合 $A$、$B$；

（2）若 $A\cup\overline{B}=\R$，求实数 $a$ 的取值范围.

\ifanswer
\solbegin
\step{（1）当 $a=1$ 时，由 $\lvert x-1\rvert<2$ 得 $-2<x-1<2$，解得 $-1<x<3$，}
\step{所以 $A=(-1,3)$.}
\step{由 $\dfrac{x+4}{x+2}>2$ 得 $2-\dfrac{x+4}{x+2}=\dfrac{x}{x+2}<0$，则有 $x(x+2)<0$，解得 $-2<x<0$，}
\step{所以 $B=(-2,0)$；}
\step{（2）由 $\lvert x-a\rvert<2$ 得 $-2<x-a<2$，解得 $a-2<x<a+2$，}
\step{所以 $A=(a-2,a+2)$，}
\step{由（1）得 $B=(-2,0)$，则 $\overline{B}=\{x\mid x\geqslant0$ 或 $x\leqslant-2\}$，}
\step{由于 $A\cup\overline{B}=\R$，所以 $\begin{cases}a+2\geqslant0\\a-2\leqslant-2\end{cases}$，解得 $-2\leqslant a\leqslant0$，}
\step{所以实数 $a$ 的取值范围是 $[-2,0]$.}
\solend
\fi

\item 某服装公司生产的衬衫每件定价为 80 元，在某城市年销售 8 万件. 现该公司计划在该市招收代理商来销售衬衫，以降低管理和营销成本. 已知代理商收取的代理费为总销售金额的 $r\%$（其中 $0<r\leqslant10$）.

（1）若单件衬衫的利润保持不变，求招代理商后每件衬衫的定价为多少元（用 $r$ 表示）；

（2）已知将每件衬衫的价格提高到 $\dfrac{80}{1-r\%}$ 元后，每年的销售量会减少 $0.62r$ 万件. 若代理商每年收取的代理费不少于 16 万元，求 $r$ 的取值范围.

\ifanswer
\solbegin
\step{（1）设招代理商后每件衬衫的定价为 $x$ 元，代理商收取总销售金额的 $r\%$，}
\step{因此公司每件衬衫的实际净收入为 $x(1-r\%)$ 元，}
\step{单件利润保持不变，等价于净收入仍为 80 元，}
\step{于是 $x(1-r\%)=80$，解得 $x=\dfrac{80}{1-r\%}$，因为 $0<r\leqslant10$，有 $1-r\%\geqslant0.9>0$，}
\step{分母恒不为零，该定价有意义，所以招代理商后每件衬衫的定价为 $\dfrac{80}{1-r\%}$ 元；}
\step{（2）价格提高到 $\dfrac{80}{1-r\%}$ 元，年销售量减少 $0.62r$ 万件，}
\step{故年销售量为 $8-0.62r$ 万件，$8-0.62r>0$，}
\step{由题意得 $\dfrac{80}{1-r\%}\cdot(8-0.62r)\cdot r\%\geqslant16$，}
\step{即 $\dfrac{80}{1-\dfrac r{100}}\cdot(8-0.62r)\cdot\dfrac r{100}=\dfrac{80r(8-0.62r)}{100-r}\geqslant16$，}
\step{由 $0<r\leqslant10$ 得 $100-r>0$，整理得 $80r(8-0.62r)\geqslant16(100-r)$，}
\step{整理得 $31r^2-410r+1000\leqslant0$，解得 $\dfrac{100}{31}\leqslant r\leqslant10$，}
\step{当 $\dfrac{100}{31}\leqslant r\leqslant10$ 时，$8-0.62r>0$ 且 $1-r\%>0$，}
\step{均符合实际经营条件，所以 $r$ 的取值范围是 $\left[\dfrac{100}{31},10\right]$.}
\solend
\fi

\item 定义 $\min\{a_1,a_2,\cdots,a_n\}$ 为 $n$ 个实数 $a_1,a_2,\cdots,a_n$ 中的最小数，$\max\{a_1,a_2,\cdots,a_n\}$ 为 $n$ 个实数 $a_1,a_2,\cdots,a_n$ 中的最大数.

（1）设 $a,b$ 都是正实数，且 $a+b=1$，求 $\max\left\{ab,\dfrac14\right\}$；

（2）解不等式：$\min\{x+1,x^2+3,\lvert x-1\rvert\}>2x-3$；

（3）设 $a,b$ 都是正实数，求 $\max\left\{a+\dfrac1b,\dfrac2a+b\right\}$ 的最小值.

\ifanswer
\solbegin
\step{（1）$ab\leqslant\dfrac{(a+b)^2}{4}=\dfrac14$，所以 $\max\left\{ab,\dfrac14\right\}=\dfrac14$.}
\step{（2）因为 $x^2+3>x+1$，$x^2+3>\lvert x-1\rvert$，}
\step{所以 $\min\{x+1,x^2+3,\lvert x-1\rvert\}=\min\{x+1,\lvert x-1\rvert\}$.}
\step{当 $x\geqslant0$ 时，$\lvert x-1\rvert\leqslant x+1$，所以 $\lvert x-1\rvert>2x-3$，解得 $0\leqslant x<2$；}
\step{当 $x<0$ 时，$\lvert x-1\rvert>x+1$，所以 $x+1>2x-3$，解得 $x<0$.}
\step{综上，不等式的解集为 $(-\infty,2)$.}
\step{（3）设 $M=\max\left\{a+\dfrac1b,\dfrac2a+b\right\}$，则 $M\geqslant a+\dfrac1b$，$M\geqslant\dfrac2a+b$，}
\step{所以 $M^2\geqslant\left(a+\dfrac1b\right)\left(\dfrac2a+b\right)=3+ab+\dfrac2{ab}\geqslant3+2\sqrt2$，}
\step{所以 $M$ 的最小值为 $\sqrt2+1$，当 $a=\sqrt2$，$b=1$ 时取等号.}
\solend
\fi

\end{enumerate}

\end{document}
"
% 紧凑版：xelatex "\def\iscompact{1}% 高一126_建青实验_期中.tex
%   —— 高一上数学"2026-2027 年建青实验高一上期中"（试卷版/答案版）
% 试卷版：xelatex 高一126_建青实验_期中.tex
% 答案版：xelatex "\def\isanswer{1}\input{高一126_建青实验_期中.tex}"
% 紧凑版：xelatex "\def\iscompact{1}\input{高一126_建青实验_期中.tex}"
%
% 原始材料：微信公众号"魔都数学加油站"文章，作者破晓，2026-11-11 18:35 发布；
% 连载"27学年高一126"，原文标题"27学年高一126——2026-2027年上海市建青实验高一上期中"；
% 原文链接：https://mp.weixin.qq.com/s/HpxpBJu4d_wOnUCIbZSQiA。
% 正文为 8 张 PNG 页（01、05、08 为 2560x3620，02、03、04、06、07 为 4762x6735），
% 题目下印有【解析】或【答案】，为讲评版。题目、解析均据原图转录；
% 卷面粉红色水印"魔都数学加油站"及公众号页脚未录入。全卷无插图。
%
% 卷首标题：原卷印作"2026-2027 年建青实验高一上期中"（公众号标题中的"上海市"
% 卷面标题行没有，本稿照卷面），年份区间按全稿惯例排 en dash；原卷未印副标题，
% 本卷为期中卷、跨集合与常用逻辑用语／不等式／函数三章，按题目内容标为
% "集合与常用逻辑用语、不等式、函数"（与 125 同为期中卷，口径相同）：
% 集合与逻辑见 Q5、Q8、Q9、Q11、Q13、Q14、Q16、Q19，不等式见 Q6、Q7、Q10、
% Q17、Q20、Q21，函数（含指数对数、二次函数零点）见 Q3、Q4、Q12、Q15、Q18。
%
% 与原卷的排版差异：
%   ① 原文从第 3 题起，填空题编号为 3—12、选择题为 13—16、解答题为 17—21，本稿保留原题号；
%   ② 原卷每题下直接印【解析】／【答案】，本稿按全稿惯例置于答案版；第 4、5 题只印
%      【答案】，另提 \ans{}；第 13—16 题答案均在【解析】末句"故选 …"，题干括号内留空、
%      不另提 \ans{}；第 3 题印一行【解析】，按 \solbegin+\step 排；第 17 题解析保留
%      原卷评分点（照 高一54／64 成例排 \dots\dots N 分）；
%   ③ 原卷未印副标题，本稿按"知识模块"口径补出（见上"卷首标题"）；
%   ④ 原卷填空横线按全稿惯例排为 \blankline[4.4em]；选择题括号排 \mbox{（\quad）}、
%      选项段前按惯例加 \keepwithprev；第 14 题原卷括号另起一行，本稿并入题干行末；
%      解答题不留作答区（原卷为讲评版，与 109—125 同）；
%   ⑤ 原卷数集排斜体 $R$，按全稿惯例改 \R；绝对值排 \lvert…\rvert，集合描述竖线排 \mid；
%      补集记号照原卷上划线排 \overline{B}；不等号按全稿惯例排 \leqslant／\geqslant；
%      子题标号原卷一律用全角括号（1）（2）（3），本稿照录（\step 行内亦同，照 高一121 成例）。
%
% 判别式字形：第 10 题解析 $\Delta=a^2-4a<0$ 原卷排 $\Delta$（等腰、高），
%   照录 \Delta（与 高一124／125 同题排 $\triangle$（等边、扁）不同，逐处照本卷字形）。
%
% 子集符号：第 8 题题干与解析两处 $S\subset P$ 均印单线、无下划线，照录 \subset；
%   全卷未出现带下划线的 \subseteq。
%
% 原卷存疑照录：
%   第 10 题题干作 $ax^2-ax+1>0$，解析却作 $ax^2+ax+1>0$（一次项符号不一致），照录、未统一；
%   第 11 题解析"此时 $M=\dfrac32,1,\dfrac12$，满足"未用集合花括号（题干 $M$ 为集合），照录；
%   第 11 题解析"因为它们是集合 $M$ 中的元素，所以存在互异性"（"存在互异性"疑为
%     "满足互异性"之误），又"当 $x_1\neq x_2\neq x_3$ 时"连写两个不等号，均照录；
%   第 16 题解析两处"惟如此才能保证"（用"惟"字）及"符合定义的参数 $m$ 的值一定大于等于 1
%     或小于等于 0"的措辞，照录；
%   第 17 题 (1) 解析"当且仅当 $a=1,b=-1$ 时"用半角逗号、第 21 题 (3) 解析
%     "当 $a=\sqrt2$，$b=1$ 时取等号"用全角逗号，均照录；
%   第 19 题解析由 $\dfrac{x+4}{x+2}>2$ 得 "$2-\dfrac{x+4}{x+2}=\dfrac{x}{x+2}<0$"
%     （以 2 减左式，与常见的"左式减 2"次序相反），照录；
%   第 15 题选项 A、B 作"该方程…"、C、D 作"方程组…"，照录；
%   第 20 题题干作"招收代理商"，(1) 却作"招代理商后"，照录。
%
% 行数复核：逐题数 \step 条数，合计 91 条（Q3：1；Q4：0；Q5：0；Q6：3；Q7：2；
% Q8：5；Q9：5；Q10：4；Q11：8；Q12：3；Q13：3；Q14：2；Q15：3；Q16：9；
% Q17：7；Q18：5；Q19：9；Q20：13；Q21：9），与原卷印刷行逐题相等
% （cases 块作一个印刷行计）。第 4、5 题只印【答案】，提 \ans{}，不计入 \step 账；
% 四道选择题答案均在解析末句，括号内留空、不另提 \ans{}。
%
\documentclass[a4paper,11pt]{article}
\usepackage{common}

\begin{document}

\examtitlest[3]{2026--2027 年建青实验高一上期中}{集合与常用逻辑用语、不等式、函数}

\bigsec{一、填空题}

\begin{enumerate}
\setcounter{enumi}{2}

\item 已知 $\log_2 x=3$，则 $x=$\blankline[4.4em].

\ifanswer
\solbegin
\step{$x=2^3=8$.}
\solend
\fi

\item 将 $\sqrt{a^2\sqrt a}$（其中 $a>0$）化为有理数指数幂的形式为\blankline[4.4em].

\ifanswer
\ans{$a^{\frac54}$}
\fi

\item 用反证法证明命题：“已知 $m$、$n\in\R$，若 $m+n\leqslant5$，则 $m\leqslant2$ 或 $n\leqslant3$” 时应假设\blankline[4.4em].

\ifanswer
\ans{$m>2$ 且 $n>3$}
\fi

\item 若 $x\in\R$，则 $\lvert x+1\rvert+\lvert x-1\rvert$ 的最小值为\blankline[4.4em].

\ifanswer
\solbegin
\step{$\lvert x+1\rvert+\lvert x-1\rvert\geqslant\lvert(x+1)-(x-1)\rvert=2$，}
\step{当且仅当 $(x+1)(x-1)\leqslant0$，即 $-1\leqslant x\leqslant1$ 时取等号，}
\step{所以 $\lvert x+1\rvert+\lvert x-1\rvert$ 的最小值为 2.}
\solend
\fi

\item 已知 $x>1$，则 $x+\dfrac1{x-1}$ 的最小值是\blankline[4.4em].

\ifanswer
\solbegin
\step{$x-1>0$，则 $x+\dfrac1{x-1}=(x-1)+\dfrac1{x-1}+1\geqslant2\sqrt{(x-1)\cdot\dfrac1{x-1}}+1=3$，}
\step{当且仅当 $x=2$ 时取等号，故 $x+\dfrac1{x-1}$ 的最小值为 3.}
\solend
\fi

\item 已知集合 $P=\{x\mid x^2+x-6=0\}$，$S=\{x\mid ax+1=0\}$ 且 $S\subset P$，则 $a$ 的值为\blankline[4.4em].

\ifanswer
\solbegin
\step{集合 $P=\{-3,2\}$，集合 $S$ 中至多有一个元素，}
\step{若集合 $S$ 为空集，即 $a=0$ 时，显然满足条件 $S\subset P$，故 $a=0$.}
\step{若集合 $S$ 非空集，即 $a\neq0$，此时 $S=\left\{-\dfrac1a\right\}$，}
\step{若 $-\dfrac1a=-3$，则 $a=\dfrac13$，若 $-\dfrac1a=2$，则 $a=-\dfrac12$；}
\step{故 $a$ 的取值为 $0$、$\dfrac13$、$-\dfrac12$.}
\solend
\fi

\item 已知 $p:\begin{cases}a>4\\b>4\end{cases}$，$q:\begin{cases}a+b>8\\ab>16\end{cases}$，则 $p$ 是 $q$ 的\blankline[4.4em]条件.

\ifanswer
\solbegin
\step{若 $a>4$，$b>4$，则 $a+b>4+4=8$，$ab>4\times4=16$，}
\step{所以 $p$ 成立一定能推出 $q$ 成立，充分性成立，}
\step{取 $a=2$，$b=9$，$a+b=11>8$，$ab=18>16$，满足 $q$；}
\step{但是 $a=2<4$，不满足 $p$，所以由 $q$ 不能推出 $p$，必要性不成立，}
\step{综上所述，$p$ 是 $q$ 的充分不必要条件.}
\solend
\fi

\item 若关于 $x$ 的不等式 $ax^2-ax+1>0$ 对一切实数 $x$ 恒成立，则实数 $a$ 的取值范围是\blankline[4.4em].

\ifanswer
\solbegin
\step{当 $a=0$ 时，不等式即 $1>0$，满足条件.}
\step{当 $a\neq0$ 时，要使不等式 $ax^2+ax+1>0$ 对一切 $x\in\R$ 恒成立，}
\step{需 $\begin{cases}a>0\\\Delta=a^2-4a<0\end{cases}$，解得 $0<a<4$.}
\step{综上，实数 $a$ 的取值范围是 $[0,4)$.}
\solend
\fi

\item 已知集合 $M=\{x\mid (x-a)(x^2-ax+a-1)=0\}$ 各元素之和等于 3，则实数 $a=$\blankline[4.4em].

\ifanswer
\solbegin
\step{方程 $(x-a)(x^2-ax+a-1)=0$，化简得 $(x-a)(x-1)[x-(a-1)]=0$，}
\step{解得 $x_1=a$，$x_2=1$，$x_3=a-1$，}
\step{因为它们是集合 $M$ 中的元素，所以存在互异性.}
\step{当 $x_1\neq x_2\neq x_3$ 时，即 $a+1+a-1=3$，解得 $a=\dfrac32$，此时 $M=\dfrac32,1,\dfrac12$，满足.}
\step{当 $a=1$ 时，$1+0\neq3$，不合题意；}
\step{当 $a=a-1$ 时，不符合题意；}
\step{当 $1=a-1$ 时，$a=2$，此时 $M=\{2,1\}$ 满足题意.}
\step{故 $a$ 的取值为 $2$ 或 $\dfrac32$.}
\solend
\fi

\item 已知 $a>b>1$ 且 $a^2-6a-2m+1=0$，$b^2-6b-2m+1=0$，则实数 $m$ 的取值范围是\blankline[4.4em].

\ifanswer
\solbegin
\step{由题意得 $a$、$b$ 是函数 $f(x)=x^2-6x-2m+1$ 两个大于 1 的不同零点，}
\step{所以 $\begin{cases}f(1)>0\\f(3)<0\end{cases}$，即 $\begin{cases}-4-2m>0\\-8-2m<0\end{cases}$，解得 $-4<m<-2$，}
\step{故 $m$ 的取值范围为 $(-4,-2)$.}
\solend
\fi

\end{enumerate}

\bigsec{二、选择题}

\begin{enumerate}
\setcounter{enumi}{12}

\item “$x<2$” 是 “$x^2<4$” 的\mbox{（\quad）}条件

\keepwithprev
A. 充分不必要\qquad\qquad B. 必要不充分

C. 充要\qquad\qquad\qquad\quad D. 既不充分也不必要

\ifanswer
\solbegin
\step{解 $x^2<4$，得 $-2<x<2$，}
\step{由 $x<2$ 得不出 $-2<x<2$，而由 $-2<x<2$ 能得出 $x<2$，}
\step{所以 “$x<2$” 是 “$x^2<4$” 的必要不充分条件. 故选 $B$.}
\solend
\fi

\item 已知集合 $A=\{x\mid x<-2\}$，$B=\{x\mid x\geqslant a\}$，且 $A\cup B=\R$，则实数 $a$ 的取值范围是\mbox{（\quad）}

\keepwithprev
A. $a\leqslant-2$\qquad\qquad B. $a<-2$\qquad\qquad C. $a\geqslant-2$\qquad\qquad D. $a>-2$

\ifanswer
\solbegin
\step{集合 $A=\{x\mid x<-2\}$，$B=\{x\mid x\geqslant a\}$，且 $A\cup B=\R$，}
\step{则实数 $a$ 的取值范围是 $a\leqslant-2$. 故选 $A$.}
\solend
\fi

\item 下列关于 $x,y$ 的方程组 $\begin{cases}y=kx+1\\y=2kx+3\end{cases}$ 的解的说法中正确的是\mbox{（\quad）}

\keepwithprev
A. 该方程一定有唯一解\qquad\qquad B. 该方程没有解

C. $k\neq0$ 时，方程组有唯一解\qquad\qquad D. $k=0$ 时，方程组有无数组解

\ifanswer
\solbegin
\step{由题意得 $kx+1=2kx+3$，即 $kx=-2$，}
\step{当 $k=0$ 时，$0=-2$ 不成立，方程组无解；}
\step{当 $k\neq0$ 时，$x=-\dfrac2k$，方程组有唯一解. 故选 $C$.}
\solend
\fi

\item 设非空集合 $S=\{x\mid m\leqslant x\leqslant n\}$ 满足：当 $x\in S$ 时，有 $x^2\in S$. 给出如下三个结论：①若 $m=1$，则 $S=\{1\}$；②若 $m=-\dfrac12$，则 $\dfrac14\leqslant n\leqslant1$；③若 $n=\dfrac12$，则 $-\dfrac{\sqrt2}{2}\leqslant m\leqslant0$. 其中正确命题的个数为\mbox{（\quad）}

\keepwithprev
A. $3$\qquad\qquad B. $2$\qquad\qquad C. $1$\qquad\qquad D. $0$

\ifanswer
\solbegin
\step{由定义设非空集合 $S=\{x\mid m\leqslant x\leqslant n\}$ 满足当 $x\in S$ 时，有 $x^2\in S$，}
\step{符合定义的参数 $m$ 的值一定大于等于 1 或小于等于 0，}
\step{惟如此才能保证 $m\in S$ 时，有 $m^2\in S$ 即 $m^2\geqslant m$，}
\step{符合条件的 $n$ 的值一定大于等于 0，小于等于 1，}
\step{惟如此才能保证 $n\in S$ 时，有 $n^2\in S$ 即 $n^2\leqslant n$，对各个命题进行判断：}
\step{对于① $m=1$，$m^2=1\in S$，故必有 $\begin{cases}n^2\leqslant n\\n\geqslant1\end{cases}$，得 $n=1$，$S=\{1\}$，}
\step{② $m=-\dfrac12$，$m^2=\dfrac14\in S$，则 $\begin{cases}n^2\leqslant n\\\dfrac14\leqslant n\end{cases}$，解得 $\dfrac14\leqslant n\leqslant1$；}
\step{对于③，若 $n=\dfrac12$，则 $\begin{cases}\dfrac12\geqslant m\\m^2\geqslant m\\\dfrac12\geqslant m^2\end{cases}$，解得 $-\dfrac{\sqrt2}{2}\leqslant m\leqslant0$，}
\step{所以正确命题有 3 个. 故选 $A$.}
\solend
\fi

\end{enumerate}

\bigsec{三、解答题}

\begin{enumerate}
\setcounter{enumi}{16}

\item 已知 $a$、$b$ 是实数.

（1）求证：$a^2+b^2\geqslant2a-2b-2$，并指出等号成立的条件；

（2）若 $ab=1$，求 $a^2+4b^2$ 的最小值.

\ifanswer
\solbegin
\step{（1）因为 $a^2+b^2-(2a-2b-2)=a^2+b^2-2a+2b+2$\dots\dots2 分}
\step{$=(a-1)^2+(b+1)^2\geqslant0$\dots\dots4 分}
\step{所以 $a^2+b^2\geqslant2a-2b-2$.}
\step{当且仅当 $a=1,b=-1$ 时，不等式中等号成立.\dots\dots7 分}
\step{（2）$a^2+4b^2=a^2+(2b)^2\geqslant2\cdot a\cdot(2b)=4ab=4$，\dots\dots5 分}
\step{当且仅当 $a=2b$，即 $\begin{cases}a=\sqrt2\\b=\dfrac{\sqrt2}{2}\end{cases}$ 或 $\begin{cases}a=-\sqrt2\\b=-\dfrac{\sqrt2}{2}\end{cases}$ 时，不等式中等号成立，}
\step{所以 $a^2+4b^2$ 的最小值为 4.\dots\dots7 分}
\solend
\fi

\item 已知关于 $x$ 的方程 $x^2+(m-3)x+2=0$，其中 $m\in\R$，有两个实数根 $\alpha$、$\beta$.

（1）若 $m=8$，求 $\dfrac1\alpha+\dfrac1\beta$ 的值；

（2）若 $\alpha^2+\beta^2=\alpha\beta+3$，求实数 $m$ 的值.

\ifanswer
\solbegin
\step{（1）若 $m=8$，则关于 $x$ 的方程为 $x^2+5x+2=0$，}
\step{则 $\alpha+\beta=-5$，$\alpha\beta=2$，所以 $\dfrac1\alpha+\dfrac1\beta=\dfrac{\alpha+\beta}{\alpha\beta}=-\dfrac52$；}
\step{（2）由题意得 $\alpha+\beta=-(m-3)=3-m$，$\alpha\beta=2$，}
\step{又 $\alpha^2+\beta^2=\alpha\beta+3$，则 $(\alpha+\beta)^2-2\alpha\beta=\alpha\beta+3$，}
\step{即 $(\alpha+\beta)^2=3\alpha\beta+3$，所以 $(3-m)^2=3\times2+3=9$，解得 $m=0$ 或 $6$.}
\solend
\fi

\item 设集合 $A=\{x\mid\lvert x-a\rvert<2\}$，$B=\left\{x\mid\dfrac{x+4}{x+2}>2\right\}$.

（1）若 $a=1$，试用区间表示集合 $A$、$B$；

（2）若 $A\cup\overline{B}=\R$，求实数 $a$ 的取值范围.

\ifanswer
\solbegin
\step{（1）当 $a=1$ 时，由 $\lvert x-1\rvert<2$ 得 $-2<x-1<2$，解得 $-1<x<3$，}
\step{所以 $A=(-1,3)$.}
\step{由 $\dfrac{x+4}{x+2}>2$ 得 $2-\dfrac{x+4}{x+2}=\dfrac{x}{x+2}<0$，则有 $x(x+2)<0$，解得 $-2<x<0$，}
\step{所以 $B=(-2,0)$；}
\step{（2）由 $\lvert x-a\rvert<2$ 得 $-2<x-a<2$，解得 $a-2<x<a+2$，}
\step{所以 $A=(a-2,a+2)$，}
\step{由（1）得 $B=(-2,0)$，则 $\overline{B}=\{x\mid x\geqslant0$ 或 $x\leqslant-2\}$，}
\step{由于 $A\cup\overline{B}=\R$，所以 $\begin{cases}a+2\geqslant0\\a-2\leqslant-2\end{cases}$，解得 $-2\leqslant a\leqslant0$，}
\step{所以实数 $a$ 的取值范围是 $[-2,0]$.}
\solend
\fi

\item 某服装公司生产的衬衫每件定价为 80 元，在某城市年销售 8 万件. 现该公司计划在该市招收代理商来销售衬衫，以降低管理和营销成本. 已知代理商收取的代理费为总销售金额的 $r\%$（其中 $0<r\leqslant10$）.

（1）若单件衬衫的利润保持不变，求招代理商后每件衬衫的定价为多少元（用 $r$ 表示）；

（2）已知将每件衬衫的价格提高到 $\dfrac{80}{1-r\%}$ 元后，每年的销售量会减少 $0.62r$ 万件. 若代理商每年收取的代理费不少于 16 万元，求 $r$ 的取值范围.

\ifanswer
\solbegin
\step{（1）设招代理商后每件衬衫的定价为 $x$ 元，代理商收取总销售金额的 $r\%$，}
\step{因此公司每件衬衫的实际净收入为 $x(1-r\%)$ 元，}
\step{单件利润保持不变，等价于净收入仍为 80 元，}
\step{于是 $x(1-r\%)=80$，解得 $x=\dfrac{80}{1-r\%}$，因为 $0<r\leqslant10$，有 $1-r\%\geqslant0.9>0$，}
\step{分母恒不为零，该定价有意义，所以招代理商后每件衬衫的定价为 $\dfrac{80}{1-r\%}$ 元；}
\step{（2）价格提高到 $\dfrac{80}{1-r\%}$ 元，年销售量减少 $0.62r$ 万件，}
\step{故年销售量为 $8-0.62r$ 万件，$8-0.62r>0$，}
\step{由题意得 $\dfrac{80}{1-r\%}\cdot(8-0.62r)\cdot r\%\geqslant16$，}
\step{即 $\dfrac{80}{1-\dfrac r{100}}\cdot(8-0.62r)\cdot\dfrac r{100}=\dfrac{80r(8-0.62r)}{100-r}\geqslant16$，}
\step{由 $0<r\leqslant10$ 得 $100-r>0$，整理得 $80r(8-0.62r)\geqslant16(100-r)$，}
\step{整理得 $31r^2-410r+1000\leqslant0$，解得 $\dfrac{100}{31}\leqslant r\leqslant10$，}
\step{当 $\dfrac{100}{31}\leqslant r\leqslant10$ 时，$8-0.62r>0$ 且 $1-r\%>0$，}
\step{均符合实际经营条件，所以 $r$ 的取值范围是 $\left[\dfrac{100}{31},10\right]$.}
\solend
\fi

\item 定义 $\min\{a_1,a_2,\cdots,a_n\}$ 为 $n$ 个实数 $a_1,a_2,\cdots,a_n$ 中的最小数，$\max\{a_1,a_2,\cdots,a_n\}$ 为 $n$ 个实数 $a_1,a_2,\cdots,a_n$ 中的最大数.

（1）设 $a,b$ 都是正实数，且 $a+b=1$，求 $\max\left\{ab,\dfrac14\right\}$；

（2）解不等式：$\min\{x+1,x^2+3,\lvert x-1\rvert\}>2x-3$；

（3）设 $a,b$ 都是正实数，求 $\max\left\{a+\dfrac1b,\dfrac2a+b\right\}$ 的最小值.

\ifanswer
\solbegin
\step{（1）$ab\leqslant\dfrac{(a+b)^2}{4}=\dfrac14$，所以 $\max\left\{ab,\dfrac14\right\}=\dfrac14$.}
\step{（2）因为 $x^2+3>x+1$，$x^2+3>\lvert x-1\rvert$，}
\step{所以 $\min\{x+1,x^2+3,\lvert x-1\rvert\}=\min\{x+1,\lvert x-1\rvert\}$.}
\step{当 $x\geqslant0$ 时，$\lvert x-1\rvert\leqslant x+1$，所以 $\lvert x-1\rvert>2x-3$，解得 $0\leqslant x<2$；}
\step{当 $x<0$ 时，$\lvert x-1\rvert>x+1$，所以 $x+1>2x-3$，解得 $x<0$.}
\step{综上，不等式的解集为 $(-\infty,2)$.}
\step{（3）设 $M=\max\left\{a+\dfrac1b,\dfrac2a+b\right\}$，则 $M\geqslant a+\dfrac1b$，$M\geqslant\dfrac2a+b$，}
\step{所以 $M^2\geqslant\left(a+\dfrac1b\right)\left(\dfrac2a+b\right)=3+ab+\dfrac2{ab}\geqslant3+2\sqrt2$，}
\step{所以 $M$ 的最小值为 $\sqrt2+1$，当 $a=\sqrt2$，$b=1$ 时取等号.}
\solend
\fi

\end{enumerate}

\end{document}
"
%
% 原始材料：微信公众号"魔都数学加油站"文章，作者破晓，2026-11-11 18:35 发布；
% 连载"27学年高一126"，原文标题"27学年高一126——2026-2027年上海市建青实验高一上期中"；
% 原文链接：https://mp.weixin.qq.com/s/HpxpBJu4d_wOnUCIbZSQiA。
% 正文为 8 张 PNG 页（01、05、08 为 2560x3620，02、03、04、06、07 为 4762x6735），
% 题目下印有【解析】或【答案】，为讲评版。题目、解析均据原图转录；
% 卷面粉红色水印"魔都数学加油站"及公众号页脚未录入。全卷无插图。
%
% 卷首标题：原卷印作"2026-2027 年建青实验高一上期中"（公众号标题中的"上海市"
% 卷面标题行没有，本稿照卷面），年份区间按全稿惯例排 en dash；原卷未印副标题，
% 本卷为期中卷、跨集合与常用逻辑用语／不等式／函数三章，按题目内容标为
% "集合与常用逻辑用语、不等式、函数"（与 125 同为期中卷，口径相同）：
% 集合与逻辑见 Q5、Q8、Q9、Q11、Q13、Q14、Q16、Q19，不等式见 Q6、Q7、Q10、
% Q17、Q20、Q21，函数（含指数对数、二次函数零点）见 Q3、Q4、Q12、Q15、Q18。
%
% 与原卷的排版差异：
%   ① 原文从第 3 题起，填空题编号为 3—12、选择题为 13—16、解答题为 17—21，本稿保留原题号；
%   ② 原卷每题下直接印【解析】／【答案】，本稿按全稿惯例置于答案版；第 4、5 题只印
%      【答案】，另提 \ans{}；第 13—16 题答案均在【解析】末句"故选 …"，题干括号内留空、
%      不另提 \ans{}；第 3 题印一行【解析】，按 \solbegin+\step 排；第 17 题解析保留
%      原卷评分点（照 高一54／64 成例排 \dots\dots N 分）；
%   ③ 原卷未印副标题，本稿按"知识模块"口径补出（见上"卷首标题"）；
%   ④ 原卷填空横线按全稿惯例排为 \blankline[4.4em]；选择题括号排 \mbox{（\quad）}、
%      选项段前按惯例加 \keepwithprev；第 14 题原卷括号另起一行，本稿并入题干行末；
%      解答题不留作答区（原卷为讲评版，与 109—125 同）；
%   ⑤ 原卷数集排斜体 $R$，按全稿惯例改 \R；绝对值排 \lvert…\rvert，集合描述竖线排 \mid；
%      补集记号照原卷上划线排 \overline{B}；不等号按全稿惯例排 \leqslant／\geqslant；
%      子题标号原卷一律用全角括号（1）（2）（3），本稿照录（\step 行内亦同，照 高一121 成例）。
%
% 判别式字形：第 10 题解析 $\Delta=a^2-4a<0$ 原卷排 $\Delta$（等腰、高），
%   照录 \Delta（与 高一124／125 同题排 $\triangle$（等边、扁）不同，逐处照本卷字形）。
%
% 子集符号：第 8 题题干与解析两处 $S\subset P$ 均印单线、无下划线，照录 \subset；
%   全卷未出现带下划线的 \subseteq。
%
% 原卷存疑照录：
%   第 10 题题干作 $ax^2-ax+1>0$，解析却作 $ax^2+ax+1>0$（一次项符号不一致），照录、未统一；
%   第 11 题解析"此时 $M=\dfrac32,1,\dfrac12$，满足"未用集合花括号（题干 $M$ 为集合），照录；
%   第 11 题解析"因为它们是集合 $M$ 中的元素，所以存在互异性"（"存在互异性"疑为
%     "满足互异性"之误），又"当 $x_1\neq x_2\neq x_3$ 时"连写两个不等号，均照录；
%   第 16 题解析两处"惟如此才能保证"（用"惟"字）及"符合定义的参数 $m$ 的值一定大于等于 1
%     或小于等于 0"的措辞，照录；
%   第 17 题 (1) 解析"当且仅当 $a=1,b=-1$ 时"用半角逗号、第 21 题 (3) 解析
%     "当 $a=\sqrt2$，$b=1$ 时取等号"用全角逗号，均照录；
%   第 19 题解析由 $\dfrac{x+4}{x+2}>2$ 得 "$2-\dfrac{x+4}{x+2}=\dfrac{x}{x+2}<0$"
%     （以 2 减左式，与常见的"左式减 2"次序相反），照录；
%   第 15 题选项 A、B 作"该方程…"、C、D 作"方程组…"，照录；
%   第 20 题题干作"招收代理商"，(1) 却作"招代理商后"，照录。
%
% 行数复核：逐题数 \step 条数，合计 91 条（Q3：1；Q4：0；Q5：0；Q6：3；Q7：2；
% Q8：5；Q9：5；Q10：4；Q11：8；Q12：3；Q13：3；Q14：2；Q15：3；Q16：9；
% Q17：7；Q18：5；Q19：9；Q20：13；Q21：9），与原卷印刷行逐题相等
% （cases 块作一个印刷行计）。第 4、5 题只印【答案】，提 \ans{}，不计入 \step 账；
% 四道选择题答案均在解析末句，括号内留空、不另提 \ans{}。
%
\documentclass[a4paper,11pt]{article}
\usepackage{common}

\begin{document}

\examtitlest[3]{2026--2027 年建青实验高一上期中}{集合与常用逻辑用语、不等式、函数}

\bigsec{一、填空题}

\begin{enumerate}
\setcounter{enumi}{2}

\item 已知 $\log_2 x=3$，则 $x=$\blankline[4.4em].

\ifanswer
\solbegin
\step{$x=2^3=8$.}
\solend
\fi

\item 将 $\sqrt{a^2\sqrt a}$（其中 $a>0$）化为有理数指数幂的形式为\blankline[4.4em].

\ifanswer
\ans{$a^{\frac54}$}
\fi

\item 用反证法证明命题：“已知 $m$、$n\in\R$，若 $m+n\leqslant5$，则 $m\leqslant2$ 或 $n\leqslant3$” 时应假设\blankline[4.4em].

\ifanswer
\ans{$m>2$ 且 $n>3$}
\fi

\item 若 $x\in\R$，则 $\lvert x+1\rvert+\lvert x-1\rvert$ 的最小值为\blankline[4.4em].

\ifanswer
\solbegin
\step{$\lvert x+1\rvert+\lvert x-1\rvert\geqslant\lvert(x+1)-(x-1)\rvert=2$，}
\step{当且仅当 $(x+1)(x-1)\leqslant0$，即 $-1\leqslant x\leqslant1$ 时取等号，}
\step{所以 $\lvert x+1\rvert+\lvert x-1\rvert$ 的最小值为 2.}
\solend
\fi

\item 已知 $x>1$，则 $x+\dfrac1{x-1}$ 的最小值是\blankline[4.4em].

\ifanswer
\solbegin
\step{$x-1>0$，则 $x+\dfrac1{x-1}=(x-1)+\dfrac1{x-1}+1\geqslant2\sqrt{(x-1)\cdot\dfrac1{x-1}}+1=3$，}
\step{当且仅当 $x=2$ 时取等号，故 $x+\dfrac1{x-1}$ 的最小值为 3.}
\solend
\fi

\item 已知集合 $P=\{x\mid x^2+x-6=0\}$，$S=\{x\mid ax+1=0\}$ 且 $S\subset P$，则 $a$ 的值为\blankline[4.4em].

\ifanswer
\solbegin
\step{集合 $P=\{-3,2\}$，集合 $S$ 中至多有一个元素，}
\step{若集合 $S$ 为空集，即 $a=0$ 时，显然满足条件 $S\subset P$，故 $a=0$.}
\step{若集合 $S$ 非空集，即 $a\neq0$，此时 $S=\left\{-\dfrac1a\right\}$，}
\step{若 $-\dfrac1a=-3$，则 $a=\dfrac13$，若 $-\dfrac1a=2$，则 $a=-\dfrac12$；}
\step{故 $a$ 的取值为 $0$、$\dfrac13$、$-\dfrac12$.}
\solend
\fi

\item 已知 $p:\begin{cases}a>4\\b>4\end{cases}$，$q:\begin{cases}a+b>8\\ab>16\end{cases}$，则 $p$ 是 $q$ 的\blankline[4.4em]条件.

\ifanswer
\solbegin
\step{若 $a>4$，$b>4$，则 $a+b>4+4=8$，$ab>4\times4=16$，}
\step{所以 $p$ 成立一定能推出 $q$ 成立，充分性成立，}
\step{取 $a=2$，$b=9$，$a+b=11>8$，$ab=18>16$，满足 $q$；}
\step{但是 $a=2<4$，不满足 $p$，所以由 $q$ 不能推出 $p$，必要性不成立，}
\step{综上所述，$p$ 是 $q$ 的充分不必要条件.}
\solend
\fi

\item 若关于 $x$ 的不等式 $ax^2-ax+1>0$ 对一切实数 $x$ 恒成立，则实数 $a$ 的取值范围是\blankline[4.4em].

\ifanswer
\solbegin
\step{当 $a=0$ 时，不等式即 $1>0$，满足条件.}
\step{当 $a\neq0$ 时，要使不等式 $ax^2+ax+1>0$ 对一切 $x\in\R$ 恒成立，}
\step{需 $\begin{cases}a>0\\\Delta=a^2-4a<0\end{cases}$，解得 $0<a<4$.}
\step{综上，实数 $a$ 的取值范围是 $[0,4)$.}
\solend
\fi

\item 已知集合 $M=\{x\mid (x-a)(x^2-ax+a-1)=0\}$ 各元素之和等于 3，则实数 $a=$\blankline[4.4em].

\ifanswer
\solbegin
\step{方程 $(x-a)(x^2-ax+a-1)=0$，化简得 $(x-a)(x-1)[x-(a-1)]=0$，}
\step{解得 $x_1=a$，$x_2=1$，$x_3=a-1$，}
\step{因为它们是集合 $M$ 中的元素，所以存在互异性.}
\step{当 $x_1\neq x_2\neq x_3$ 时，即 $a+1+a-1=3$，解得 $a=\dfrac32$，此时 $M=\dfrac32,1,\dfrac12$，满足.}
\step{当 $a=1$ 时，$1+0\neq3$，不合题意；}
\step{当 $a=a-1$ 时，不符合题意；}
\step{当 $1=a-1$ 时，$a=2$，此时 $M=\{2,1\}$ 满足题意.}
\step{故 $a$ 的取值为 $2$ 或 $\dfrac32$.}
\solend
\fi

\item 已知 $a>b>1$ 且 $a^2-6a-2m+1=0$，$b^2-6b-2m+1=0$，则实数 $m$ 的取值范围是\blankline[4.4em].

\ifanswer
\solbegin
\step{由题意得 $a$、$b$ 是函数 $f(x)=x^2-6x-2m+1$ 两个大于 1 的不同零点，}
\step{所以 $\begin{cases}f(1)>0\\f(3)<0\end{cases}$，即 $\begin{cases}-4-2m>0\\-8-2m<0\end{cases}$，解得 $-4<m<-2$，}
\step{故 $m$ 的取值范围为 $(-4,-2)$.}
\solend
\fi

\end{enumerate}

\bigsec{二、选择题}

\begin{enumerate}
\setcounter{enumi}{12}

\item “$x<2$” 是 “$x^2<4$” 的\mbox{（\quad）}条件

\keepwithprev
A. 充分不必要\qquad\qquad B. 必要不充分

C. 充要\qquad\qquad\qquad\quad D. 既不充分也不必要

\ifanswer
\solbegin
\step{解 $x^2<4$，得 $-2<x<2$，}
\step{由 $x<2$ 得不出 $-2<x<2$，而由 $-2<x<2$ 能得出 $x<2$，}
\step{所以 “$x<2$” 是 “$x^2<4$” 的必要不充分条件. 故选 $B$.}
\solend
\fi

\item 已知集合 $A=\{x\mid x<-2\}$，$B=\{x\mid x\geqslant a\}$，且 $A\cup B=\R$，则实数 $a$ 的取值范围是\mbox{（\quad）}

\keepwithprev
A. $a\leqslant-2$\qquad\qquad B. $a<-2$\qquad\qquad C. $a\geqslant-2$\qquad\qquad D. $a>-2$

\ifanswer
\solbegin
\step{集合 $A=\{x\mid x<-2\}$，$B=\{x\mid x\geqslant a\}$，且 $A\cup B=\R$，}
\step{则实数 $a$ 的取值范围是 $a\leqslant-2$. 故选 $A$.}
\solend
\fi

\item 下列关于 $x,y$ 的方程组 $\begin{cases}y=kx+1\\y=2kx+3\end{cases}$ 的解的说法中正确的是\mbox{（\quad）}

\keepwithprev
A. 该方程一定有唯一解\qquad\qquad B. 该方程没有解

C. $k\neq0$ 时，方程组有唯一解\qquad\qquad D. $k=0$ 时，方程组有无数组解

\ifanswer
\solbegin
\step{由题意得 $kx+1=2kx+3$，即 $kx=-2$，}
\step{当 $k=0$ 时，$0=-2$ 不成立，方程组无解；}
\step{当 $k\neq0$ 时，$x=-\dfrac2k$，方程组有唯一解. 故选 $C$.}
\solend
\fi

\item 设非空集合 $S=\{x\mid m\leqslant x\leqslant n\}$ 满足：当 $x\in S$ 时，有 $x^2\in S$. 给出如下三个结论：①若 $m=1$，则 $S=\{1\}$；②若 $m=-\dfrac12$，则 $\dfrac14\leqslant n\leqslant1$；③若 $n=\dfrac12$，则 $-\dfrac{\sqrt2}{2}\leqslant m\leqslant0$. 其中正确命题的个数为\mbox{（\quad）}

\keepwithprev
A. $3$\qquad\qquad B. $2$\qquad\qquad C. $1$\qquad\qquad D. $0$

\ifanswer
\solbegin
\step{由定义设非空集合 $S=\{x\mid m\leqslant x\leqslant n\}$ 满足当 $x\in S$ 时，有 $x^2\in S$，}
\step{符合定义的参数 $m$ 的值一定大于等于 1 或小于等于 0，}
\step{惟如此才能保证 $m\in S$ 时，有 $m^2\in S$ 即 $m^2\geqslant m$，}
\step{符合条件的 $n$ 的值一定大于等于 0，小于等于 1，}
\step{惟如此才能保证 $n\in S$ 时，有 $n^2\in S$ 即 $n^2\leqslant n$，对各个命题进行判断：}
\step{对于① $m=1$，$m^2=1\in S$，故必有 $\begin{cases}n^2\leqslant n\\n\geqslant1\end{cases}$，得 $n=1$，$S=\{1\}$，}
\step{② $m=-\dfrac12$，$m^2=\dfrac14\in S$，则 $\begin{cases}n^2\leqslant n\\\dfrac14\leqslant n\end{cases}$，解得 $\dfrac14\leqslant n\leqslant1$；}
\step{对于③，若 $n=\dfrac12$，则 $\begin{cases}\dfrac12\geqslant m\\m^2\geqslant m\\\dfrac12\geqslant m^2\end{cases}$，解得 $-\dfrac{\sqrt2}{2}\leqslant m\leqslant0$，}
\step{所以正确命题有 3 个. 故选 $A$.}
\solend
\fi

\end{enumerate}

\bigsec{三、解答题}

\begin{enumerate}
\setcounter{enumi}{16}

\item 已知 $a$、$b$ 是实数.

（1）求证：$a^2+b^2\geqslant2a-2b-2$，并指出等号成立的条件；

（2）若 $ab=1$，求 $a^2+4b^2$ 的最小值.

\ifanswer
\solbegin
\step{（1）因为 $a^2+b^2-(2a-2b-2)=a^2+b^2-2a+2b+2$\dots\dots2 分}
\step{$=(a-1)^2+(b+1)^2\geqslant0$\dots\dots4 分}
\step{所以 $a^2+b^2\geqslant2a-2b-2$.}
\step{当且仅当 $a=1,b=-1$ 时，不等式中等号成立.\dots\dots7 分}
\step{（2）$a^2+4b^2=a^2+(2b)^2\geqslant2\cdot a\cdot(2b)=4ab=4$，\dots\dots5 分}
\step{当且仅当 $a=2b$，即 $\begin{cases}a=\sqrt2\\b=\dfrac{\sqrt2}{2}\end{cases}$ 或 $\begin{cases}a=-\sqrt2\\b=-\dfrac{\sqrt2}{2}\end{cases}$ 时，不等式中等号成立，}
\step{所以 $a^2+4b^2$ 的最小值为 4.\dots\dots7 分}
\solend
\fi

\item 已知关于 $x$ 的方程 $x^2+(m-3)x+2=0$，其中 $m\in\R$，有两个实数根 $\alpha$、$\beta$.

（1）若 $m=8$，求 $\dfrac1\alpha+\dfrac1\beta$ 的值；

（2）若 $\alpha^2+\beta^2=\alpha\beta+3$，求实数 $m$ 的值.

\ifanswer
\solbegin
\step{（1）若 $m=8$，则关于 $x$ 的方程为 $x^2+5x+2=0$，}
\step{则 $\alpha+\beta=-5$，$\alpha\beta=2$，所以 $\dfrac1\alpha+\dfrac1\beta=\dfrac{\alpha+\beta}{\alpha\beta}=-\dfrac52$；}
\step{（2）由题意得 $\alpha+\beta=-(m-3)=3-m$，$\alpha\beta=2$，}
\step{又 $\alpha^2+\beta^2=\alpha\beta+3$，则 $(\alpha+\beta)^2-2\alpha\beta=\alpha\beta+3$，}
\step{即 $(\alpha+\beta)^2=3\alpha\beta+3$，所以 $(3-m)^2=3\times2+3=9$，解得 $m=0$ 或 $6$.}
\solend
\fi

\item 设集合 $A=\{x\mid\lvert x-a\rvert<2\}$，$B=\left\{x\mid\dfrac{x+4}{x+2}>2\right\}$.

（1）若 $a=1$，试用区间表示集合 $A$、$B$；

（2）若 $A\cup\overline{B}=\R$，求实数 $a$ 的取值范围.

\ifanswer
\solbegin
\step{（1）当 $a=1$ 时，由 $\lvert x-1\rvert<2$ 得 $-2<x-1<2$，解得 $-1<x<3$，}
\step{所以 $A=(-1,3)$.}
\step{由 $\dfrac{x+4}{x+2}>2$ 得 $2-\dfrac{x+4}{x+2}=\dfrac{x}{x+2}<0$，则有 $x(x+2)<0$，解得 $-2<x<0$，}
\step{所以 $B=(-2,0)$；}
\step{（2）由 $\lvert x-a\rvert<2$ 得 $-2<x-a<2$，解得 $a-2<x<a+2$，}
\step{所以 $A=(a-2,a+2)$，}
\step{由（1）得 $B=(-2,0)$，则 $\overline{B}=\{x\mid x\geqslant0$ 或 $x\leqslant-2\}$，}
\step{由于 $A\cup\overline{B}=\R$，所以 $\begin{cases}a+2\geqslant0\\a-2\leqslant-2\end{cases}$，解得 $-2\leqslant a\leqslant0$，}
\step{所以实数 $a$ 的取值范围是 $[-2,0]$.}
\solend
\fi

\item 某服装公司生产的衬衫每件定价为 80 元，在某城市年销售 8 万件. 现该公司计划在该市招收代理商来销售衬衫，以降低管理和营销成本. 已知代理商收取的代理费为总销售金额的 $r\%$（其中 $0<r\leqslant10$）.

（1）若单件衬衫的利润保持不变，求招代理商后每件衬衫的定价为多少元（用 $r$ 表示）；

（2）已知将每件衬衫的价格提高到 $\dfrac{80}{1-r\%}$ 元后，每年的销售量会减少 $0.62r$ 万件. 若代理商每年收取的代理费不少于 16 万元，求 $r$ 的取值范围.

\ifanswer
\solbegin
\step{（1）设招代理商后每件衬衫的定价为 $x$ 元，代理商收取总销售金额的 $r\%$，}
\step{因此公司每件衬衫的实际净收入为 $x(1-r\%)$ 元，}
\step{单件利润保持不变，等价于净收入仍为 80 元，}
\step{于是 $x(1-r\%)=80$，解得 $x=\dfrac{80}{1-r\%}$，因为 $0<r\leqslant10$，有 $1-r\%\geqslant0.9>0$，}
\step{分母恒不为零，该定价有意义，所以招代理商后每件衬衫的定价为 $\dfrac{80}{1-r\%}$ 元；}
\step{（2）价格提高到 $\dfrac{80}{1-r\%}$ 元，年销售量减少 $0.62r$ 万件，}
\step{故年销售量为 $8-0.62r$ 万件，$8-0.62r>0$，}
\step{由题意得 $\dfrac{80}{1-r\%}\cdot(8-0.62r)\cdot r\%\geqslant16$，}
\step{即 $\dfrac{80}{1-\dfrac r{100}}\cdot(8-0.62r)\cdot\dfrac r{100}=\dfrac{80r(8-0.62r)}{100-r}\geqslant16$，}
\step{由 $0<r\leqslant10$ 得 $100-r>0$，整理得 $80r(8-0.62r)\geqslant16(100-r)$，}
\step{整理得 $31r^2-410r+1000\leqslant0$，解得 $\dfrac{100}{31}\leqslant r\leqslant10$，}
\step{当 $\dfrac{100}{31}\leqslant r\leqslant10$ 时，$8-0.62r>0$ 且 $1-r\%>0$，}
\step{均符合实际经营条件，所以 $r$ 的取值范围是 $\left[\dfrac{100}{31},10\right]$.}
\solend
\fi

\item 定义 $\min\{a_1,a_2,\cdots,a_n\}$ 为 $n$ 个实数 $a_1,a_2,\cdots,a_n$ 中的最小数，$\max\{a_1,a_2,\cdots,a_n\}$ 为 $n$ 个实数 $a_1,a_2,\cdots,a_n$ 中的最大数.

（1）设 $a,b$ 都是正实数，且 $a+b=1$，求 $\max\left\{ab,\dfrac14\right\}$；

（2）解不等式：$\min\{x+1,x^2+3,\lvert x-1\rvert\}>2x-3$；

（3）设 $a,b$ 都是正实数，求 $\max\left\{a+\dfrac1b,\dfrac2a+b\right\}$ 的最小值.

\ifanswer
\solbegin
\step{（1）$ab\leqslant\dfrac{(a+b)^2}{4}=\dfrac14$，所以 $\max\left\{ab,\dfrac14\right\}=\dfrac14$.}
\step{（2）因为 $x^2+3>x+1$，$x^2+3>\lvert x-1\rvert$，}
\step{所以 $\min\{x+1,x^2+3,\lvert x-1\rvert\}=\min\{x+1,\lvert x-1\rvert\}$.}
\step{当 $x\geqslant0$ 时，$\lvert x-1\rvert\leqslant x+1$，所以 $\lvert x-1\rvert>2x-3$，解得 $0\leqslant x<2$；}
\step{当 $x<0$ 时，$\lvert x-1\rvert>x+1$，所以 $x+1>2x-3$，解得 $x<0$.}
\step{综上，不等式的解集为 $(-\infty,2)$.}
\step{（3）设 $M=\max\left\{a+\dfrac1b,\dfrac2a+b\right\}$，则 $M\geqslant a+\dfrac1b$，$M\geqslant\dfrac2a+b$，}
\step{所以 $M^2\geqslant\left(a+\dfrac1b\right)\left(\dfrac2a+b\right)=3+ab+\dfrac2{ab}\geqslant3+2\sqrt2$，}
\step{所以 $M$ 的最小值为 $\sqrt2+1$，当 $a=\sqrt2$，$b=1$ 时取等号.}
\solend
\fi

\end{enumerate}

\end{document}
"
% 紧凑版：xelatex "\def\iscompact{1}% 高一126_建青实验_期中.tex
%   —— 高一上数学"2026-2027 年建青实验高一上期中"（试卷版/答案版）
% 试卷版：xelatex 高一126_建青实验_期中.tex
% 答案版：xelatex "\def\isanswer{1}% 高一126_建青实验_期中.tex
%   —— 高一上数学"2026-2027 年建青实验高一上期中"（试卷版/答案版）
% 试卷版：xelatex 高一126_建青实验_期中.tex
% 答案版：xelatex "\def\isanswer{1}\input{高一126_建青实验_期中.tex}"
% 紧凑版：xelatex "\def\iscompact{1}\input{高一126_建青实验_期中.tex}"
%
% 原始材料：微信公众号"魔都数学加油站"文章，作者破晓，2026-11-11 18:35 发布；
% 连载"27学年高一126"，原文标题"27学年高一126——2026-2027年上海市建青实验高一上期中"；
% 原文链接：https://mp.weixin.qq.com/s/HpxpBJu4d_wOnUCIbZSQiA。
% 正文为 8 张 PNG 页（01、05、08 为 2560x3620，02、03、04、06、07 为 4762x6735），
% 题目下印有【解析】或【答案】，为讲评版。题目、解析均据原图转录；
% 卷面粉红色水印"魔都数学加油站"及公众号页脚未录入。全卷无插图。
%
% 卷首标题：原卷印作"2026-2027 年建青实验高一上期中"（公众号标题中的"上海市"
% 卷面标题行没有，本稿照卷面），年份区间按全稿惯例排 en dash；原卷未印副标题，
% 本卷为期中卷、跨集合与常用逻辑用语／不等式／函数三章，按题目内容标为
% "集合与常用逻辑用语、不等式、函数"（与 125 同为期中卷，口径相同）：
% 集合与逻辑见 Q5、Q8、Q9、Q11、Q13、Q14、Q16、Q19，不等式见 Q6、Q7、Q10、
% Q17、Q20、Q21，函数（含指数对数、二次函数零点）见 Q3、Q4、Q12、Q15、Q18。
%
% 与原卷的排版差异：
%   ① 原文从第 3 题起，填空题编号为 3—12、选择题为 13—16、解答题为 17—21，本稿保留原题号；
%   ② 原卷每题下直接印【解析】／【答案】，本稿按全稿惯例置于答案版；第 4、5 题只印
%      【答案】，另提 \ans{}；第 13—16 题答案均在【解析】末句"故选 …"，题干括号内留空、
%      不另提 \ans{}；第 3 题印一行【解析】，按 \solbegin+\step 排；第 17 题解析保留
%      原卷评分点（照 高一54／64 成例排 \dots\dots N 分）；
%   ③ 原卷未印副标题，本稿按"知识模块"口径补出（见上"卷首标题"）；
%   ④ 原卷填空横线按全稿惯例排为 \blankline[4.4em]；选择题括号排 \mbox{（\quad）}、
%      选项段前按惯例加 \keepwithprev；第 14 题原卷括号另起一行，本稿并入题干行末；
%      解答题不留作答区（原卷为讲评版，与 109—125 同）；
%   ⑤ 原卷数集排斜体 $R$，按全稿惯例改 \R；绝对值排 \lvert…\rvert，集合描述竖线排 \mid；
%      补集记号照原卷上划线排 \overline{B}；不等号按全稿惯例排 \leqslant／\geqslant；
%      子题标号原卷一律用全角括号（1）（2）（3），本稿照录（\step 行内亦同，照 高一121 成例）。
%
% 判别式字形：第 10 题解析 $\Delta=a^2-4a<0$ 原卷排 $\Delta$（等腰、高），
%   照录 \Delta（与 高一124／125 同题排 $\triangle$（等边、扁）不同，逐处照本卷字形）。
%
% 子集符号：第 8 题题干与解析两处 $S\subset P$ 均印单线、无下划线，照录 \subset；
%   全卷未出现带下划线的 \subseteq。
%
% 原卷存疑照录：
%   第 10 题题干作 $ax^2-ax+1>0$，解析却作 $ax^2+ax+1>0$（一次项符号不一致），照录、未统一；
%   第 11 题解析"此时 $M=\dfrac32,1,\dfrac12$，满足"未用集合花括号（题干 $M$ 为集合），照录；
%   第 11 题解析"因为它们是集合 $M$ 中的元素，所以存在互异性"（"存在互异性"疑为
%     "满足互异性"之误），又"当 $x_1\neq x_2\neq x_3$ 时"连写两个不等号，均照录；
%   第 16 题解析两处"惟如此才能保证"（用"惟"字）及"符合定义的参数 $m$ 的值一定大于等于 1
%     或小于等于 0"的措辞，照录；
%   第 17 题 (1) 解析"当且仅当 $a=1,b=-1$ 时"用半角逗号、第 21 题 (3) 解析
%     "当 $a=\sqrt2$，$b=1$ 时取等号"用全角逗号，均照录；
%   第 19 题解析由 $\dfrac{x+4}{x+2}>2$ 得 "$2-\dfrac{x+4}{x+2}=\dfrac{x}{x+2}<0$"
%     （以 2 减左式，与常见的"左式减 2"次序相反），照录；
%   第 15 题选项 A、B 作"该方程…"、C、D 作"方程组…"，照录；
%   第 20 题题干作"招收代理商"，(1) 却作"招代理商后"，照录。
%
% 行数复核：逐题数 \step 条数，合计 91 条（Q3：1；Q4：0；Q5：0；Q6：3；Q7：2；
% Q8：5；Q9：5；Q10：4；Q11：8；Q12：3；Q13：3；Q14：2；Q15：3；Q16：9；
% Q17：7；Q18：5；Q19：9；Q20：13；Q21：9），与原卷印刷行逐题相等
% （cases 块作一个印刷行计）。第 4、5 题只印【答案】，提 \ans{}，不计入 \step 账；
% 四道选择题答案均在解析末句，括号内留空、不另提 \ans{}。
%
\documentclass[a4paper,11pt]{article}
\usepackage{common}

\begin{document}

\examtitlest[3]{2026--2027 年建青实验高一上期中}{集合与常用逻辑用语、不等式、函数}

\bigsec{一、填空题}

\begin{enumerate}
\setcounter{enumi}{2}

\item 已知 $\log_2 x=3$，则 $x=$\blankline[4.4em].

\ifanswer
\solbegin
\step{$x=2^3=8$.}
\solend
\fi

\item 将 $\sqrt{a^2\sqrt a}$（其中 $a>0$）化为有理数指数幂的形式为\blankline[4.4em].

\ifanswer
\ans{$a^{\frac54}$}
\fi

\item 用反证法证明命题：“已知 $m$、$n\in\R$，若 $m+n\leqslant5$，则 $m\leqslant2$ 或 $n\leqslant3$” 时应假设\blankline[4.4em].

\ifanswer
\ans{$m>2$ 且 $n>3$}
\fi

\item 若 $x\in\R$，则 $\lvert x+1\rvert+\lvert x-1\rvert$ 的最小值为\blankline[4.4em].

\ifanswer
\solbegin
\step{$\lvert x+1\rvert+\lvert x-1\rvert\geqslant\lvert(x+1)-(x-1)\rvert=2$，}
\step{当且仅当 $(x+1)(x-1)\leqslant0$，即 $-1\leqslant x\leqslant1$ 时取等号，}
\step{所以 $\lvert x+1\rvert+\lvert x-1\rvert$ 的最小值为 2.}
\solend
\fi

\item 已知 $x>1$，则 $x+\dfrac1{x-1}$ 的最小值是\blankline[4.4em].

\ifanswer
\solbegin
\step{$x-1>0$，则 $x+\dfrac1{x-1}=(x-1)+\dfrac1{x-1}+1\geqslant2\sqrt{(x-1)\cdot\dfrac1{x-1}}+1=3$，}
\step{当且仅当 $x=2$ 时取等号，故 $x+\dfrac1{x-1}$ 的最小值为 3.}
\solend
\fi

\item 已知集合 $P=\{x\mid x^2+x-6=0\}$，$S=\{x\mid ax+1=0\}$ 且 $S\subset P$，则 $a$ 的值为\blankline[4.4em].

\ifanswer
\solbegin
\step{集合 $P=\{-3,2\}$，集合 $S$ 中至多有一个元素，}
\step{若集合 $S$ 为空集，即 $a=0$ 时，显然满足条件 $S\subset P$，故 $a=0$.}
\step{若集合 $S$ 非空集，即 $a\neq0$，此时 $S=\left\{-\dfrac1a\right\}$，}
\step{若 $-\dfrac1a=-3$，则 $a=\dfrac13$，若 $-\dfrac1a=2$，则 $a=-\dfrac12$；}
\step{故 $a$ 的取值为 $0$、$\dfrac13$、$-\dfrac12$.}
\solend
\fi

\item 已知 $p:\begin{cases}a>4\\b>4\end{cases}$，$q:\begin{cases}a+b>8\\ab>16\end{cases}$，则 $p$ 是 $q$ 的\blankline[4.4em]条件.

\ifanswer
\solbegin
\step{若 $a>4$，$b>4$，则 $a+b>4+4=8$，$ab>4\times4=16$，}
\step{所以 $p$ 成立一定能推出 $q$ 成立，充分性成立，}
\step{取 $a=2$，$b=9$，$a+b=11>8$，$ab=18>16$，满足 $q$；}
\step{但是 $a=2<4$，不满足 $p$，所以由 $q$ 不能推出 $p$，必要性不成立，}
\step{综上所述，$p$ 是 $q$ 的充分不必要条件.}
\solend
\fi

\item 若关于 $x$ 的不等式 $ax^2-ax+1>0$ 对一切实数 $x$ 恒成立，则实数 $a$ 的取值范围是\blankline[4.4em].

\ifanswer
\solbegin
\step{当 $a=0$ 时，不等式即 $1>0$，满足条件.}
\step{当 $a\neq0$ 时，要使不等式 $ax^2+ax+1>0$ 对一切 $x\in\R$ 恒成立，}
\step{需 $\begin{cases}a>0\\\Delta=a^2-4a<0\end{cases}$，解得 $0<a<4$.}
\step{综上，实数 $a$ 的取值范围是 $[0,4)$.}
\solend
\fi

\item 已知集合 $M=\{x\mid (x-a)(x^2-ax+a-1)=0\}$ 各元素之和等于 3，则实数 $a=$\blankline[4.4em].

\ifanswer
\solbegin
\step{方程 $(x-a)(x^2-ax+a-1)=0$，化简得 $(x-a)(x-1)[x-(a-1)]=0$，}
\step{解得 $x_1=a$，$x_2=1$，$x_3=a-1$，}
\step{因为它们是集合 $M$ 中的元素，所以存在互异性.}
\step{当 $x_1\neq x_2\neq x_3$ 时，即 $a+1+a-1=3$，解得 $a=\dfrac32$，此时 $M=\dfrac32,1,\dfrac12$，满足.}
\step{当 $a=1$ 时，$1+0\neq3$，不合题意；}
\step{当 $a=a-1$ 时，不符合题意；}
\step{当 $1=a-1$ 时，$a=2$，此时 $M=\{2,1\}$ 满足题意.}
\step{故 $a$ 的取值为 $2$ 或 $\dfrac32$.}
\solend
\fi

\item 已知 $a>b>1$ 且 $a^2-6a-2m+1=0$，$b^2-6b-2m+1=0$，则实数 $m$ 的取值范围是\blankline[4.4em].

\ifanswer
\solbegin
\step{由题意得 $a$、$b$ 是函数 $f(x)=x^2-6x-2m+1$ 两个大于 1 的不同零点，}
\step{所以 $\begin{cases}f(1)>0\\f(3)<0\end{cases}$，即 $\begin{cases}-4-2m>0\\-8-2m<0\end{cases}$，解得 $-4<m<-2$，}
\step{故 $m$ 的取值范围为 $(-4,-2)$.}
\solend
\fi

\end{enumerate}

\bigsec{二、选择题}

\begin{enumerate}
\setcounter{enumi}{12}

\item “$x<2$” 是 “$x^2<4$” 的\mbox{（\quad）}条件

\keepwithprev
A. 充分不必要\qquad\qquad B. 必要不充分

C. 充要\qquad\qquad\qquad\quad D. 既不充分也不必要

\ifanswer
\solbegin
\step{解 $x^2<4$，得 $-2<x<2$，}
\step{由 $x<2$ 得不出 $-2<x<2$，而由 $-2<x<2$ 能得出 $x<2$，}
\step{所以 “$x<2$” 是 “$x^2<4$” 的必要不充分条件. 故选 $B$.}
\solend
\fi

\item 已知集合 $A=\{x\mid x<-2\}$，$B=\{x\mid x\geqslant a\}$，且 $A\cup B=\R$，则实数 $a$ 的取值范围是\mbox{（\quad）}

\keepwithprev
A. $a\leqslant-2$\qquad\qquad B. $a<-2$\qquad\qquad C. $a\geqslant-2$\qquad\qquad D. $a>-2$

\ifanswer
\solbegin
\step{集合 $A=\{x\mid x<-2\}$，$B=\{x\mid x\geqslant a\}$，且 $A\cup B=\R$，}
\step{则实数 $a$ 的取值范围是 $a\leqslant-2$. 故选 $A$.}
\solend
\fi

\item 下列关于 $x,y$ 的方程组 $\begin{cases}y=kx+1\\y=2kx+3\end{cases}$ 的解的说法中正确的是\mbox{（\quad）}

\keepwithprev
A. 该方程一定有唯一解\qquad\qquad B. 该方程没有解

C. $k\neq0$ 时，方程组有唯一解\qquad\qquad D. $k=0$ 时，方程组有无数组解

\ifanswer
\solbegin
\step{由题意得 $kx+1=2kx+3$，即 $kx=-2$，}
\step{当 $k=0$ 时，$0=-2$ 不成立，方程组无解；}
\step{当 $k\neq0$ 时，$x=-\dfrac2k$，方程组有唯一解. 故选 $C$.}
\solend
\fi

\item 设非空集合 $S=\{x\mid m\leqslant x\leqslant n\}$ 满足：当 $x\in S$ 时，有 $x^2\in S$. 给出如下三个结论：①若 $m=1$，则 $S=\{1\}$；②若 $m=-\dfrac12$，则 $\dfrac14\leqslant n\leqslant1$；③若 $n=\dfrac12$，则 $-\dfrac{\sqrt2}{2}\leqslant m\leqslant0$. 其中正确命题的个数为\mbox{（\quad）}

\keepwithprev
A. $3$\qquad\qquad B. $2$\qquad\qquad C. $1$\qquad\qquad D. $0$

\ifanswer
\solbegin
\step{由定义设非空集合 $S=\{x\mid m\leqslant x\leqslant n\}$ 满足当 $x\in S$ 时，有 $x^2\in S$，}
\step{符合定义的参数 $m$ 的值一定大于等于 1 或小于等于 0，}
\step{惟如此才能保证 $m\in S$ 时，有 $m^2\in S$ 即 $m^2\geqslant m$，}
\step{符合条件的 $n$ 的值一定大于等于 0，小于等于 1，}
\step{惟如此才能保证 $n\in S$ 时，有 $n^2\in S$ 即 $n^2\leqslant n$，对各个命题进行判断：}
\step{对于① $m=1$，$m^2=1\in S$，故必有 $\begin{cases}n^2\leqslant n\\n\geqslant1\end{cases}$，得 $n=1$，$S=\{1\}$，}
\step{② $m=-\dfrac12$，$m^2=\dfrac14\in S$，则 $\begin{cases}n^2\leqslant n\\\dfrac14\leqslant n\end{cases}$，解得 $\dfrac14\leqslant n\leqslant1$；}
\step{对于③，若 $n=\dfrac12$，则 $\begin{cases}\dfrac12\geqslant m\\m^2\geqslant m\\\dfrac12\geqslant m^2\end{cases}$，解得 $-\dfrac{\sqrt2}{2}\leqslant m\leqslant0$，}
\step{所以正确命题有 3 个. 故选 $A$.}
\solend
\fi

\end{enumerate}

\bigsec{三、解答题}

\begin{enumerate}
\setcounter{enumi}{16}

\item 已知 $a$、$b$ 是实数.

（1）求证：$a^2+b^2\geqslant2a-2b-2$，并指出等号成立的条件；

（2）若 $ab=1$，求 $a^2+4b^2$ 的最小值.

\ifanswer
\solbegin
\step{（1）因为 $a^2+b^2-(2a-2b-2)=a^2+b^2-2a+2b+2$\dots\dots2 分}
\step{$=(a-1)^2+(b+1)^2\geqslant0$\dots\dots4 分}
\step{所以 $a^2+b^2\geqslant2a-2b-2$.}
\step{当且仅当 $a=1,b=-1$ 时，不等式中等号成立.\dots\dots7 分}
\step{（2）$a^2+4b^2=a^2+(2b)^2\geqslant2\cdot a\cdot(2b)=4ab=4$，\dots\dots5 分}
\step{当且仅当 $a=2b$，即 $\begin{cases}a=\sqrt2\\b=\dfrac{\sqrt2}{2}\end{cases}$ 或 $\begin{cases}a=-\sqrt2\\b=-\dfrac{\sqrt2}{2}\end{cases}$ 时，不等式中等号成立，}
\step{所以 $a^2+4b^2$ 的最小值为 4.\dots\dots7 分}
\solend
\fi

\item 已知关于 $x$ 的方程 $x^2+(m-3)x+2=0$，其中 $m\in\R$，有两个实数根 $\alpha$、$\beta$.

（1）若 $m=8$，求 $\dfrac1\alpha+\dfrac1\beta$ 的值；

（2）若 $\alpha^2+\beta^2=\alpha\beta+3$，求实数 $m$ 的值.

\ifanswer
\solbegin
\step{（1）若 $m=8$，则关于 $x$ 的方程为 $x^2+5x+2=0$，}
\step{则 $\alpha+\beta=-5$，$\alpha\beta=2$，所以 $\dfrac1\alpha+\dfrac1\beta=\dfrac{\alpha+\beta}{\alpha\beta}=-\dfrac52$；}
\step{（2）由题意得 $\alpha+\beta=-(m-3)=3-m$，$\alpha\beta=2$，}
\step{又 $\alpha^2+\beta^2=\alpha\beta+3$，则 $(\alpha+\beta)^2-2\alpha\beta=\alpha\beta+3$，}
\step{即 $(\alpha+\beta)^2=3\alpha\beta+3$，所以 $(3-m)^2=3\times2+3=9$，解得 $m=0$ 或 $6$.}
\solend
\fi

\item 设集合 $A=\{x\mid\lvert x-a\rvert<2\}$，$B=\left\{x\mid\dfrac{x+4}{x+2}>2\right\}$.

（1）若 $a=1$，试用区间表示集合 $A$、$B$；

（2）若 $A\cup\overline{B}=\R$，求实数 $a$ 的取值范围.

\ifanswer
\solbegin
\step{（1）当 $a=1$ 时，由 $\lvert x-1\rvert<2$ 得 $-2<x-1<2$，解得 $-1<x<3$，}
\step{所以 $A=(-1,3)$.}
\step{由 $\dfrac{x+4}{x+2}>2$ 得 $2-\dfrac{x+4}{x+2}=\dfrac{x}{x+2}<0$，则有 $x(x+2)<0$，解得 $-2<x<0$，}
\step{所以 $B=(-2,0)$；}
\step{（2）由 $\lvert x-a\rvert<2$ 得 $-2<x-a<2$，解得 $a-2<x<a+2$，}
\step{所以 $A=(a-2,a+2)$，}
\step{由（1）得 $B=(-2,0)$，则 $\overline{B}=\{x\mid x\geqslant0$ 或 $x\leqslant-2\}$，}
\step{由于 $A\cup\overline{B}=\R$，所以 $\begin{cases}a+2\geqslant0\\a-2\leqslant-2\end{cases}$，解得 $-2\leqslant a\leqslant0$，}
\step{所以实数 $a$ 的取值范围是 $[-2,0]$.}
\solend
\fi

\item 某服装公司生产的衬衫每件定价为 80 元，在某城市年销售 8 万件. 现该公司计划在该市招收代理商来销售衬衫，以降低管理和营销成本. 已知代理商收取的代理费为总销售金额的 $r\%$（其中 $0<r\leqslant10$）.

（1）若单件衬衫的利润保持不变，求招代理商后每件衬衫的定价为多少元（用 $r$ 表示）；

（2）已知将每件衬衫的价格提高到 $\dfrac{80}{1-r\%}$ 元后，每年的销售量会减少 $0.62r$ 万件. 若代理商每年收取的代理费不少于 16 万元，求 $r$ 的取值范围.

\ifanswer
\solbegin
\step{（1）设招代理商后每件衬衫的定价为 $x$ 元，代理商收取总销售金额的 $r\%$，}
\step{因此公司每件衬衫的实际净收入为 $x(1-r\%)$ 元，}
\step{单件利润保持不变，等价于净收入仍为 80 元，}
\step{于是 $x(1-r\%)=80$，解得 $x=\dfrac{80}{1-r\%}$，因为 $0<r\leqslant10$，有 $1-r\%\geqslant0.9>0$，}
\step{分母恒不为零，该定价有意义，所以招代理商后每件衬衫的定价为 $\dfrac{80}{1-r\%}$ 元；}
\step{（2）价格提高到 $\dfrac{80}{1-r\%}$ 元，年销售量减少 $0.62r$ 万件，}
\step{故年销售量为 $8-0.62r$ 万件，$8-0.62r>0$，}
\step{由题意得 $\dfrac{80}{1-r\%}\cdot(8-0.62r)\cdot r\%\geqslant16$，}
\step{即 $\dfrac{80}{1-\dfrac r{100}}\cdot(8-0.62r)\cdot\dfrac r{100}=\dfrac{80r(8-0.62r)}{100-r}\geqslant16$，}
\step{由 $0<r\leqslant10$ 得 $100-r>0$，整理得 $80r(8-0.62r)\geqslant16(100-r)$，}
\step{整理得 $31r^2-410r+1000\leqslant0$，解得 $\dfrac{100}{31}\leqslant r\leqslant10$，}
\step{当 $\dfrac{100}{31}\leqslant r\leqslant10$ 时，$8-0.62r>0$ 且 $1-r\%>0$，}
\step{均符合实际经营条件，所以 $r$ 的取值范围是 $\left[\dfrac{100}{31},10\right]$.}
\solend
\fi

\item 定义 $\min\{a_1,a_2,\cdots,a_n\}$ 为 $n$ 个实数 $a_1,a_2,\cdots,a_n$ 中的最小数，$\max\{a_1,a_2,\cdots,a_n\}$ 为 $n$ 个实数 $a_1,a_2,\cdots,a_n$ 中的最大数.

（1）设 $a,b$ 都是正实数，且 $a+b=1$，求 $\max\left\{ab,\dfrac14\right\}$；

（2）解不等式：$\min\{x+1,x^2+3,\lvert x-1\rvert\}>2x-3$；

（3）设 $a,b$ 都是正实数，求 $\max\left\{a+\dfrac1b,\dfrac2a+b\right\}$ 的最小值.

\ifanswer
\solbegin
\step{（1）$ab\leqslant\dfrac{(a+b)^2}{4}=\dfrac14$，所以 $\max\left\{ab,\dfrac14\right\}=\dfrac14$.}
\step{（2）因为 $x^2+3>x+1$，$x^2+3>\lvert x-1\rvert$，}
\step{所以 $\min\{x+1,x^2+3,\lvert x-1\rvert\}=\min\{x+1,\lvert x-1\rvert\}$.}
\step{当 $x\geqslant0$ 时，$\lvert x-1\rvert\leqslant x+1$，所以 $\lvert x-1\rvert>2x-3$，解得 $0\leqslant x<2$；}
\step{当 $x<0$ 时，$\lvert x-1\rvert>x+1$，所以 $x+1>2x-3$，解得 $x<0$.}
\step{综上，不等式的解集为 $(-\infty,2)$.}
\step{（3）设 $M=\max\left\{a+\dfrac1b,\dfrac2a+b\right\}$，则 $M\geqslant a+\dfrac1b$，$M\geqslant\dfrac2a+b$，}
\step{所以 $M^2\geqslant\left(a+\dfrac1b\right)\left(\dfrac2a+b\right)=3+ab+\dfrac2{ab}\geqslant3+2\sqrt2$，}
\step{所以 $M$ 的最小值为 $\sqrt2+1$，当 $a=\sqrt2$，$b=1$ 时取等号.}
\solend
\fi

\end{enumerate}

\end{document}
"
% 紧凑版：xelatex "\def\iscompact{1}% 高一126_建青实验_期中.tex
%   —— 高一上数学"2026-2027 年建青实验高一上期中"（试卷版/答案版）
% 试卷版：xelatex 高一126_建青实验_期中.tex
% 答案版：xelatex "\def\isanswer{1}\input{高一126_建青实验_期中.tex}"
% 紧凑版：xelatex "\def\iscompact{1}\input{高一126_建青实验_期中.tex}"
%
% 原始材料：微信公众号"魔都数学加油站"文章，作者破晓，2026-11-11 18:35 发布；
% 连载"27学年高一126"，原文标题"27学年高一126——2026-2027年上海市建青实验高一上期中"；
% 原文链接：https://mp.weixin.qq.com/s/HpxpBJu4d_wOnUCIbZSQiA。
% 正文为 8 张 PNG 页（01、05、08 为 2560x3620，02、03、04、06、07 为 4762x6735），
% 题目下印有【解析】或【答案】，为讲评版。题目、解析均据原图转录；
% 卷面粉红色水印"魔都数学加油站"及公众号页脚未录入。全卷无插图。
%
% 卷首标题：原卷印作"2026-2027 年建青实验高一上期中"（公众号标题中的"上海市"
% 卷面标题行没有，本稿照卷面），年份区间按全稿惯例排 en dash；原卷未印副标题，
% 本卷为期中卷、跨集合与常用逻辑用语／不等式／函数三章，按题目内容标为
% "集合与常用逻辑用语、不等式、函数"（与 125 同为期中卷，口径相同）：
% 集合与逻辑见 Q5、Q8、Q9、Q11、Q13、Q14、Q16、Q19，不等式见 Q6、Q7、Q10、
% Q17、Q20、Q21，函数（含指数对数、二次函数零点）见 Q3、Q4、Q12、Q15、Q18。
%
% 与原卷的排版差异：
%   ① 原文从第 3 题起，填空题编号为 3—12、选择题为 13—16、解答题为 17—21，本稿保留原题号；
%   ② 原卷每题下直接印【解析】／【答案】，本稿按全稿惯例置于答案版；第 4、5 题只印
%      【答案】，另提 \ans{}；第 13—16 题答案均在【解析】末句"故选 …"，题干括号内留空、
%      不另提 \ans{}；第 3 题印一行【解析】，按 \solbegin+\step 排；第 17 题解析保留
%      原卷评分点（照 高一54／64 成例排 \dots\dots N 分）；
%   ③ 原卷未印副标题，本稿按"知识模块"口径补出（见上"卷首标题"）；
%   ④ 原卷填空横线按全稿惯例排为 \blankline[4.4em]；选择题括号排 \mbox{（\quad）}、
%      选项段前按惯例加 \keepwithprev；第 14 题原卷括号另起一行，本稿并入题干行末；
%      解答题不留作答区（原卷为讲评版，与 109—125 同）；
%   ⑤ 原卷数集排斜体 $R$，按全稿惯例改 \R；绝对值排 \lvert…\rvert，集合描述竖线排 \mid；
%      补集记号照原卷上划线排 \overline{B}；不等号按全稿惯例排 \leqslant／\geqslant；
%      子题标号原卷一律用全角括号（1）（2）（3），本稿照录（\step 行内亦同，照 高一121 成例）。
%
% 判别式字形：第 10 题解析 $\Delta=a^2-4a<0$ 原卷排 $\Delta$（等腰、高），
%   照录 \Delta（与 高一124／125 同题排 $\triangle$（等边、扁）不同，逐处照本卷字形）。
%
% 子集符号：第 8 题题干与解析两处 $S\subset P$ 均印单线、无下划线，照录 \subset；
%   全卷未出现带下划线的 \subseteq。
%
% 原卷存疑照录：
%   第 10 题题干作 $ax^2-ax+1>0$，解析却作 $ax^2+ax+1>0$（一次项符号不一致），照录、未统一；
%   第 11 题解析"此时 $M=\dfrac32,1,\dfrac12$，满足"未用集合花括号（题干 $M$ 为集合），照录；
%   第 11 题解析"因为它们是集合 $M$ 中的元素，所以存在互异性"（"存在互异性"疑为
%     "满足互异性"之误），又"当 $x_1\neq x_2\neq x_3$ 时"连写两个不等号，均照录；
%   第 16 题解析两处"惟如此才能保证"（用"惟"字）及"符合定义的参数 $m$ 的值一定大于等于 1
%     或小于等于 0"的措辞，照录；
%   第 17 题 (1) 解析"当且仅当 $a=1,b=-1$ 时"用半角逗号、第 21 题 (3) 解析
%     "当 $a=\sqrt2$，$b=1$ 时取等号"用全角逗号，均照录；
%   第 19 题解析由 $\dfrac{x+4}{x+2}>2$ 得 "$2-\dfrac{x+4}{x+2}=\dfrac{x}{x+2}<0$"
%     （以 2 减左式，与常见的"左式减 2"次序相反），照录；
%   第 15 题选项 A、B 作"该方程…"、C、D 作"方程组…"，照录；
%   第 20 题题干作"招收代理商"，(1) 却作"招代理商后"，照录。
%
% 行数复核：逐题数 \step 条数，合计 91 条（Q3：1；Q4：0；Q5：0；Q6：3；Q7：2；
% Q8：5；Q9：5；Q10：4；Q11：8；Q12：3；Q13：3；Q14：2；Q15：3；Q16：9；
% Q17：7；Q18：5；Q19：9；Q20：13；Q21：9），与原卷印刷行逐题相等
% （cases 块作一个印刷行计）。第 4、5 题只印【答案】，提 \ans{}，不计入 \step 账；
% 四道选择题答案均在解析末句，括号内留空、不另提 \ans{}。
%
\documentclass[a4paper,11pt]{article}
\usepackage{common}

\begin{document}

\examtitlest[3]{2026--2027 年建青实验高一上期中}{集合与常用逻辑用语、不等式、函数}

\bigsec{一、填空题}

\begin{enumerate}
\setcounter{enumi}{2}

\item 已知 $\log_2 x=3$，则 $x=$\blankline[4.4em].

\ifanswer
\solbegin
\step{$x=2^3=8$.}
\solend
\fi

\item 将 $\sqrt{a^2\sqrt a}$（其中 $a>0$）化为有理数指数幂的形式为\blankline[4.4em].

\ifanswer
\ans{$a^{\frac54}$}
\fi

\item 用反证法证明命题：“已知 $m$、$n\in\R$，若 $m+n\leqslant5$，则 $m\leqslant2$ 或 $n\leqslant3$” 时应假设\blankline[4.4em].

\ifanswer
\ans{$m>2$ 且 $n>3$}
\fi

\item 若 $x\in\R$，则 $\lvert x+1\rvert+\lvert x-1\rvert$ 的最小值为\blankline[4.4em].

\ifanswer
\solbegin
\step{$\lvert x+1\rvert+\lvert x-1\rvert\geqslant\lvert(x+1)-(x-1)\rvert=2$，}
\step{当且仅当 $(x+1)(x-1)\leqslant0$，即 $-1\leqslant x\leqslant1$ 时取等号，}
\step{所以 $\lvert x+1\rvert+\lvert x-1\rvert$ 的最小值为 2.}
\solend
\fi

\item 已知 $x>1$，则 $x+\dfrac1{x-1}$ 的最小值是\blankline[4.4em].

\ifanswer
\solbegin
\step{$x-1>0$，则 $x+\dfrac1{x-1}=(x-1)+\dfrac1{x-1}+1\geqslant2\sqrt{(x-1)\cdot\dfrac1{x-1}}+1=3$，}
\step{当且仅当 $x=2$ 时取等号，故 $x+\dfrac1{x-1}$ 的最小值为 3.}
\solend
\fi

\item 已知集合 $P=\{x\mid x^2+x-6=0\}$，$S=\{x\mid ax+1=0\}$ 且 $S\subset P$，则 $a$ 的值为\blankline[4.4em].

\ifanswer
\solbegin
\step{集合 $P=\{-3,2\}$，集合 $S$ 中至多有一个元素，}
\step{若集合 $S$ 为空集，即 $a=0$ 时，显然满足条件 $S\subset P$，故 $a=0$.}
\step{若集合 $S$ 非空集，即 $a\neq0$，此时 $S=\left\{-\dfrac1a\right\}$，}
\step{若 $-\dfrac1a=-3$，则 $a=\dfrac13$，若 $-\dfrac1a=2$，则 $a=-\dfrac12$；}
\step{故 $a$ 的取值为 $0$、$\dfrac13$、$-\dfrac12$.}
\solend
\fi

\item 已知 $p:\begin{cases}a>4\\b>4\end{cases}$，$q:\begin{cases}a+b>8\\ab>16\end{cases}$，则 $p$ 是 $q$ 的\blankline[4.4em]条件.

\ifanswer
\solbegin
\step{若 $a>4$，$b>4$，则 $a+b>4+4=8$，$ab>4\times4=16$，}
\step{所以 $p$ 成立一定能推出 $q$ 成立，充分性成立，}
\step{取 $a=2$，$b=9$，$a+b=11>8$，$ab=18>16$，满足 $q$；}
\step{但是 $a=2<4$，不满足 $p$，所以由 $q$ 不能推出 $p$，必要性不成立，}
\step{综上所述，$p$ 是 $q$ 的充分不必要条件.}
\solend
\fi

\item 若关于 $x$ 的不等式 $ax^2-ax+1>0$ 对一切实数 $x$ 恒成立，则实数 $a$ 的取值范围是\blankline[4.4em].

\ifanswer
\solbegin
\step{当 $a=0$ 时，不等式即 $1>0$，满足条件.}
\step{当 $a\neq0$ 时，要使不等式 $ax^2+ax+1>0$ 对一切 $x\in\R$ 恒成立，}
\step{需 $\begin{cases}a>0\\\Delta=a^2-4a<0\end{cases}$，解得 $0<a<4$.}
\step{综上，实数 $a$ 的取值范围是 $[0,4)$.}
\solend
\fi

\item 已知集合 $M=\{x\mid (x-a)(x^2-ax+a-1)=0\}$ 各元素之和等于 3，则实数 $a=$\blankline[4.4em].

\ifanswer
\solbegin
\step{方程 $(x-a)(x^2-ax+a-1)=0$，化简得 $(x-a)(x-1)[x-(a-1)]=0$，}
\step{解得 $x_1=a$，$x_2=1$，$x_3=a-1$，}
\step{因为它们是集合 $M$ 中的元素，所以存在互异性.}
\step{当 $x_1\neq x_2\neq x_3$ 时，即 $a+1+a-1=3$，解得 $a=\dfrac32$，此时 $M=\dfrac32,1,\dfrac12$，满足.}
\step{当 $a=1$ 时，$1+0\neq3$，不合题意；}
\step{当 $a=a-1$ 时，不符合题意；}
\step{当 $1=a-1$ 时，$a=2$，此时 $M=\{2,1\}$ 满足题意.}
\step{故 $a$ 的取值为 $2$ 或 $\dfrac32$.}
\solend
\fi

\item 已知 $a>b>1$ 且 $a^2-6a-2m+1=0$，$b^2-6b-2m+1=0$，则实数 $m$ 的取值范围是\blankline[4.4em].

\ifanswer
\solbegin
\step{由题意得 $a$、$b$ 是函数 $f(x)=x^2-6x-2m+1$ 两个大于 1 的不同零点，}
\step{所以 $\begin{cases}f(1)>0\\f(3)<0\end{cases}$，即 $\begin{cases}-4-2m>0\\-8-2m<0\end{cases}$，解得 $-4<m<-2$，}
\step{故 $m$ 的取值范围为 $(-4,-2)$.}
\solend
\fi

\end{enumerate}

\bigsec{二、选择题}

\begin{enumerate}
\setcounter{enumi}{12}

\item “$x<2$” 是 “$x^2<4$” 的\mbox{（\quad）}条件

\keepwithprev
A. 充分不必要\qquad\qquad B. 必要不充分

C. 充要\qquad\qquad\qquad\quad D. 既不充分也不必要

\ifanswer
\solbegin
\step{解 $x^2<4$，得 $-2<x<2$，}
\step{由 $x<2$ 得不出 $-2<x<2$，而由 $-2<x<2$ 能得出 $x<2$，}
\step{所以 “$x<2$” 是 “$x^2<4$” 的必要不充分条件. 故选 $B$.}
\solend
\fi

\item 已知集合 $A=\{x\mid x<-2\}$，$B=\{x\mid x\geqslant a\}$，且 $A\cup B=\R$，则实数 $a$ 的取值范围是\mbox{（\quad）}

\keepwithprev
A. $a\leqslant-2$\qquad\qquad B. $a<-2$\qquad\qquad C. $a\geqslant-2$\qquad\qquad D. $a>-2$

\ifanswer
\solbegin
\step{集合 $A=\{x\mid x<-2\}$，$B=\{x\mid x\geqslant a\}$，且 $A\cup B=\R$，}
\step{则实数 $a$ 的取值范围是 $a\leqslant-2$. 故选 $A$.}
\solend
\fi

\item 下列关于 $x,y$ 的方程组 $\begin{cases}y=kx+1\\y=2kx+3\end{cases}$ 的解的说法中正确的是\mbox{（\quad）}

\keepwithprev
A. 该方程一定有唯一解\qquad\qquad B. 该方程没有解

C. $k\neq0$ 时，方程组有唯一解\qquad\qquad D. $k=0$ 时，方程组有无数组解

\ifanswer
\solbegin
\step{由题意得 $kx+1=2kx+3$，即 $kx=-2$，}
\step{当 $k=0$ 时，$0=-2$ 不成立，方程组无解；}
\step{当 $k\neq0$ 时，$x=-\dfrac2k$，方程组有唯一解. 故选 $C$.}
\solend
\fi

\item 设非空集合 $S=\{x\mid m\leqslant x\leqslant n\}$ 满足：当 $x\in S$ 时，有 $x^2\in S$. 给出如下三个结论：①若 $m=1$，则 $S=\{1\}$；②若 $m=-\dfrac12$，则 $\dfrac14\leqslant n\leqslant1$；③若 $n=\dfrac12$，则 $-\dfrac{\sqrt2}{2}\leqslant m\leqslant0$. 其中正确命题的个数为\mbox{（\quad）}

\keepwithprev
A. $3$\qquad\qquad B. $2$\qquad\qquad C. $1$\qquad\qquad D. $0$

\ifanswer
\solbegin
\step{由定义设非空集合 $S=\{x\mid m\leqslant x\leqslant n\}$ 满足当 $x\in S$ 时，有 $x^2\in S$，}
\step{符合定义的参数 $m$ 的值一定大于等于 1 或小于等于 0，}
\step{惟如此才能保证 $m\in S$ 时，有 $m^2\in S$ 即 $m^2\geqslant m$，}
\step{符合条件的 $n$ 的值一定大于等于 0，小于等于 1，}
\step{惟如此才能保证 $n\in S$ 时，有 $n^2\in S$ 即 $n^2\leqslant n$，对各个命题进行判断：}
\step{对于① $m=1$，$m^2=1\in S$，故必有 $\begin{cases}n^2\leqslant n\\n\geqslant1\end{cases}$，得 $n=1$，$S=\{1\}$，}
\step{② $m=-\dfrac12$，$m^2=\dfrac14\in S$，则 $\begin{cases}n^2\leqslant n\\\dfrac14\leqslant n\end{cases}$，解得 $\dfrac14\leqslant n\leqslant1$；}
\step{对于③，若 $n=\dfrac12$，则 $\begin{cases}\dfrac12\geqslant m\\m^2\geqslant m\\\dfrac12\geqslant m^2\end{cases}$，解得 $-\dfrac{\sqrt2}{2}\leqslant m\leqslant0$，}
\step{所以正确命题有 3 个. 故选 $A$.}
\solend
\fi

\end{enumerate}

\bigsec{三、解答题}

\begin{enumerate}
\setcounter{enumi}{16}

\item 已知 $a$、$b$ 是实数.

（1）求证：$a^2+b^2\geqslant2a-2b-2$，并指出等号成立的条件；

（2）若 $ab=1$，求 $a^2+4b^2$ 的最小值.

\ifanswer
\solbegin
\step{（1）因为 $a^2+b^2-(2a-2b-2)=a^2+b^2-2a+2b+2$\dots\dots2 分}
\step{$=(a-1)^2+(b+1)^2\geqslant0$\dots\dots4 分}
\step{所以 $a^2+b^2\geqslant2a-2b-2$.}
\step{当且仅当 $a=1,b=-1$ 时，不等式中等号成立.\dots\dots7 分}
\step{（2）$a^2+4b^2=a^2+(2b)^2\geqslant2\cdot a\cdot(2b)=4ab=4$，\dots\dots5 分}
\step{当且仅当 $a=2b$，即 $\begin{cases}a=\sqrt2\\b=\dfrac{\sqrt2}{2}\end{cases}$ 或 $\begin{cases}a=-\sqrt2\\b=-\dfrac{\sqrt2}{2}\end{cases}$ 时，不等式中等号成立，}
\step{所以 $a^2+4b^2$ 的最小值为 4.\dots\dots7 分}
\solend
\fi

\item 已知关于 $x$ 的方程 $x^2+(m-3)x+2=0$，其中 $m\in\R$，有两个实数根 $\alpha$、$\beta$.

（1）若 $m=8$，求 $\dfrac1\alpha+\dfrac1\beta$ 的值；

（2）若 $\alpha^2+\beta^2=\alpha\beta+3$，求实数 $m$ 的值.

\ifanswer
\solbegin
\step{（1）若 $m=8$，则关于 $x$ 的方程为 $x^2+5x+2=0$，}
\step{则 $\alpha+\beta=-5$，$\alpha\beta=2$，所以 $\dfrac1\alpha+\dfrac1\beta=\dfrac{\alpha+\beta}{\alpha\beta}=-\dfrac52$；}
\step{（2）由题意得 $\alpha+\beta=-(m-3)=3-m$，$\alpha\beta=2$，}
\step{又 $\alpha^2+\beta^2=\alpha\beta+3$，则 $(\alpha+\beta)^2-2\alpha\beta=\alpha\beta+3$，}
\step{即 $(\alpha+\beta)^2=3\alpha\beta+3$，所以 $(3-m)^2=3\times2+3=9$，解得 $m=0$ 或 $6$.}
\solend
\fi

\item 设集合 $A=\{x\mid\lvert x-a\rvert<2\}$，$B=\left\{x\mid\dfrac{x+4}{x+2}>2\right\}$.

（1）若 $a=1$，试用区间表示集合 $A$、$B$；

（2）若 $A\cup\overline{B}=\R$，求实数 $a$ 的取值范围.

\ifanswer
\solbegin
\step{（1）当 $a=1$ 时，由 $\lvert x-1\rvert<2$ 得 $-2<x-1<2$，解得 $-1<x<3$，}
\step{所以 $A=(-1,3)$.}
\step{由 $\dfrac{x+4}{x+2}>2$ 得 $2-\dfrac{x+4}{x+2}=\dfrac{x}{x+2}<0$，则有 $x(x+2)<0$，解得 $-2<x<0$，}
\step{所以 $B=(-2,0)$；}
\step{（2）由 $\lvert x-a\rvert<2$ 得 $-2<x-a<2$，解得 $a-2<x<a+2$，}
\step{所以 $A=(a-2,a+2)$，}
\step{由（1）得 $B=(-2,0)$，则 $\overline{B}=\{x\mid x\geqslant0$ 或 $x\leqslant-2\}$，}
\step{由于 $A\cup\overline{B}=\R$，所以 $\begin{cases}a+2\geqslant0\\a-2\leqslant-2\end{cases}$，解得 $-2\leqslant a\leqslant0$，}
\step{所以实数 $a$ 的取值范围是 $[-2,0]$.}
\solend
\fi

\item 某服装公司生产的衬衫每件定价为 80 元，在某城市年销售 8 万件. 现该公司计划在该市招收代理商来销售衬衫，以降低管理和营销成本. 已知代理商收取的代理费为总销售金额的 $r\%$（其中 $0<r\leqslant10$）.

（1）若单件衬衫的利润保持不变，求招代理商后每件衬衫的定价为多少元（用 $r$ 表示）；

（2）已知将每件衬衫的价格提高到 $\dfrac{80}{1-r\%}$ 元后，每年的销售量会减少 $0.62r$ 万件. 若代理商每年收取的代理费不少于 16 万元，求 $r$ 的取值范围.

\ifanswer
\solbegin
\step{（1）设招代理商后每件衬衫的定价为 $x$ 元，代理商收取总销售金额的 $r\%$，}
\step{因此公司每件衬衫的实际净收入为 $x(1-r\%)$ 元，}
\step{单件利润保持不变，等价于净收入仍为 80 元，}
\step{于是 $x(1-r\%)=80$，解得 $x=\dfrac{80}{1-r\%}$，因为 $0<r\leqslant10$，有 $1-r\%\geqslant0.9>0$，}
\step{分母恒不为零，该定价有意义，所以招代理商后每件衬衫的定价为 $\dfrac{80}{1-r\%}$ 元；}
\step{（2）价格提高到 $\dfrac{80}{1-r\%}$ 元，年销售量减少 $0.62r$ 万件，}
\step{故年销售量为 $8-0.62r$ 万件，$8-0.62r>0$，}
\step{由题意得 $\dfrac{80}{1-r\%}\cdot(8-0.62r)\cdot r\%\geqslant16$，}
\step{即 $\dfrac{80}{1-\dfrac r{100}}\cdot(8-0.62r)\cdot\dfrac r{100}=\dfrac{80r(8-0.62r)}{100-r}\geqslant16$，}
\step{由 $0<r\leqslant10$ 得 $100-r>0$，整理得 $80r(8-0.62r)\geqslant16(100-r)$，}
\step{整理得 $31r^2-410r+1000\leqslant0$，解得 $\dfrac{100}{31}\leqslant r\leqslant10$，}
\step{当 $\dfrac{100}{31}\leqslant r\leqslant10$ 时，$8-0.62r>0$ 且 $1-r\%>0$，}
\step{均符合实际经营条件，所以 $r$ 的取值范围是 $\left[\dfrac{100}{31},10\right]$.}
\solend
\fi

\item 定义 $\min\{a_1,a_2,\cdots,a_n\}$ 为 $n$ 个实数 $a_1,a_2,\cdots,a_n$ 中的最小数，$\max\{a_1,a_2,\cdots,a_n\}$ 为 $n$ 个实数 $a_1,a_2,\cdots,a_n$ 中的最大数.

（1）设 $a,b$ 都是正实数，且 $a+b=1$，求 $\max\left\{ab,\dfrac14\right\}$；

（2）解不等式：$\min\{x+1,x^2+3,\lvert x-1\rvert\}>2x-3$；

（3）设 $a,b$ 都是正实数，求 $\max\left\{a+\dfrac1b,\dfrac2a+b\right\}$ 的最小值.

\ifanswer
\solbegin
\step{（1）$ab\leqslant\dfrac{(a+b)^2}{4}=\dfrac14$，所以 $\max\left\{ab,\dfrac14\right\}=\dfrac14$.}
\step{（2）因为 $x^2+3>x+1$，$x^2+3>\lvert x-1\rvert$，}
\step{所以 $\min\{x+1,x^2+3,\lvert x-1\rvert\}=\min\{x+1,\lvert x-1\rvert\}$.}
\step{当 $x\geqslant0$ 时，$\lvert x-1\rvert\leqslant x+1$，所以 $\lvert x-1\rvert>2x-3$，解得 $0\leqslant x<2$；}
\step{当 $x<0$ 时，$\lvert x-1\rvert>x+1$，所以 $x+1>2x-3$，解得 $x<0$.}
\step{综上，不等式的解集为 $(-\infty,2)$.}
\step{（3）设 $M=\max\left\{a+\dfrac1b,\dfrac2a+b\right\}$，则 $M\geqslant a+\dfrac1b$，$M\geqslant\dfrac2a+b$，}
\step{所以 $M^2\geqslant\left(a+\dfrac1b\right)\left(\dfrac2a+b\right)=3+ab+\dfrac2{ab}\geqslant3+2\sqrt2$，}
\step{所以 $M$ 的最小值为 $\sqrt2+1$，当 $a=\sqrt2$，$b=1$ 时取等号.}
\solend
\fi

\end{enumerate}

\end{document}
"
%
% 原始材料：微信公众号"魔都数学加油站"文章，作者破晓，2026-11-11 18:35 发布；
% 连载"27学年高一126"，原文标题"27学年高一126——2026-2027年上海市建青实验高一上期中"；
% 原文链接：https://mp.weixin.qq.com/s/HpxpBJu4d_wOnUCIbZSQiA。
% 正文为 8 张 PNG 页（01、05、08 为 2560x3620，02、03、04、06、07 为 4762x6735），
% 题目下印有【解析】或【答案】，为讲评版。题目、解析均据原图转录；
% 卷面粉红色水印"魔都数学加油站"及公众号页脚未录入。全卷无插图。
%
% 卷首标题：原卷印作"2026-2027 年建青实验高一上期中"（公众号标题中的"上海市"
% 卷面标题行没有，本稿照卷面），年份区间按全稿惯例排 en dash；原卷未印副标题，
% 本卷为期中卷、跨集合与常用逻辑用语／不等式／函数三章，按题目内容标为
% "集合与常用逻辑用语、不等式、函数"（与 125 同为期中卷，口径相同）：
% 集合与逻辑见 Q5、Q8、Q9、Q11、Q13、Q14、Q16、Q19，不等式见 Q6、Q7、Q10、
% Q17、Q20、Q21，函数（含指数对数、二次函数零点）见 Q3、Q4、Q12、Q15、Q18。
%
% 与原卷的排版差异：
%   ① 原文从第 3 题起，填空题编号为 3—12、选择题为 13—16、解答题为 17—21，本稿保留原题号；
%   ② 原卷每题下直接印【解析】／【答案】，本稿按全稿惯例置于答案版；第 4、5 题只印
%      【答案】，另提 \ans{}；第 13—16 题答案均在【解析】末句"故选 …"，题干括号内留空、
%      不另提 \ans{}；第 3 题印一行【解析】，按 \solbegin+\step 排；第 17 题解析保留
%      原卷评分点（照 高一54／64 成例排 \dots\dots N 分）；
%   ③ 原卷未印副标题，本稿按"知识模块"口径补出（见上"卷首标题"）；
%   ④ 原卷填空横线按全稿惯例排为 \blankline[4.4em]；选择题括号排 \mbox{（\quad）}、
%      选项段前按惯例加 \keepwithprev；第 14 题原卷括号另起一行，本稿并入题干行末；
%      解答题不留作答区（原卷为讲评版，与 109—125 同）；
%   ⑤ 原卷数集排斜体 $R$，按全稿惯例改 \R；绝对值排 \lvert…\rvert，集合描述竖线排 \mid；
%      补集记号照原卷上划线排 \overline{B}；不等号按全稿惯例排 \leqslant／\geqslant；
%      子题标号原卷一律用全角括号（1）（2）（3），本稿照录（\step 行内亦同，照 高一121 成例）。
%
% 判别式字形：第 10 题解析 $\Delta=a^2-4a<0$ 原卷排 $\Delta$（等腰、高），
%   照录 \Delta（与 高一124／125 同题排 $\triangle$（等边、扁）不同，逐处照本卷字形）。
%
% 子集符号：第 8 题题干与解析两处 $S\subset P$ 均印单线、无下划线，照录 \subset；
%   全卷未出现带下划线的 \subseteq。
%
% 原卷存疑照录：
%   第 10 题题干作 $ax^2-ax+1>0$，解析却作 $ax^2+ax+1>0$（一次项符号不一致），照录、未统一；
%   第 11 题解析"此时 $M=\dfrac32,1,\dfrac12$，满足"未用集合花括号（题干 $M$ 为集合），照录；
%   第 11 题解析"因为它们是集合 $M$ 中的元素，所以存在互异性"（"存在互异性"疑为
%     "满足互异性"之误），又"当 $x_1\neq x_2\neq x_3$ 时"连写两个不等号，均照录；
%   第 16 题解析两处"惟如此才能保证"（用"惟"字）及"符合定义的参数 $m$ 的值一定大于等于 1
%     或小于等于 0"的措辞，照录；
%   第 17 题 (1) 解析"当且仅当 $a=1,b=-1$ 时"用半角逗号、第 21 题 (3) 解析
%     "当 $a=\sqrt2$，$b=1$ 时取等号"用全角逗号，均照录；
%   第 19 题解析由 $\dfrac{x+4}{x+2}>2$ 得 "$2-\dfrac{x+4}{x+2}=\dfrac{x}{x+2}<0$"
%     （以 2 减左式，与常见的"左式减 2"次序相反），照录；
%   第 15 题选项 A、B 作"该方程…"、C、D 作"方程组…"，照录；
%   第 20 题题干作"招收代理商"，(1) 却作"招代理商后"，照录。
%
% 行数复核：逐题数 \step 条数，合计 91 条（Q3：1；Q4：0；Q5：0；Q6：3；Q7：2；
% Q8：5；Q9：5；Q10：4；Q11：8；Q12：3；Q13：3；Q14：2；Q15：3；Q16：9；
% Q17：7；Q18：5；Q19：9；Q20：13；Q21：9），与原卷印刷行逐题相等
% （cases 块作一个印刷行计）。第 4、5 题只印【答案】，提 \ans{}，不计入 \step 账；
% 四道选择题答案均在解析末句，括号内留空、不另提 \ans{}。
%
\documentclass[a4paper,11pt]{article}
\usepackage{common}

\begin{document}

\examtitlest[3]{2026--2027 年建青实验高一上期中}{集合与常用逻辑用语、不等式、函数}

\bigsec{一、填空题}

\begin{enumerate}
\setcounter{enumi}{2}

\item 已知 $\log_2 x=3$，则 $x=$\blankline[4.4em].

\ifanswer
\solbegin
\step{$x=2^3=8$.}
\solend
\fi

\item 将 $\sqrt{a^2\sqrt a}$（其中 $a>0$）化为有理数指数幂的形式为\blankline[4.4em].

\ifanswer
\ans{$a^{\frac54}$}
\fi

\item 用反证法证明命题：“已知 $m$、$n\in\R$，若 $m+n\leqslant5$，则 $m\leqslant2$ 或 $n\leqslant3$” 时应假设\blankline[4.4em].

\ifanswer
\ans{$m>2$ 且 $n>3$}
\fi

\item 若 $x\in\R$，则 $\lvert x+1\rvert+\lvert x-1\rvert$ 的最小值为\blankline[4.4em].

\ifanswer
\solbegin
\step{$\lvert x+1\rvert+\lvert x-1\rvert\geqslant\lvert(x+1)-(x-1)\rvert=2$，}
\step{当且仅当 $(x+1)(x-1)\leqslant0$，即 $-1\leqslant x\leqslant1$ 时取等号，}
\step{所以 $\lvert x+1\rvert+\lvert x-1\rvert$ 的最小值为 2.}
\solend
\fi

\item 已知 $x>1$，则 $x+\dfrac1{x-1}$ 的最小值是\blankline[4.4em].

\ifanswer
\solbegin
\step{$x-1>0$，则 $x+\dfrac1{x-1}=(x-1)+\dfrac1{x-1}+1\geqslant2\sqrt{(x-1)\cdot\dfrac1{x-1}}+1=3$，}
\step{当且仅当 $x=2$ 时取等号，故 $x+\dfrac1{x-1}$ 的最小值为 3.}
\solend
\fi

\item 已知集合 $P=\{x\mid x^2+x-6=0\}$，$S=\{x\mid ax+1=0\}$ 且 $S\subset P$，则 $a$ 的值为\blankline[4.4em].

\ifanswer
\solbegin
\step{集合 $P=\{-3,2\}$，集合 $S$ 中至多有一个元素，}
\step{若集合 $S$ 为空集，即 $a=0$ 时，显然满足条件 $S\subset P$，故 $a=0$.}
\step{若集合 $S$ 非空集，即 $a\neq0$，此时 $S=\left\{-\dfrac1a\right\}$，}
\step{若 $-\dfrac1a=-3$，则 $a=\dfrac13$，若 $-\dfrac1a=2$，则 $a=-\dfrac12$；}
\step{故 $a$ 的取值为 $0$、$\dfrac13$、$-\dfrac12$.}
\solend
\fi

\item 已知 $p:\begin{cases}a>4\\b>4\end{cases}$，$q:\begin{cases}a+b>8\\ab>16\end{cases}$，则 $p$ 是 $q$ 的\blankline[4.4em]条件.

\ifanswer
\solbegin
\step{若 $a>4$，$b>4$，则 $a+b>4+4=8$，$ab>4\times4=16$，}
\step{所以 $p$ 成立一定能推出 $q$ 成立，充分性成立，}
\step{取 $a=2$，$b=9$，$a+b=11>8$，$ab=18>16$，满足 $q$；}
\step{但是 $a=2<4$，不满足 $p$，所以由 $q$ 不能推出 $p$，必要性不成立，}
\step{综上所述，$p$ 是 $q$ 的充分不必要条件.}
\solend
\fi

\item 若关于 $x$ 的不等式 $ax^2-ax+1>0$ 对一切实数 $x$ 恒成立，则实数 $a$ 的取值范围是\blankline[4.4em].

\ifanswer
\solbegin
\step{当 $a=0$ 时，不等式即 $1>0$，满足条件.}
\step{当 $a\neq0$ 时，要使不等式 $ax^2+ax+1>0$ 对一切 $x\in\R$ 恒成立，}
\step{需 $\begin{cases}a>0\\\Delta=a^2-4a<0\end{cases}$，解得 $0<a<4$.}
\step{综上，实数 $a$ 的取值范围是 $[0,4)$.}
\solend
\fi

\item 已知集合 $M=\{x\mid (x-a)(x^2-ax+a-1)=0\}$ 各元素之和等于 3，则实数 $a=$\blankline[4.4em].

\ifanswer
\solbegin
\step{方程 $(x-a)(x^2-ax+a-1)=0$，化简得 $(x-a)(x-1)[x-(a-1)]=0$，}
\step{解得 $x_1=a$，$x_2=1$，$x_3=a-1$，}
\step{因为它们是集合 $M$ 中的元素，所以存在互异性.}
\step{当 $x_1\neq x_2\neq x_3$ 时，即 $a+1+a-1=3$，解得 $a=\dfrac32$，此时 $M=\dfrac32,1,\dfrac12$，满足.}
\step{当 $a=1$ 时，$1+0\neq3$，不合题意；}
\step{当 $a=a-1$ 时，不符合题意；}
\step{当 $1=a-1$ 时，$a=2$，此时 $M=\{2,1\}$ 满足题意.}
\step{故 $a$ 的取值为 $2$ 或 $\dfrac32$.}
\solend
\fi

\item 已知 $a>b>1$ 且 $a^2-6a-2m+1=0$，$b^2-6b-2m+1=0$，则实数 $m$ 的取值范围是\blankline[4.4em].

\ifanswer
\solbegin
\step{由题意得 $a$、$b$ 是函数 $f(x)=x^2-6x-2m+1$ 两个大于 1 的不同零点，}
\step{所以 $\begin{cases}f(1)>0\\f(3)<0\end{cases}$，即 $\begin{cases}-4-2m>0\\-8-2m<0\end{cases}$，解得 $-4<m<-2$，}
\step{故 $m$ 的取值范围为 $(-4,-2)$.}
\solend
\fi

\end{enumerate}

\bigsec{二、选择题}

\begin{enumerate}
\setcounter{enumi}{12}

\item “$x<2$” 是 “$x^2<4$” 的\mbox{（\quad）}条件

\keepwithprev
A. 充分不必要\qquad\qquad B. 必要不充分

C. 充要\qquad\qquad\qquad\quad D. 既不充分也不必要

\ifanswer
\solbegin
\step{解 $x^2<4$，得 $-2<x<2$，}
\step{由 $x<2$ 得不出 $-2<x<2$，而由 $-2<x<2$ 能得出 $x<2$，}
\step{所以 “$x<2$” 是 “$x^2<4$” 的必要不充分条件. 故选 $B$.}
\solend
\fi

\item 已知集合 $A=\{x\mid x<-2\}$，$B=\{x\mid x\geqslant a\}$，且 $A\cup B=\R$，则实数 $a$ 的取值范围是\mbox{（\quad）}

\keepwithprev
A. $a\leqslant-2$\qquad\qquad B. $a<-2$\qquad\qquad C. $a\geqslant-2$\qquad\qquad D. $a>-2$

\ifanswer
\solbegin
\step{集合 $A=\{x\mid x<-2\}$，$B=\{x\mid x\geqslant a\}$，且 $A\cup B=\R$，}
\step{则实数 $a$ 的取值范围是 $a\leqslant-2$. 故选 $A$.}
\solend
\fi

\item 下列关于 $x,y$ 的方程组 $\begin{cases}y=kx+1\\y=2kx+3\end{cases}$ 的解的说法中正确的是\mbox{（\quad）}

\keepwithprev
A. 该方程一定有唯一解\qquad\qquad B. 该方程没有解

C. $k\neq0$ 时，方程组有唯一解\qquad\qquad D. $k=0$ 时，方程组有无数组解

\ifanswer
\solbegin
\step{由题意得 $kx+1=2kx+3$，即 $kx=-2$，}
\step{当 $k=0$ 时，$0=-2$ 不成立，方程组无解；}
\step{当 $k\neq0$ 时，$x=-\dfrac2k$，方程组有唯一解. 故选 $C$.}
\solend
\fi

\item 设非空集合 $S=\{x\mid m\leqslant x\leqslant n\}$ 满足：当 $x\in S$ 时，有 $x^2\in S$. 给出如下三个结论：①若 $m=1$，则 $S=\{1\}$；②若 $m=-\dfrac12$，则 $\dfrac14\leqslant n\leqslant1$；③若 $n=\dfrac12$，则 $-\dfrac{\sqrt2}{2}\leqslant m\leqslant0$. 其中正确命题的个数为\mbox{（\quad）}

\keepwithprev
A. $3$\qquad\qquad B. $2$\qquad\qquad C. $1$\qquad\qquad D. $0$

\ifanswer
\solbegin
\step{由定义设非空集合 $S=\{x\mid m\leqslant x\leqslant n\}$ 满足当 $x\in S$ 时，有 $x^2\in S$，}
\step{符合定义的参数 $m$ 的值一定大于等于 1 或小于等于 0，}
\step{惟如此才能保证 $m\in S$ 时，有 $m^2\in S$ 即 $m^2\geqslant m$，}
\step{符合条件的 $n$ 的值一定大于等于 0，小于等于 1，}
\step{惟如此才能保证 $n\in S$ 时，有 $n^2\in S$ 即 $n^2\leqslant n$，对各个命题进行判断：}
\step{对于① $m=1$，$m^2=1\in S$，故必有 $\begin{cases}n^2\leqslant n\\n\geqslant1\end{cases}$，得 $n=1$，$S=\{1\}$，}
\step{② $m=-\dfrac12$，$m^2=\dfrac14\in S$，则 $\begin{cases}n^2\leqslant n\\\dfrac14\leqslant n\end{cases}$，解得 $\dfrac14\leqslant n\leqslant1$；}
\step{对于③，若 $n=\dfrac12$，则 $\begin{cases}\dfrac12\geqslant m\\m^2\geqslant m\\\dfrac12\geqslant m^2\end{cases}$，解得 $-\dfrac{\sqrt2}{2}\leqslant m\leqslant0$，}
\step{所以正确命题有 3 个. 故选 $A$.}
\solend
\fi

\end{enumerate}

\bigsec{三、解答题}

\begin{enumerate}
\setcounter{enumi}{16}

\item 已知 $a$、$b$ 是实数.

（1）求证：$a^2+b^2\geqslant2a-2b-2$，并指出等号成立的条件；

（2）若 $ab=1$，求 $a^2+4b^2$ 的最小值.

\ifanswer
\solbegin
\step{（1）因为 $a^2+b^2-(2a-2b-2)=a^2+b^2-2a+2b+2$\dots\dots2 分}
\step{$=(a-1)^2+(b+1)^2\geqslant0$\dots\dots4 分}
\step{所以 $a^2+b^2\geqslant2a-2b-2$.}
\step{当且仅当 $a=1,b=-1$ 时，不等式中等号成立.\dots\dots7 分}
\step{（2）$a^2+4b^2=a^2+(2b)^2\geqslant2\cdot a\cdot(2b)=4ab=4$，\dots\dots5 分}
\step{当且仅当 $a=2b$，即 $\begin{cases}a=\sqrt2\\b=\dfrac{\sqrt2}{2}\end{cases}$ 或 $\begin{cases}a=-\sqrt2\\b=-\dfrac{\sqrt2}{2}\end{cases}$ 时，不等式中等号成立，}
\step{所以 $a^2+4b^2$ 的最小值为 4.\dots\dots7 分}
\solend
\fi

\item 已知关于 $x$ 的方程 $x^2+(m-3)x+2=0$，其中 $m\in\R$，有两个实数根 $\alpha$、$\beta$.

（1）若 $m=8$，求 $\dfrac1\alpha+\dfrac1\beta$ 的值；

（2）若 $\alpha^2+\beta^2=\alpha\beta+3$，求实数 $m$ 的值.

\ifanswer
\solbegin
\step{（1）若 $m=8$，则关于 $x$ 的方程为 $x^2+5x+2=0$，}
\step{则 $\alpha+\beta=-5$，$\alpha\beta=2$，所以 $\dfrac1\alpha+\dfrac1\beta=\dfrac{\alpha+\beta}{\alpha\beta}=-\dfrac52$；}
\step{（2）由题意得 $\alpha+\beta=-(m-3)=3-m$，$\alpha\beta=2$，}
\step{又 $\alpha^2+\beta^2=\alpha\beta+3$，则 $(\alpha+\beta)^2-2\alpha\beta=\alpha\beta+3$，}
\step{即 $(\alpha+\beta)^2=3\alpha\beta+3$，所以 $(3-m)^2=3\times2+3=9$，解得 $m=0$ 或 $6$.}
\solend
\fi

\item 设集合 $A=\{x\mid\lvert x-a\rvert<2\}$，$B=\left\{x\mid\dfrac{x+4}{x+2}>2\right\}$.

（1）若 $a=1$，试用区间表示集合 $A$、$B$；

（2）若 $A\cup\overline{B}=\R$，求实数 $a$ 的取值范围.

\ifanswer
\solbegin
\step{（1）当 $a=1$ 时，由 $\lvert x-1\rvert<2$ 得 $-2<x-1<2$，解得 $-1<x<3$，}
\step{所以 $A=(-1,3)$.}
\step{由 $\dfrac{x+4}{x+2}>2$ 得 $2-\dfrac{x+4}{x+2}=\dfrac{x}{x+2}<0$，则有 $x(x+2)<0$，解得 $-2<x<0$，}
\step{所以 $B=(-2,0)$；}
\step{（2）由 $\lvert x-a\rvert<2$ 得 $-2<x-a<2$，解得 $a-2<x<a+2$，}
\step{所以 $A=(a-2,a+2)$，}
\step{由（1）得 $B=(-2,0)$，则 $\overline{B}=\{x\mid x\geqslant0$ 或 $x\leqslant-2\}$，}
\step{由于 $A\cup\overline{B}=\R$，所以 $\begin{cases}a+2\geqslant0\\a-2\leqslant-2\end{cases}$，解得 $-2\leqslant a\leqslant0$，}
\step{所以实数 $a$ 的取值范围是 $[-2,0]$.}
\solend
\fi

\item 某服装公司生产的衬衫每件定价为 80 元，在某城市年销售 8 万件. 现该公司计划在该市招收代理商来销售衬衫，以降低管理和营销成本. 已知代理商收取的代理费为总销售金额的 $r\%$（其中 $0<r\leqslant10$）.

（1）若单件衬衫的利润保持不变，求招代理商后每件衬衫的定价为多少元（用 $r$ 表示）；

（2）已知将每件衬衫的价格提高到 $\dfrac{80}{1-r\%}$ 元后，每年的销售量会减少 $0.62r$ 万件. 若代理商每年收取的代理费不少于 16 万元，求 $r$ 的取值范围.

\ifanswer
\solbegin
\step{（1）设招代理商后每件衬衫的定价为 $x$ 元，代理商收取总销售金额的 $r\%$，}
\step{因此公司每件衬衫的实际净收入为 $x(1-r\%)$ 元，}
\step{单件利润保持不变，等价于净收入仍为 80 元，}
\step{于是 $x(1-r\%)=80$，解得 $x=\dfrac{80}{1-r\%}$，因为 $0<r\leqslant10$，有 $1-r\%\geqslant0.9>0$，}
\step{分母恒不为零，该定价有意义，所以招代理商后每件衬衫的定价为 $\dfrac{80}{1-r\%}$ 元；}
\step{（2）价格提高到 $\dfrac{80}{1-r\%}$ 元，年销售量减少 $0.62r$ 万件，}
\step{故年销售量为 $8-0.62r$ 万件，$8-0.62r>0$，}
\step{由题意得 $\dfrac{80}{1-r\%}\cdot(8-0.62r)\cdot r\%\geqslant16$，}
\step{即 $\dfrac{80}{1-\dfrac r{100}}\cdot(8-0.62r)\cdot\dfrac r{100}=\dfrac{80r(8-0.62r)}{100-r}\geqslant16$，}
\step{由 $0<r\leqslant10$ 得 $100-r>0$，整理得 $80r(8-0.62r)\geqslant16(100-r)$，}
\step{整理得 $31r^2-410r+1000\leqslant0$，解得 $\dfrac{100}{31}\leqslant r\leqslant10$，}
\step{当 $\dfrac{100}{31}\leqslant r\leqslant10$ 时，$8-0.62r>0$ 且 $1-r\%>0$，}
\step{均符合实际经营条件，所以 $r$ 的取值范围是 $\left[\dfrac{100}{31},10\right]$.}
\solend
\fi

\item 定义 $\min\{a_1,a_2,\cdots,a_n\}$ 为 $n$ 个实数 $a_1,a_2,\cdots,a_n$ 中的最小数，$\max\{a_1,a_2,\cdots,a_n\}$ 为 $n$ 个实数 $a_1,a_2,\cdots,a_n$ 中的最大数.

（1）设 $a,b$ 都是正实数，且 $a+b=1$，求 $\max\left\{ab,\dfrac14\right\}$；

（2）解不等式：$\min\{x+1,x^2+3,\lvert x-1\rvert\}>2x-3$；

（3）设 $a,b$ 都是正实数，求 $\max\left\{a+\dfrac1b,\dfrac2a+b\right\}$ 的最小值.

\ifanswer
\solbegin
\step{（1）$ab\leqslant\dfrac{(a+b)^2}{4}=\dfrac14$，所以 $\max\left\{ab,\dfrac14\right\}=\dfrac14$.}
\step{（2）因为 $x^2+3>x+1$，$x^2+3>\lvert x-1\rvert$，}
\step{所以 $\min\{x+1,x^2+3,\lvert x-1\rvert\}=\min\{x+1,\lvert x-1\rvert\}$.}
\step{当 $x\geqslant0$ 时，$\lvert x-1\rvert\leqslant x+1$，所以 $\lvert x-1\rvert>2x-3$，解得 $0\leqslant x<2$；}
\step{当 $x<0$ 时，$\lvert x-1\rvert>x+1$，所以 $x+1>2x-3$，解得 $x<0$.}
\step{综上，不等式的解集为 $(-\infty,2)$.}
\step{（3）设 $M=\max\left\{a+\dfrac1b,\dfrac2a+b\right\}$，则 $M\geqslant a+\dfrac1b$，$M\geqslant\dfrac2a+b$，}
\step{所以 $M^2\geqslant\left(a+\dfrac1b\right)\left(\dfrac2a+b\right)=3+ab+\dfrac2{ab}\geqslant3+2\sqrt2$，}
\step{所以 $M$ 的最小值为 $\sqrt2+1$，当 $a=\sqrt2$，$b=1$ 时取等号.}
\solend
\fi

\end{enumerate}

\end{document}
"
%
% 原始材料：微信公众号"魔都数学加油站"文章，作者破晓，2026-11-11 18:35 发布；
% 连载"27学年高一126"，原文标题"27学年高一126——2026-2027年上海市建青实验高一上期中"；
% 原文链接：https://mp.weixin.qq.com/s/HpxpBJu4d_wOnUCIbZSQiA。
% 正文为 8 张 PNG 页（01、05、08 为 2560x3620，02、03、04、06、07 为 4762x6735），
% 题目下印有【解析】或【答案】，为讲评版。题目、解析均据原图转录；
% 卷面粉红色水印"魔都数学加油站"及公众号页脚未录入。全卷无插图。
%
% 卷首标题：原卷印作"2026-2027 年建青实验高一上期中"（公众号标题中的"上海市"
% 卷面标题行没有，本稿照卷面），年份区间按全稿惯例排 en dash；原卷未印副标题，
% 本卷为期中卷、跨集合与常用逻辑用语／不等式／函数三章，按题目内容标为
% "集合与常用逻辑用语、不等式、函数"（与 125 同为期中卷，口径相同）：
% 集合与逻辑见 Q5、Q8、Q9、Q11、Q13、Q14、Q16、Q19，不等式见 Q6、Q7、Q10、
% Q17、Q20、Q21，函数（含指数对数、二次函数零点）见 Q3、Q4、Q12、Q15、Q18。
%
% 与原卷的排版差异：
%   ① 原文从第 3 题起，填空题编号为 3—12、选择题为 13—16、解答题为 17—21，本稿保留原题号；
%   ② 原卷每题下直接印【解析】／【答案】，本稿按全稿惯例置于答案版；第 4、5 题只印
%      【答案】，另提 \ans{}；第 13—16 题答案均在【解析】末句"故选 …"，题干括号内留空、
%      不另提 \ans{}；第 3 题印一行【解析】，按 \solbegin+\step 排；第 17 题解析保留
%      原卷评分点（照 高一54／64 成例排 \dots\dots N 分）；
%   ③ 原卷未印副标题，本稿按"知识模块"口径补出（见上"卷首标题"）；
%   ④ 原卷填空横线按全稿惯例排为 \blankline[4.4em]；选择题括号排 \mbox{（\quad）}、
%      选项段前按惯例加 \keepwithprev；第 14 题原卷括号另起一行，本稿并入题干行末；
%      解答题不留作答区（原卷为讲评版，与 109—125 同）；
%   ⑤ 原卷数集排斜体 $R$，按全稿惯例改 \R；绝对值排 \lvert…\rvert，集合描述竖线排 \mid；
%      补集记号照原卷上划线排 \overline{B}；不等号按全稿惯例排 \leqslant／\geqslant；
%      子题标号原卷一律用全角括号（1）（2）（3），本稿照录（\step 行内亦同，照 高一121 成例）。
%
% 判别式字形：第 10 题解析 $\Delta=a^2-4a<0$ 原卷排 $\Delta$（等腰、高），
%   照录 \Delta（与 高一124／125 同题排 $\triangle$（等边、扁）不同，逐处照本卷字形）。
%
% 子集符号：第 8 题题干与解析两处 $S\subset P$ 均印单线、无下划线，照录 \subset；
%   全卷未出现带下划线的 \subseteq。
%
% 原卷存疑照录：
%   第 10 题题干作 $ax^2-ax+1>0$，解析却作 $ax^2+ax+1>0$（一次项符号不一致），照录、未统一；
%   第 11 题解析"此时 $M=\dfrac32,1,\dfrac12$，满足"未用集合花括号（题干 $M$ 为集合），照录；
%   第 11 题解析"因为它们是集合 $M$ 中的元素，所以存在互异性"（"存在互异性"疑为
%     "满足互异性"之误），又"当 $x_1\neq x_2\neq x_3$ 时"连写两个不等号，均照录；
%   第 16 题解析两处"惟如此才能保证"（用"惟"字）及"符合定义的参数 $m$ 的值一定大于等于 1
%     或小于等于 0"的措辞，照录；
%   第 17 题 (1) 解析"当且仅当 $a=1,b=-1$ 时"用半角逗号、第 21 题 (3) 解析
%     "当 $a=\sqrt2$，$b=1$ 时取等号"用全角逗号，均照录；
%   第 19 题解析由 $\dfrac{x+4}{x+2}>2$ 得 "$2-\dfrac{x+4}{x+2}=\dfrac{x}{x+2}<0$"
%     （以 2 减左式，与常见的"左式减 2"次序相反），照录；
%   第 15 题选项 A、B 作"该方程…"、C、D 作"方程组…"，照录；
%   第 20 题题干作"招收代理商"，(1) 却作"招代理商后"，照录。
%
% 行数复核：逐题数 \step 条数，合计 91 条（Q3：1；Q4：0；Q5：0；Q6：3；Q7：2；
% Q8：5；Q9：5；Q10：4；Q11：8；Q12：3；Q13：3；Q14：2；Q15：3；Q16：9；
% Q17：7；Q18：5；Q19：9；Q20：13；Q21：9），与原卷印刷行逐题相等
% （cases 块作一个印刷行计）。第 4、5 题只印【答案】，提 \ans{}，不计入 \step 账；
% 四道选择题答案均在解析末句，括号内留空、不另提 \ans{}。
%
\documentclass[a4paper,11pt]{article}
\usepackage{common}

\begin{document}

\examtitlest[3]{2026--2027 年建青实验高一上期中}{集合与常用逻辑用语、不等式、函数}

\bigsec{一、填空题}

\begin{enumerate}
\setcounter{enumi}{2}

\item 已知 $\log_2 x=3$，则 $x=$\blankline[4.4em].

\ifanswer
\solbegin
\step{$x=2^3=8$.}
\solend
\fi

\item 将 $\sqrt{a^2\sqrt a}$（其中 $a>0$）化为有理数指数幂的形式为\blankline[4.4em].

\ifanswer
\ans{$a^{\frac54}$}
\fi

\item 用反证法证明命题：“已知 $m$、$n\in\R$，若 $m+n\leqslant5$，则 $m\leqslant2$ 或 $n\leqslant3$” 时应假设\blankline[4.4em].

\ifanswer
\ans{$m>2$ 且 $n>3$}
\fi

\item 若 $x\in\R$，则 $\lvert x+1\rvert+\lvert x-1\rvert$ 的最小值为\blankline[4.4em].

\ifanswer
\solbegin
\step{$\lvert x+1\rvert+\lvert x-1\rvert\geqslant\lvert(x+1)-(x-1)\rvert=2$，}
\step{当且仅当 $(x+1)(x-1)\leqslant0$，即 $-1\leqslant x\leqslant1$ 时取等号，}
\step{所以 $\lvert x+1\rvert+\lvert x-1\rvert$ 的最小值为 2.}
\solend
\fi

\item 已知 $x>1$，则 $x+\dfrac1{x-1}$ 的最小值是\blankline[4.4em].

\ifanswer
\solbegin
\step{$x-1>0$，则 $x+\dfrac1{x-1}=(x-1)+\dfrac1{x-1}+1\geqslant2\sqrt{(x-1)\cdot\dfrac1{x-1}}+1=3$，}
\step{当且仅当 $x=2$ 时取等号，故 $x+\dfrac1{x-1}$ 的最小值为 3.}
\solend
\fi

\item 已知集合 $P=\{x\mid x^2+x-6=0\}$，$S=\{x\mid ax+1=0\}$ 且 $S\subset P$，则 $a$ 的值为\blankline[4.4em].

\ifanswer
\solbegin
\step{集合 $P=\{-3,2\}$，集合 $S$ 中至多有一个元素，}
\step{若集合 $S$ 为空集，即 $a=0$ 时，显然满足条件 $S\subset P$，故 $a=0$.}
\step{若集合 $S$ 非空集，即 $a\neq0$，此时 $S=\left\{-\dfrac1a\right\}$，}
\step{若 $-\dfrac1a=-3$，则 $a=\dfrac13$，若 $-\dfrac1a=2$，则 $a=-\dfrac12$；}
\step{故 $a$ 的取值为 $0$、$\dfrac13$、$-\dfrac12$.}
\solend
\fi

\item 已知 $p:\begin{cases}a>4\\b>4\end{cases}$，$q:\begin{cases}a+b>8\\ab>16\end{cases}$，则 $p$ 是 $q$ 的\blankline[4.4em]条件.

\ifanswer
\solbegin
\step{若 $a>4$，$b>4$，则 $a+b>4+4=8$，$ab>4\times4=16$，}
\step{所以 $p$ 成立一定能推出 $q$ 成立，充分性成立，}
\step{取 $a=2$，$b=9$，$a+b=11>8$，$ab=18>16$，满足 $q$；}
\step{但是 $a=2<4$，不满足 $p$，所以由 $q$ 不能推出 $p$，必要性不成立，}
\step{综上所述，$p$ 是 $q$ 的充分不必要条件.}
\solend
\fi

\item 若关于 $x$ 的不等式 $ax^2-ax+1>0$ 对一切实数 $x$ 恒成立，则实数 $a$ 的取值范围是\blankline[4.4em].

\ifanswer
\solbegin
\step{当 $a=0$ 时，不等式即 $1>0$，满足条件.}
\step{当 $a\neq0$ 时，要使不等式 $ax^2+ax+1>0$ 对一切 $x\in\R$ 恒成立，}
\step{需 $\begin{cases}a>0\\\Delta=a^2-4a<0\end{cases}$，解得 $0<a<4$.}
\step{综上，实数 $a$ 的取值范围是 $[0,4)$.}
\solend
\fi

\item 已知集合 $M=\{x\mid (x-a)(x^2-ax+a-1)=0\}$ 各元素之和等于 3，则实数 $a=$\blankline[4.4em].

\ifanswer
\solbegin
\step{方程 $(x-a)(x^2-ax+a-1)=0$，化简得 $(x-a)(x-1)[x-(a-1)]=0$，}
\step{解得 $x_1=a$，$x_2=1$，$x_3=a-1$，}
\step{因为它们是集合 $M$ 中的元素，所以存在互异性.}
\step{当 $x_1\neq x_2\neq x_3$ 时，即 $a+1+a-1=3$，解得 $a=\dfrac32$，此时 $M=\dfrac32,1,\dfrac12$，满足.}
\step{当 $a=1$ 时，$1+0\neq3$，不合题意；}
\step{当 $a=a-1$ 时，不符合题意；}
\step{当 $1=a-1$ 时，$a=2$，此时 $M=\{2,1\}$ 满足题意.}
\step{故 $a$ 的取值为 $2$ 或 $\dfrac32$.}
\solend
\fi

\item 已知 $a>b>1$ 且 $a^2-6a-2m+1=0$，$b^2-6b-2m+1=0$，则实数 $m$ 的取值范围是\blankline[4.4em].

\ifanswer
\solbegin
\step{由题意得 $a$、$b$ 是函数 $f(x)=x^2-6x-2m+1$ 两个大于 1 的不同零点，}
\step{所以 $\begin{cases}f(1)>0\\f(3)<0\end{cases}$，即 $\begin{cases}-4-2m>0\\-8-2m<0\end{cases}$，解得 $-4<m<-2$，}
\step{故 $m$ 的取值范围为 $(-4,-2)$.}
\solend
\fi

\end{enumerate}

\bigsec{二、选择题}

\begin{enumerate}
\setcounter{enumi}{12}

\item “$x<2$” 是 “$x^2<4$” 的\mbox{（\quad）}条件

\keepwithprev
A. 充分不必要\qquad\qquad B. 必要不充分

C. 充要\qquad\qquad\qquad\quad D. 既不充分也不必要

\ifanswer
\solbegin
\step{解 $x^2<4$，得 $-2<x<2$，}
\step{由 $x<2$ 得不出 $-2<x<2$，而由 $-2<x<2$ 能得出 $x<2$，}
\step{所以 “$x<2$” 是 “$x^2<4$” 的必要不充分条件. 故选 $B$.}
\solend
\fi

\item 已知集合 $A=\{x\mid x<-2\}$，$B=\{x\mid x\geqslant a\}$，且 $A\cup B=\R$，则实数 $a$ 的取值范围是\mbox{（\quad）}

\keepwithprev
A. $a\leqslant-2$\qquad\qquad B. $a<-2$\qquad\qquad C. $a\geqslant-2$\qquad\qquad D. $a>-2$

\ifanswer
\solbegin
\step{集合 $A=\{x\mid x<-2\}$，$B=\{x\mid x\geqslant a\}$，且 $A\cup B=\R$，}
\step{则实数 $a$ 的取值范围是 $a\leqslant-2$. 故选 $A$.}
\solend
\fi

\item 下列关于 $x,y$ 的方程组 $\begin{cases}y=kx+1\\y=2kx+3\end{cases}$ 的解的说法中正确的是\mbox{（\quad）}

\keepwithprev
A. 该方程一定有唯一解\qquad\qquad B. 该方程没有解

C. $k\neq0$ 时，方程组有唯一解\qquad\qquad D. $k=0$ 时，方程组有无数组解

\ifanswer
\solbegin
\step{由题意得 $kx+1=2kx+3$，即 $kx=-2$，}
\step{当 $k=0$ 时，$0=-2$ 不成立，方程组无解；}
\step{当 $k\neq0$ 时，$x=-\dfrac2k$，方程组有唯一解. 故选 $C$.}
\solend
\fi

\item 设非空集合 $S=\{x\mid m\leqslant x\leqslant n\}$ 满足：当 $x\in S$ 时，有 $x^2\in S$. 给出如下三个结论：①若 $m=1$，则 $S=\{1\}$；②若 $m=-\dfrac12$，则 $\dfrac14\leqslant n\leqslant1$；③若 $n=\dfrac12$，则 $-\dfrac{\sqrt2}{2}\leqslant m\leqslant0$. 其中正确命题的个数为\mbox{（\quad）}

\keepwithprev
A. $3$\qquad\qquad B. $2$\qquad\qquad C. $1$\qquad\qquad D. $0$

\ifanswer
\solbegin
\step{由定义设非空集合 $S=\{x\mid m\leqslant x\leqslant n\}$ 满足当 $x\in S$ 时，有 $x^2\in S$，}
\step{符合定义的参数 $m$ 的值一定大于等于 1 或小于等于 0，}
\step{惟如此才能保证 $m\in S$ 时，有 $m^2\in S$ 即 $m^2\geqslant m$，}
\step{符合条件的 $n$ 的值一定大于等于 0，小于等于 1，}
\step{惟如此才能保证 $n\in S$ 时，有 $n^2\in S$ 即 $n^2\leqslant n$，对各个命题进行判断：}
\step{对于① $m=1$，$m^2=1\in S$，故必有 $\begin{cases}n^2\leqslant n\\n\geqslant1\end{cases}$，得 $n=1$，$S=\{1\}$，}
\step{② $m=-\dfrac12$，$m^2=\dfrac14\in S$，则 $\begin{cases}n^2\leqslant n\\\dfrac14\leqslant n\end{cases}$，解得 $\dfrac14\leqslant n\leqslant1$；}
\step{对于③，若 $n=\dfrac12$，则 $\begin{cases}\dfrac12\geqslant m\\m^2\geqslant m\\\dfrac12\geqslant m^2\end{cases}$，解得 $-\dfrac{\sqrt2}{2}\leqslant m\leqslant0$，}
\step{所以正确命题有 3 个. 故选 $A$.}
\solend
\fi

\end{enumerate}

\bigsec{三、解答题}

\begin{enumerate}
\setcounter{enumi}{16}

\item 已知 $a$、$b$ 是实数.

（1）求证：$a^2+b^2\geqslant2a-2b-2$，并指出等号成立的条件；

（2）若 $ab=1$，求 $a^2+4b^2$ 的最小值.

\ifanswer
\solbegin
\step{（1）因为 $a^2+b^2-(2a-2b-2)=a^2+b^2-2a+2b+2$\dots\dots2 分}
\step{$=(a-1)^2+(b+1)^2\geqslant0$\dots\dots4 分}
\step{所以 $a^2+b^2\geqslant2a-2b-2$.}
\step{当且仅当 $a=1,b=-1$ 时，不等式中等号成立.\dots\dots7 分}
\step{（2）$a^2+4b^2=a^2+(2b)^2\geqslant2\cdot a\cdot(2b)=4ab=4$，\dots\dots5 分}
\step{当且仅当 $a=2b$，即 $\begin{cases}a=\sqrt2\\b=\dfrac{\sqrt2}{2}\end{cases}$ 或 $\begin{cases}a=-\sqrt2\\b=-\dfrac{\sqrt2}{2}\end{cases}$ 时，不等式中等号成立，}
\step{所以 $a^2+4b^2$ 的最小值为 4.\dots\dots7 分}
\solend
\fi

\item 已知关于 $x$ 的方程 $x^2+(m-3)x+2=0$，其中 $m\in\R$，有两个实数根 $\alpha$、$\beta$.

（1）若 $m=8$，求 $\dfrac1\alpha+\dfrac1\beta$ 的值；

（2）若 $\alpha^2+\beta^2=\alpha\beta+3$，求实数 $m$ 的值.

\ifanswer
\solbegin
\step{（1）若 $m=8$，则关于 $x$ 的方程为 $x^2+5x+2=0$，}
\step{则 $\alpha+\beta=-5$，$\alpha\beta=2$，所以 $\dfrac1\alpha+\dfrac1\beta=\dfrac{\alpha+\beta}{\alpha\beta}=-\dfrac52$；}
\step{（2）由题意得 $\alpha+\beta=-(m-3)=3-m$，$\alpha\beta=2$，}
\step{又 $\alpha^2+\beta^2=\alpha\beta+3$，则 $(\alpha+\beta)^2-2\alpha\beta=\alpha\beta+3$，}
\step{即 $(\alpha+\beta)^2=3\alpha\beta+3$，所以 $(3-m)^2=3\times2+3=9$，解得 $m=0$ 或 $6$.}
\solend
\fi

\item 设集合 $A=\{x\mid\lvert x-a\rvert<2\}$，$B=\left\{x\mid\dfrac{x+4}{x+2}>2\right\}$.

（1）若 $a=1$，试用区间表示集合 $A$、$B$；

（2）若 $A\cup\overline{B}=\R$，求实数 $a$ 的取值范围.

\ifanswer
\solbegin
\step{（1）当 $a=1$ 时，由 $\lvert x-1\rvert<2$ 得 $-2<x-1<2$，解得 $-1<x<3$，}
\step{所以 $A=(-1,3)$.}
\step{由 $\dfrac{x+4}{x+2}>2$ 得 $2-\dfrac{x+4}{x+2}=\dfrac{x}{x+2}<0$，则有 $x(x+2)<0$，解得 $-2<x<0$，}
\step{所以 $B=(-2,0)$；}
\step{（2）由 $\lvert x-a\rvert<2$ 得 $-2<x-a<2$，解得 $a-2<x<a+2$，}
\step{所以 $A=(a-2,a+2)$，}
\step{由（1）得 $B=(-2,0)$，则 $\overline{B}=\{x\mid x\geqslant0$ 或 $x\leqslant-2\}$，}
\step{由于 $A\cup\overline{B}=\R$，所以 $\begin{cases}a+2\geqslant0\\a-2\leqslant-2\end{cases}$，解得 $-2\leqslant a\leqslant0$，}
\step{所以实数 $a$ 的取值范围是 $[-2,0]$.}
\solend
\fi

\item 某服装公司生产的衬衫每件定价为 80 元，在某城市年销售 8 万件. 现该公司计划在该市招收代理商来销售衬衫，以降低管理和营销成本. 已知代理商收取的代理费为总销售金额的 $r\%$（其中 $0<r\leqslant10$）.

（1）若单件衬衫的利润保持不变，求招代理商后每件衬衫的定价为多少元（用 $r$ 表示）；

（2）已知将每件衬衫的价格提高到 $\dfrac{80}{1-r\%}$ 元后，每年的销售量会减少 $0.62r$ 万件. 若代理商每年收取的代理费不少于 16 万元，求 $r$ 的取值范围.

\ifanswer
\solbegin
\step{（1）设招代理商后每件衬衫的定价为 $x$ 元，代理商收取总销售金额的 $r\%$，}
\step{因此公司每件衬衫的实际净收入为 $x(1-r\%)$ 元，}
\step{单件利润保持不变，等价于净收入仍为 80 元，}
\step{于是 $x(1-r\%)=80$，解得 $x=\dfrac{80}{1-r\%}$，因为 $0<r\leqslant10$，有 $1-r\%\geqslant0.9>0$，}
\step{分母恒不为零，该定价有意义，所以招代理商后每件衬衫的定价为 $\dfrac{80}{1-r\%}$ 元；}
\step{（2）价格提高到 $\dfrac{80}{1-r\%}$ 元，年销售量减少 $0.62r$ 万件，}
\step{故年销售量为 $8-0.62r$ 万件，$8-0.62r>0$，}
\step{由题意得 $\dfrac{80}{1-r\%}\cdot(8-0.62r)\cdot r\%\geqslant16$，}
\step{即 $\dfrac{80}{1-\dfrac r{100}}\cdot(8-0.62r)\cdot\dfrac r{100}=\dfrac{80r(8-0.62r)}{100-r}\geqslant16$，}
\step{由 $0<r\leqslant10$ 得 $100-r>0$，整理得 $80r(8-0.62r)\geqslant16(100-r)$，}
\step{整理得 $31r^2-410r+1000\leqslant0$，解得 $\dfrac{100}{31}\leqslant r\leqslant10$，}
\step{当 $\dfrac{100}{31}\leqslant r\leqslant10$ 时，$8-0.62r>0$ 且 $1-r\%>0$，}
\step{均符合实际经营条件，所以 $r$ 的取值范围是 $\left[\dfrac{100}{31},10\right]$.}
\solend
\fi

\item 定义 $\min\{a_1,a_2,\cdots,a_n\}$ 为 $n$ 个实数 $a_1,a_2,\cdots,a_n$ 中的最小数，$\max\{a_1,a_2,\cdots,a_n\}$ 为 $n$ 个实数 $a_1,a_2,\cdots,a_n$ 中的最大数.

（1）设 $a,b$ 都是正实数，且 $a+b=1$，求 $\max\left\{ab,\dfrac14\right\}$；

（2）解不等式：$\min\{x+1,x^2+3,\lvert x-1\rvert\}>2x-3$；

（3）设 $a,b$ 都是正实数，求 $\max\left\{a+\dfrac1b,\dfrac2a+b\right\}$ 的最小值.

\ifanswer
\solbegin
\step{（1）$ab\leqslant\dfrac{(a+b)^2}{4}=\dfrac14$，所以 $\max\left\{ab,\dfrac14\right\}=\dfrac14$.}
\step{（2）因为 $x^2+3>x+1$，$x^2+3>\lvert x-1\rvert$，}
\step{所以 $\min\{x+1,x^2+3,\lvert x-1\rvert\}=\min\{x+1,\lvert x-1\rvert\}$.}
\step{当 $x\geqslant0$ 时，$\lvert x-1\rvert\leqslant x+1$，所以 $\lvert x-1\rvert>2x-3$，解得 $0\leqslant x<2$；}
\step{当 $x<0$ 时，$\lvert x-1\rvert>x+1$，所以 $x+1>2x-3$，解得 $x<0$.}
\step{综上，不等式的解集为 $(-\infty,2)$.}
\step{（3）设 $M=\max\left\{a+\dfrac1b,\dfrac2a+b\right\}$，则 $M\geqslant a+\dfrac1b$，$M\geqslant\dfrac2a+b$，}
\step{所以 $M^2\geqslant\left(a+\dfrac1b\right)\left(\dfrac2a+b\right)=3+ab+\dfrac2{ab}\geqslant3+2\sqrt2$，}
\step{所以 $M$ 的最小值为 $\sqrt2+1$，当 $a=\sqrt2$，$b=1$ 时取等号.}
\solend
\fi

\end{enumerate}

\end{document}
"
% 紧凑版：xelatex "\def\iscompact{1}% 高一126_建青实验_期中.tex
%   —— 高一上数学"2026-2027 年建青实验高一上期中"（试卷版/答案版）
% 试卷版：xelatex 高一126_建青实验_期中.tex
% 答案版：xelatex "\def\isanswer{1}% 高一126_建青实验_期中.tex
%   —— 高一上数学"2026-2027 年建青实验高一上期中"（试卷版/答案版）
% 试卷版：xelatex 高一126_建青实验_期中.tex
% 答案版：xelatex "\def\isanswer{1}% 高一126_建青实验_期中.tex
%   —— 高一上数学"2026-2027 年建青实验高一上期中"（试卷版/答案版）
% 试卷版：xelatex 高一126_建青实验_期中.tex
% 答案版：xelatex "\def\isanswer{1}\input{高一126_建青实验_期中.tex}"
% 紧凑版：xelatex "\def\iscompact{1}\input{高一126_建青实验_期中.tex}"
%
% 原始材料：微信公众号"魔都数学加油站"文章，作者破晓，2026-11-11 18:35 发布；
% 连载"27学年高一126"，原文标题"27学年高一126——2026-2027年上海市建青实验高一上期中"；
% 原文链接：https://mp.weixin.qq.com/s/HpxpBJu4d_wOnUCIbZSQiA。
% 正文为 8 张 PNG 页（01、05、08 为 2560x3620，02、03、04、06、07 为 4762x6735），
% 题目下印有【解析】或【答案】，为讲评版。题目、解析均据原图转录；
% 卷面粉红色水印"魔都数学加油站"及公众号页脚未录入。全卷无插图。
%
% 卷首标题：原卷印作"2026-2027 年建青实验高一上期中"（公众号标题中的"上海市"
% 卷面标题行没有，本稿照卷面），年份区间按全稿惯例排 en dash；原卷未印副标题，
% 本卷为期中卷、跨集合与常用逻辑用语／不等式／函数三章，按题目内容标为
% "集合与常用逻辑用语、不等式、函数"（与 125 同为期中卷，口径相同）：
% 集合与逻辑见 Q5、Q8、Q9、Q11、Q13、Q14、Q16、Q19，不等式见 Q6、Q7、Q10、
% Q17、Q20、Q21，函数（含指数对数、二次函数零点）见 Q3、Q4、Q12、Q15、Q18。
%
% 与原卷的排版差异：
%   ① 原文从第 3 题起，填空题编号为 3—12、选择题为 13—16、解答题为 17—21，本稿保留原题号；
%   ② 原卷每题下直接印【解析】／【答案】，本稿按全稿惯例置于答案版；第 4、5 题只印
%      【答案】，另提 \ans{}；第 13—16 题答案均在【解析】末句"故选 …"，题干括号内留空、
%      不另提 \ans{}；第 3 题印一行【解析】，按 \solbegin+\step 排；第 17 题解析保留
%      原卷评分点（照 高一54／64 成例排 \dots\dots N 分）；
%   ③ 原卷未印副标题，本稿按"知识模块"口径补出（见上"卷首标题"）；
%   ④ 原卷填空横线按全稿惯例排为 \blankline[4.4em]；选择题括号排 \mbox{（\quad）}、
%      选项段前按惯例加 \keepwithprev；第 14 题原卷括号另起一行，本稿并入题干行末；
%      解答题不留作答区（原卷为讲评版，与 109—125 同）；
%   ⑤ 原卷数集排斜体 $R$，按全稿惯例改 \R；绝对值排 \lvert…\rvert，集合描述竖线排 \mid；
%      补集记号照原卷上划线排 \overline{B}；不等号按全稿惯例排 \leqslant／\geqslant；
%      子题标号原卷一律用全角括号（1）（2）（3），本稿照录（\step 行内亦同，照 高一121 成例）。
%
% 判别式字形：第 10 题解析 $\Delta=a^2-4a<0$ 原卷排 $\Delta$（等腰、高），
%   照录 \Delta（与 高一124／125 同题排 $\triangle$（等边、扁）不同，逐处照本卷字形）。
%
% 子集符号：第 8 题题干与解析两处 $S\subset P$ 均印单线、无下划线，照录 \subset；
%   全卷未出现带下划线的 \subseteq。
%
% 原卷存疑照录：
%   第 10 题题干作 $ax^2-ax+1>0$，解析却作 $ax^2+ax+1>0$（一次项符号不一致），照录、未统一；
%   第 11 题解析"此时 $M=\dfrac32,1,\dfrac12$，满足"未用集合花括号（题干 $M$ 为集合），照录；
%   第 11 题解析"因为它们是集合 $M$ 中的元素，所以存在互异性"（"存在互异性"疑为
%     "满足互异性"之误），又"当 $x_1\neq x_2\neq x_3$ 时"连写两个不等号，均照录；
%   第 16 题解析两处"惟如此才能保证"（用"惟"字）及"符合定义的参数 $m$ 的值一定大于等于 1
%     或小于等于 0"的措辞，照录；
%   第 17 题 (1) 解析"当且仅当 $a=1,b=-1$ 时"用半角逗号、第 21 题 (3) 解析
%     "当 $a=\sqrt2$，$b=1$ 时取等号"用全角逗号，均照录；
%   第 19 题解析由 $\dfrac{x+4}{x+2}>2$ 得 "$2-\dfrac{x+4}{x+2}=\dfrac{x}{x+2}<0$"
%     （以 2 减左式，与常见的"左式减 2"次序相反），照录；
%   第 15 题选项 A、B 作"该方程…"、C、D 作"方程组…"，照录；
%   第 20 题题干作"招收代理商"，(1) 却作"招代理商后"，照录。
%
% 行数复核：逐题数 \step 条数，合计 91 条（Q3：1；Q4：0；Q5：0；Q6：3；Q7：2；
% Q8：5；Q9：5；Q10：4；Q11：8；Q12：3；Q13：3；Q14：2；Q15：3；Q16：9；
% Q17：7；Q18：5；Q19：9；Q20：13；Q21：9），与原卷印刷行逐题相等
% （cases 块作一个印刷行计）。第 4、5 题只印【答案】，提 \ans{}，不计入 \step 账；
% 四道选择题答案均在解析末句，括号内留空、不另提 \ans{}。
%
\documentclass[a4paper,11pt]{article}
\usepackage{common}

\begin{document}

\examtitlest[3]{2026--2027 年建青实验高一上期中}{集合与常用逻辑用语、不等式、函数}

\bigsec{一、填空题}

\begin{enumerate}
\setcounter{enumi}{2}

\item 已知 $\log_2 x=3$，则 $x=$\blankline[4.4em].

\ifanswer
\solbegin
\step{$x=2^3=8$.}
\solend
\fi

\item 将 $\sqrt{a^2\sqrt a}$（其中 $a>0$）化为有理数指数幂的形式为\blankline[4.4em].

\ifanswer
\ans{$a^{\frac54}$}
\fi

\item 用反证法证明命题：“已知 $m$、$n\in\R$，若 $m+n\leqslant5$，则 $m\leqslant2$ 或 $n\leqslant3$” 时应假设\blankline[4.4em].

\ifanswer
\ans{$m>2$ 且 $n>3$}
\fi

\item 若 $x\in\R$，则 $\lvert x+1\rvert+\lvert x-1\rvert$ 的最小值为\blankline[4.4em].

\ifanswer
\solbegin
\step{$\lvert x+1\rvert+\lvert x-1\rvert\geqslant\lvert(x+1)-(x-1)\rvert=2$，}
\step{当且仅当 $(x+1)(x-1)\leqslant0$，即 $-1\leqslant x\leqslant1$ 时取等号，}
\step{所以 $\lvert x+1\rvert+\lvert x-1\rvert$ 的最小值为 2.}
\solend
\fi

\item 已知 $x>1$，则 $x+\dfrac1{x-1}$ 的最小值是\blankline[4.4em].

\ifanswer
\solbegin
\step{$x-1>0$，则 $x+\dfrac1{x-1}=(x-1)+\dfrac1{x-1}+1\geqslant2\sqrt{(x-1)\cdot\dfrac1{x-1}}+1=3$，}
\step{当且仅当 $x=2$ 时取等号，故 $x+\dfrac1{x-1}$ 的最小值为 3.}
\solend
\fi

\item 已知集合 $P=\{x\mid x^2+x-6=0\}$，$S=\{x\mid ax+1=0\}$ 且 $S\subset P$，则 $a$ 的值为\blankline[4.4em].

\ifanswer
\solbegin
\step{集合 $P=\{-3,2\}$，集合 $S$ 中至多有一个元素，}
\step{若集合 $S$ 为空集，即 $a=0$ 时，显然满足条件 $S\subset P$，故 $a=0$.}
\step{若集合 $S$ 非空集，即 $a\neq0$，此时 $S=\left\{-\dfrac1a\right\}$，}
\step{若 $-\dfrac1a=-3$，则 $a=\dfrac13$，若 $-\dfrac1a=2$，则 $a=-\dfrac12$；}
\step{故 $a$ 的取值为 $0$、$\dfrac13$、$-\dfrac12$.}
\solend
\fi

\item 已知 $p:\begin{cases}a>4\\b>4\end{cases}$，$q:\begin{cases}a+b>8\\ab>16\end{cases}$，则 $p$ 是 $q$ 的\blankline[4.4em]条件.

\ifanswer
\solbegin
\step{若 $a>4$，$b>4$，则 $a+b>4+4=8$，$ab>4\times4=16$，}
\step{所以 $p$ 成立一定能推出 $q$ 成立，充分性成立，}
\step{取 $a=2$，$b=9$，$a+b=11>8$，$ab=18>16$，满足 $q$；}
\step{但是 $a=2<4$，不满足 $p$，所以由 $q$ 不能推出 $p$，必要性不成立，}
\step{综上所述，$p$ 是 $q$ 的充分不必要条件.}
\solend
\fi

\item 若关于 $x$ 的不等式 $ax^2-ax+1>0$ 对一切实数 $x$ 恒成立，则实数 $a$ 的取值范围是\blankline[4.4em].

\ifanswer
\solbegin
\step{当 $a=0$ 时，不等式即 $1>0$，满足条件.}
\step{当 $a\neq0$ 时，要使不等式 $ax^2+ax+1>0$ 对一切 $x\in\R$ 恒成立，}
\step{需 $\begin{cases}a>0\\\Delta=a^2-4a<0\end{cases}$，解得 $0<a<4$.}
\step{综上，实数 $a$ 的取值范围是 $[0,4)$.}
\solend
\fi

\item 已知集合 $M=\{x\mid (x-a)(x^2-ax+a-1)=0\}$ 各元素之和等于 3，则实数 $a=$\blankline[4.4em].

\ifanswer
\solbegin
\step{方程 $(x-a)(x^2-ax+a-1)=0$，化简得 $(x-a)(x-1)[x-(a-1)]=0$，}
\step{解得 $x_1=a$，$x_2=1$，$x_3=a-1$，}
\step{因为它们是集合 $M$ 中的元素，所以存在互异性.}
\step{当 $x_1\neq x_2\neq x_3$ 时，即 $a+1+a-1=3$，解得 $a=\dfrac32$，此时 $M=\dfrac32,1,\dfrac12$，满足.}
\step{当 $a=1$ 时，$1+0\neq3$，不合题意；}
\step{当 $a=a-1$ 时，不符合题意；}
\step{当 $1=a-1$ 时，$a=2$，此时 $M=\{2,1\}$ 满足题意.}
\step{故 $a$ 的取值为 $2$ 或 $\dfrac32$.}
\solend
\fi

\item 已知 $a>b>1$ 且 $a^2-6a-2m+1=0$，$b^2-6b-2m+1=0$，则实数 $m$ 的取值范围是\blankline[4.4em].

\ifanswer
\solbegin
\step{由题意得 $a$、$b$ 是函数 $f(x)=x^2-6x-2m+1$ 两个大于 1 的不同零点，}
\step{所以 $\begin{cases}f(1)>0\\f(3)<0\end{cases}$，即 $\begin{cases}-4-2m>0\\-8-2m<0\end{cases}$，解得 $-4<m<-2$，}
\step{故 $m$ 的取值范围为 $(-4,-2)$.}
\solend
\fi

\end{enumerate}

\bigsec{二、选择题}

\begin{enumerate}
\setcounter{enumi}{12}

\item “$x<2$” 是 “$x^2<4$” 的\mbox{（\quad）}条件

\keepwithprev
A. 充分不必要\qquad\qquad B. 必要不充分

C. 充要\qquad\qquad\qquad\quad D. 既不充分也不必要

\ifanswer
\solbegin
\step{解 $x^2<4$，得 $-2<x<2$，}
\step{由 $x<2$ 得不出 $-2<x<2$，而由 $-2<x<2$ 能得出 $x<2$，}
\step{所以 “$x<2$” 是 “$x^2<4$” 的必要不充分条件. 故选 $B$.}
\solend
\fi

\item 已知集合 $A=\{x\mid x<-2\}$，$B=\{x\mid x\geqslant a\}$，且 $A\cup B=\R$，则实数 $a$ 的取值范围是\mbox{（\quad）}

\keepwithprev
A. $a\leqslant-2$\qquad\qquad B. $a<-2$\qquad\qquad C. $a\geqslant-2$\qquad\qquad D. $a>-2$

\ifanswer
\solbegin
\step{集合 $A=\{x\mid x<-2\}$，$B=\{x\mid x\geqslant a\}$，且 $A\cup B=\R$，}
\step{则实数 $a$ 的取值范围是 $a\leqslant-2$. 故选 $A$.}
\solend
\fi

\item 下列关于 $x,y$ 的方程组 $\begin{cases}y=kx+1\\y=2kx+3\end{cases}$ 的解的说法中正确的是\mbox{（\quad）}

\keepwithprev
A. 该方程一定有唯一解\qquad\qquad B. 该方程没有解

C. $k\neq0$ 时，方程组有唯一解\qquad\qquad D. $k=0$ 时，方程组有无数组解

\ifanswer
\solbegin
\step{由题意得 $kx+1=2kx+3$，即 $kx=-2$，}
\step{当 $k=0$ 时，$0=-2$ 不成立，方程组无解；}
\step{当 $k\neq0$ 时，$x=-\dfrac2k$，方程组有唯一解. 故选 $C$.}
\solend
\fi

\item 设非空集合 $S=\{x\mid m\leqslant x\leqslant n\}$ 满足：当 $x\in S$ 时，有 $x^2\in S$. 给出如下三个结论：①若 $m=1$，则 $S=\{1\}$；②若 $m=-\dfrac12$，则 $\dfrac14\leqslant n\leqslant1$；③若 $n=\dfrac12$，则 $-\dfrac{\sqrt2}{2}\leqslant m\leqslant0$. 其中正确命题的个数为\mbox{（\quad）}

\keepwithprev
A. $3$\qquad\qquad B. $2$\qquad\qquad C. $1$\qquad\qquad D. $0$

\ifanswer
\solbegin
\step{由定义设非空集合 $S=\{x\mid m\leqslant x\leqslant n\}$ 满足当 $x\in S$ 时，有 $x^2\in S$，}
\step{符合定义的参数 $m$ 的值一定大于等于 1 或小于等于 0，}
\step{惟如此才能保证 $m\in S$ 时，有 $m^2\in S$ 即 $m^2\geqslant m$，}
\step{符合条件的 $n$ 的值一定大于等于 0，小于等于 1，}
\step{惟如此才能保证 $n\in S$ 时，有 $n^2\in S$ 即 $n^2\leqslant n$，对各个命题进行判断：}
\step{对于① $m=1$，$m^2=1\in S$，故必有 $\begin{cases}n^2\leqslant n\\n\geqslant1\end{cases}$，得 $n=1$，$S=\{1\}$，}
\step{② $m=-\dfrac12$，$m^2=\dfrac14\in S$，则 $\begin{cases}n^2\leqslant n\\\dfrac14\leqslant n\end{cases}$，解得 $\dfrac14\leqslant n\leqslant1$；}
\step{对于③，若 $n=\dfrac12$，则 $\begin{cases}\dfrac12\geqslant m\\m^2\geqslant m\\\dfrac12\geqslant m^2\end{cases}$，解得 $-\dfrac{\sqrt2}{2}\leqslant m\leqslant0$，}
\step{所以正确命题有 3 个. 故选 $A$.}
\solend
\fi

\end{enumerate}

\bigsec{三、解答题}

\begin{enumerate}
\setcounter{enumi}{16}

\item 已知 $a$、$b$ 是实数.

（1）求证：$a^2+b^2\geqslant2a-2b-2$，并指出等号成立的条件；

（2）若 $ab=1$，求 $a^2+4b^2$ 的最小值.

\ifanswer
\solbegin
\step{（1）因为 $a^2+b^2-(2a-2b-2)=a^2+b^2-2a+2b+2$\dots\dots2 分}
\step{$=(a-1)^2+(b+1)^2\geqslant0$\dots\dots4 分}
\step{所以 $a^2+b^2\geqslant2a-2b-2$.}
\step{当且仅当 $a=1,b=-1$ 时，不等式中等号成立.\dots\dots7 分}
\step{（2）$a^2+4b^2=a^2+(2b)^2\geqslant2\cdot a\cdot(2b)=4ab=4$，\dots\dots5 分}
\step{当且仅当 $a=2b$，即 $\begin{cases}a=\sqrt2\\b=\dfrac{\sqrt2}{2}\end{cases}$ 或 $\begin{cases}a=-\sqrt2\\b=-\dfrac{\sqrt2}{2}\end{cases}$ 时，不等式中等号成立，}
\step{所以 $a^2+4b^2$ 的最小值为 4.\dots\dots7 分}
\solend
\fi

\item 已知关于 $x$ 的方程 $x^2+(m-3)x+2=0$，其中 $m\in\R$，有两个实数根 $\alpha$、$\beta$.

（1）若 $m=8$，求 $\dfrac1\alpha+\dfrac1\beta$ 的值；

（2）若 $\alpha^2+\beta^2=\alpha\beta+3$，求实数 $m$ 的值.

\ifanswer
\solbegin
\step{（1）若 $m=8$，则关于 $x$ 的方程为 $x^2+5x+2=0$，}
\step{则 $\alpha+\beta=-5$，$\alpha\beta=2$，所以 $\dfrac1\alpha+\dfrac1\beta=\dfrac{\alpha+\beta}{\alpha\beta}=-\dfrac52$；}
\step{（2）由题意得 $\alpha+\beta=-(m-3)=3-m$，$\alpha\beta=2$，}
\step{又 $\alpha^2+\beta^2=\alpha\beta+3$，则 $(\alpha+\beta)^2-2\alpha\beta=\alpha\beta+3$，}
\step{即 $(\alpha+\beta)^2=3\alpha\beta+3$，所以 $(3-m)^2=3\times2+3=9$，解得 $m=0$ 或 $6$.}
\solend
\fi

\item 设集合 $A=\{x\mid\lvert x-a\rvert<2\}$，$B=\left\{x\mid\dfrac{x+4}{x+2}>2\right\}$.

（1）若 $a=1$，试用区间表示集合 $A$、$B$；

（2）若 $A\cup\overline{B}=\R$，求实数 $a$ 的取值范围.

\ifanswer
\solbegin
\step{（1）当 $a=1$ 时，由 $\lvert x-1\rvert<2$ 得 $-2<x-1<2$，解得 $-1<x<3$，}
\step{所以 $A=(-1,3)$.}
\step{由 $\dfrac{x+4}{x+2}>2$ 得 $2-\dfrac{x+4}{x+2}=\dfrac{x}{x+2}<0$，则有 $x(x+2)<0$，解得 $-2<x<0$，}
\step{所以 $B=(-2,0)$；}
\step{（2）由 $\lvert x-a\rvert<2$ 得 $-2<x-a<2$，解得 $a-2<x<a+2$，}
\step{所以 $A=(a-2,a+2)$，}
\step{由（1）得 $B=(-2,0)$，则 $\overline{B}=\{x\mid x\geqslant0$ 或 $x\leqslant-2\}$，}
\step{由于 $A\cup\overline{B}=\R$，所以 $\begin{cases}a+2\geqslant0\\a-2\leqslant-2\end{cases}$，解得 $-2\leqslant a\leqslant0$，}
\step{所以实数 $a$ 的取值范围是 $[-2,0]$.}
\solend
\fi

\item 某服装公司生产的衬衫每件定价为 80 元，在某城市年销售 8 万件. 现该公司计划在该市招收代理商来销售衬衫，以降低管理和营销成本. 已知代理商收取的代理费为总销售金额的 $r\%$（其中 $0<r\leqslant10$）.

（1）若单件衬衫的利润保持不变，求招代理商后每件衬衫的定价为多少元（用 $r$ 表示）；

（2）已知将每件衬衫的价格提高到 $\dfrac{80}{1-r\%}$ 元后，每年的销售量会减少 $0.62r$ 万件. 若代理商每年收取的代理费不少于 16 万元，求 $r$ 的取值范围.

\ifanswer
\solbegin
\step{（1）设招代理商后每件衬衫的定价为 $x$ 元，代理商收取总销售金额的 $r\%$，}
\step{因此公司每件衬衫的实际净收入为 $x(1-r\%)$ 元，}
\step{单件利润保持不变，等价于净收入仍为 80 元，}
\step{于是 $x(1-r\%)=80$，解得 $x=\dfrac{80}{1-r\%}$，因为 $0<r\leqslant10$，有 $1-r\%\geqslant0.9>0$，}
\step{分母恒不为零，该定价有意义，所以招代理商后每件衬衫的定价为 $\dfrac{80}{1-r\%}$ 元；}
\step{（2）价格提高到 $\dfrac{80}{1-r\%}$ 元，年销售量减少 $0.62r$ 万件，}
\step{故年销售量为 $8-0.62r$ 万件，$8-0.62r>0$，}
\step{由题意得 $\dfrac{80}{1-r\%}\cdot(8-0.62r)\cdot r\%\geqslant16$，}
\step{即 $\dfrac{80}{1-\dfrac r{100}}\cdot(8-0.62r)\cdot\dfrac r{100}=\dfrac{80r(8-0.62r)}{100-r}\geqslant16$，}
\step{由 $0<r\leqslant10$ 得 $100-r>0$，整理得 $80r(8-0.62r)\geqslant16(100-r)$，}
\step{整理得 $31r^2-410r+1000\leqslant0$，解得 $\dfrac{100}{31}\leqslant r\leqslant10$，}
\step{当 $\dfrac{100}{31}\leqslant r\leqslant10$ 时，$8-0.62r>0$ 且 $1-r\%>0$，}
\step{均符合实际经营条件，所以 $r$ 的取值范围是 $\left[\dfrac{100}{31},10\right]$.}
\solend
\fi

\item 定义 $\min\{a_1,a_2,\cdots,a_n\}$ 为 $n$ 个实数 $a_1,a_2,\cdots,a_n$ 中的最小数，$\max\{a_1,a_2,\cdots,a_n\}$ 为 $n$ 个实数 $a_1,a_2,\cdots,a_n$ 中的最大数.

（1）设 $a,b$ 都是正实数，且 $a+b=1$，求 $\max\left\{ab,\dfrac14\right\}$；

（2）解不等式：$\min\{x+1,x^2+3,\lvert x-1\rvert\}>2x-3$；

（3）设 $a,b$ 都是正实数，求 $\max\left\{a+\dfrac1b,\dfrac2a+b\right\}$ 的最小值.

\ifanswer
\solbegin
\step{（1）$ab\leqslant\dfrac{(a+b)^2}{4}=\dfrac14$，所以 $\max\left\{ab,\dfrac14\right\}=\dfrac14$.}
\step{（2）因为 $x^2+3>x+1$，$x^2+3>\lvert x-1\rvert$，}
\step{所以 $\min\{x+1,x^2+3,\lvert x-1\rvert\}=\min\{x+1,\lvert x-1\rvert\}$.}
\step{当 $x\geqslant0$ 时，$\lvert x-1\rvert\leqslant x+1$，所以 $\lvert x-1\rvert>2x-3$，解得 $0\leqslant x<2$；}
\step{当 $x<0$ 时，$\lvert x-1\rvert>x+1$，所以 $x+1>2x-3$，解得 $x<0$.}
\step{综上，不等式的解集为 $(-\infty,2)$.}
\step{（3）设 $M=\max\left\{a+\dfrac1b,\dfrac2a+b\right\}$，则 $M\geqslant a+\dfrac1b$，$M\geqslant\dfrac2a+b$，}
\step{所以 $M^2\geqslant\left(a+\dfrac1b\right)\left(\dfrac2a+b\right)=3+ab+\dfrac2{ab}\geqslant3+2\sqrt2$，}
\step{所以 $M$ 的最小值为 $\sqrt2+1$，当 $a=\sqrt2$，$b=1$ 时取等号.}
\solend
\fi

\end{enumerate}

\end{document}
"
% 紧凑版：xelatex "\def\iscompact{1}% 高一126_建青实验_期中.tex
%   —— 高一上数学"2026-2027 年建青实验高一上期中"（试卷版/答案版）
% 试卷版：xelatex 高一126_建青实验_期中.tex
% 答案版：xelatex "\def\isanswer{1}\input{高一126_建青实验_期中.tex}"
% 紧凑版：xelatex "\def\iscompact{1}\input{高一126_建青实验_期中.tex}"
%
% 原始材料：微信公众号"魔都数学加油站"文章，作者破晓，2026-11-11 18:35 发布；
% 连载"27学年高一126"，原文标题"27学年高一126——2026-2027年上海市建青实验高一上期中"；
% 原文链接：https://mp.weixin.qq.com/s/HpxpBJu4d_wOnUCIbZSQiA。
% 正文为 8 张 PNG 页（01、05、08 为 2560x3620，02、03、04、06、07 为 4762x6735），
% 题目下印有【解析】或【答案】，为讲评版。题目、解析均据原图转录；
% 卷面粉红色水印"魔都数学加油站"及公众号页脚未录入。全卷无插图。
%
% 卷首标题：原卷印作"2026-2027 年建青实验高一上期中"（公众号标题中的"上海市"
% 卷面标题行没有，本稿照卷面），年份区间按全稿惯例排 en dash；原卷未印副标题，
% 本卷为期中卷、跨集合与常用逻辑用语／不等式／函数三章，按题目内容标为
% "集合与常用逻辑用语、不等式、函数"（与 125 同为期中卷，口径相同）：
% 集合与逻辑见 Q5、Q8、Q9、Q11、Q13、Q14、Q16、Q19，不等式见 Q6、Q7、Q10、
% Q17、Q20、Q21，函数（含指数对数、二次函数零点）见 Q3、Q4、Q12、Q15、Q18。
%
% 与原卷的排版差异：
%   ① 原文从第 3 题起，填空题编号为 3—12、选择题为 13—16、解答题为 17—21，本稿保留原题号；
%   ② 原卷每题下直接印【解析】／【答案】，本稿按全稿惯例置于答案版；第 4、5 题只印
%      【答案】，另提 \ans{}；第 13—16 题答案均在【解析】末句"故选 …"，题干括号内留空、
%      不另提 \ans{}；第 3 题印一行【解析】，按 \solbegin+\step 排；第 17 题解析保留
%      原卷评分点（照 高一54／64 成例排 \dots\dots N 分）；
%   ③ 原卷未印副标题，本稿按"知识模块"口径补出（见上"卷首标题"）；
%   ④ 原卷填空横线按全稿惯例排为 \blankline[4.4em]；选择题括号排 \mbox{（\quad）}、
%      选项段前按惯例加 \keepwithprev；第 14 题原卷括号另起一行，本稿并入题干行末；
%      解答题不留作答区（原卷为讲评版，与 109—125 同）；
%   ⑤ 原卷数集排斜体 $R$，按全稿惯例改 \R；绝对值排 \lvert…\rvert，集合描述竖线排 \mid；
%      补集记号照原卷上划线排 \overline{B}；不等号按全稿惯例排 \leqslant／\geqslant；
%      子题标号原卷一律用全角括号（1）（2）（3），本稿照录（\step 行内亦同，照 高一121 成例）。
%
% 判别式字形：第 10 题解析 $\Delta=a^2-4a<0$ 原卷排 $\Delta$（等腰、高），
%   照录 \Delta（与 高一124／125 同题排 $\triangle$（等边、扁）不同，逐处照本卷字形）。
%
% 子集符号：第 8 题题干与解析两处 $S\subset P$ 均印单线、无下划线，照录 \subset；
%   全卷未出现带下划线的 \subseteq。
%
% 原卷存疑照录：
%   第 10 题题干作 $ax^2-ax+1>0$，解析却作 $ax^2+ax+1>0$（一次项符号不一致），照录、未统一；
%   第 11 题解析"此时 $M=\dfrac32,1,\dfrac12$，满足"未用集合花括号（题干 $M$ 为集合），照录；
%   第 11 题解析"因为它们是集合 $M$ 中的元素，所以存在互异性"（"存在互异性"疑为
%     "满足互异性"之误），又"当 $x_1\neq x_2\neq x_3$ 时"连写两个不等号，均照录；
%   第 16 题解析两处"惟如此才能保证"（用"惟"字）及"符合定义的参数 $m$ 的值一定大于等于 1
%     或小于等于 0"的措辞，照录；
%   第 17 题 (1) 解析"当且仅当 $a=1,b=-1$ 时"用半角逗号、第 21 题 (3) 解析
%     "当 $a=\sqrt2$，$b=1$ 时取等号"用全角逗号，均照录；
%   第 19 题解析由 $\dfrac{x+4}{x+2}>2$ 得 "$2-\dfrac{x+4}{x+2}=\dfrac{x}{x+2}<0$"
%     （以 2 减左式，与常见的"左式减 2"次序相反），照录；
%   第 15 题选项 A、B 作"该方程…"、C、D 作"方程组…"，照录；
%   第 20 题题干作"招收代理商"，(1) 却作"招代理商后"，照录。
%
% 行数复核：逐题数 \step 条数，合计 91 条（Q3：1；Q4：0；Q5：0；Q6：3；Q7：2；
% Q8：5；Q9：5；Q10：4；Q11：8；Q12：3；Q13：3；Q14：2；Q15：3；Q16：9；
% Q17：7；Q18：5；Q19：9；Q20：13；Q21：9），与原卷印刷行逐题相等
% （cases 块作一个印刷行计）。第 4、5 题只印【答案】，提 \ans{}，不计入 \step 账；
% 四道选择题答案均在解析末句，括号内留空、不另提 \ans{}。
%
\documentclass[a4paper,11pt]{article}
\usepackage{common}

\begin{document}

\examtitlest[3]{2026--2027 年建青实验高一上期中}{集合与常用逻辑用语、不等式、函数}

\bigsec{一、填空题}

\begin{enumerate}
\setcounter{enumi}{2}

\item 已知 $\log_2 x=3$，则 $x=$\blankline[4.4em].

\ifanswer
\solbegin
\step{$x=2^3=8$.}
\solend
\fi

\item 将 $\sqrt{a^2\sqrt a}$（其中 $a>0$）化为有理数指数幂的形式为\blankline[4.4em].

\ifanswer
\ans{$a^{\frac54}$}
\fi

\item 用反证法证明命题：“已知 $m$、$n\in\R$，若 $m+n\leqslant5$，则 $m\leqslant2$ 或 $n\leqslant3$” 时应假设\blankline[4.4em].

\ifanswer
\ans{$m>2$ 且 $n>3$}
\fi

\item 若 $x\in\R$，则 $\lvert x+1\rvert+\lvert x-1\rvert$ 的最小值为\blankline[4.4em].

\ifanswer
\solbegin
\step{$\lvert x+1\rvert+\lvert x-1\rvert\geqslant\lvert(x+1)-(x-1)\rvert=2$，}
\step{当且仅当 $(x+1)(x-1)\leqslant0$，即 $-1\leqslant x\leqslant1$ 时取等号，}
\step{所以 $\lvert x+1\rvert+\lvert x-1\rvert$ 的最小值为 2.}
\solend
\fi

\item 已知 $x>1$，则 $x+\dfrac1{x-1}$ 的最小值是\blankline[4.4em].

\ifanswer
\solbegin
\step{$x-1>0$，则 $x+\dfrac1{x-1}=(x-1)+\dfrac1{x-1}+1\geqslant2\sqrt{(x-1)\cdot\dfrac1{x-1}}+1=3$，}
\step{当且仅当 $x=2$ 时取等号，故 $x+\dfrac1{x-1}$ 的最小值为 3.}
\solend
\fi

\item 已知集合 $P=\{x\mid x^2+x-6=0\}$，$S=\{x\mid ax+1=0\}$ 且 $S\subset P$，则 $a$ 的值为\blankline[4.4em].

\ifanswer
\solbegin
\step{集合 $P=\{-3,2\}$，集合 $S$ 中至多有一个元素，}
\step{若集合 $S$ 为空集，即 $a=0$ 时，显然满足条件 $S\subset P$，故 $a=0$.}
\step{若集合 $S$ 非空集，即 $a\neq0$，此时 $S=\left\{-\dfrac1a\right\}$，}
\step{若 $-\dfrac1a=-3$，则 $a=\dfrac13$，若 $-\dfrac1a=2$，则 $a=-\dfrac12$；}
\step{故 $a$ 的取值为 $0$、$\dfrac13$、$-\dfrac12$.}
\solend
\fi

\item 已知 $p:\begin{cases}a>4\\b>4\end{cases}$，$q:\begin{cases}a+b>8\\ab>16\end{cases}$，则 $p$ 是 $q$ 的\blankline[4.4em]条件.

\ifanswer
\solbegin
\step{若 $a>4$，$b>4$，则 $a+b>4+4=8$，$ab>4\times4=16$，}
\step{所以 $p$ 成立一定能推出 $q$ 成立，充分性成立，}
\step{取 $a=2$，$b=9$，$a+b=11>8$，$ab=18>16$，满足 $q$；}
\step{但是 $a=2<4$，不满足 $p$，所以由 $q$ 不能推出 $p$，必要性不成立，}
\step{综上所述，$p$ 是 $q$ 的充分不必要条件.}
\solend
\fi

\item 若关于 $x$ 的不等式 $ax^2-ax+1>0$ 对一切实数 $x$ 恒成立，则实数 $a$ 的取值范围是\blankline[4.4em].

\ifanswer
\solbegin
\step{当 $a=0$ 时，不等式即 $1>0$，满足条件.}
\step{当 $a\neq0$ 时，要使不等式 $ax^2+ax+1>0$ 对一切 $x\in\R$ 恒成立，}
\step{需 $\begin{cases}a>0\\\Delta=a^2-4a<0\end{cases}$，解得 $0<a<4$.}
\step{综上，实数 $a$ 的取值范围是 $[0,4)$.}
\solend
\fi

\item 已知集合 $M=\{x\mid (x-a)(x^2-ax+a-1)=0\}$ 各元素之和等于 3，则实数 $a=$\blankline[4.4em].

\ifanswer
\solbegin
\step{方程 $(x-a)(x^2-ax+a-1)=0$，化简得 $(x-a)(x-1)[x-(a-1)]=0$，}
\step{解得 $x_1=a$，$x_2=1$，$x_3=a-1$，}
\step{因为它们是集合 $M$ 中的元素，所以存在互异性.}
\step{当 $x_1\neq x_2\neq x_3$ 时，即 $a+1+a-1=3$，解得 $a=\dfrac32$，此时 $M=\dfrac32,1,\dfrac12$，满足.}
\step{当 $a=1$ 时，$1+0\neq3$，不合题意；}
\step{当 $a=a-1$ 时，不符合题意；}
\step{当 $1=a-1$ 时，$a=2$，此时 $M=\{2,1\}$ 满足题意.}
\step{故 $a$ 的取值为 $2$ 或 $\dfrac32$.}
\solend
\fi

\item 已知 $a>b>1$ 且 $a^2-6a-2m+1=0$，$b^2-6b-2m+1=0$，则实数 $m$ 的取值范围是\blankline[4.4em].

\ifanswer
\solbegin
\step{由题意得 $a$、$b$ 是函数 $f(x)=x^2-6x-2m+1$ 两个大于 1 的不同零点，}
\step{所以 $\begin{cases}f(1)>0\\f(3)<0\end{cases}$，即 $\begin{cases}-4-2m>0\\-8-2m<0\end{cases}$，解得 $-4<m<-2$，}
\step{故 $m$ 的取值范围为 $(-4,-2)$.}
\solend
\fi

\end{enumerate}

\bigsec{二、选择题}

\begin{enumerate}
\setcounter{enumi}{12}

\item “$x<2$” 是 “$x^2<4$” 的\mbox{（\quad）}条件

\keepwithprev
A. 充分不必要\qquad\qquad B. 必要不充分

C. 充要\qquad\qquad\qquad\quad D. 既不充分也不必要

\ifanswer
\solbegin
\step{解 $x^2<4$，得 $-2<x<2$，}
\step{由 $x<2$ 得不出 $-2<x<2$，而由 $-2<x<2$ 能得出 $x<2$，}
\step{所以 “$x<2$” 是 “$x^2<4$” 的必要不充分条件. 故选 $B$.}
\solend
\fi

\item 已知集合 $A=\{x\mid x<-2\}$，$B=\{x\mid x\geqslant a\}$，且 $A\cup B=\R$，则实数 $a$ 的取值范围是\mbox{（\quad）}

\keepwithprev
A. $a\leqslant-2$\qquad\qquad B. $a<-2$\qquad\qquad C. $a\geqslant-2$\qquad\qquad D. $a>-2$

\ifanswer
\solbegin
\step{集合 $A=\{x\mid x<-2\}$，$B=\{x\mid x\geqslant a\}$，且 $A\cup B=\R$，}
\step{则实数 $a$ 的取值范围是 $a\leqslant-2$. 故选 $A$.}
\solend
\fi

\item 下列关于 $x,y$ 的方程组 $\begin{cases}y=kx+1\\y=2kx+3\end{cases}$ 的解的说法中正确的是\mbox{（\quad）}

\keepwithprev
A. 该方程一定有唯一解\qquad\qquad B. 该方程没有解

C. $k\neq0$ 时，方程组有唯一解\qquad\qquad D. $k=0$ 时，方程组有无数组解

\ifanswer
\solbegin
\step{由题意得 $kx+1=2kx+3$，即 $kx=-2$，}
\step{当 $k=0$ 时，$0=-2$ 不成立，方程组无解；}
\step{当 $k\neq0$ 时，$x=-\dfrac2k$，方程组有唯一解. 故选 $C$.}
\solend
\fi

\item 设非空集合 $S=\{x\mid m\leqslant x\leqslant n\}$ 满足：当 $x\in S$ 时，有 $x^2\in S$. 给出如下三个结论：①若 $m=1$，则 $S=\{1\}$；②若 $m=-\dfrac12$，则 $\dfrac14\leqslant n\leqslant1$；③若 $n=\dfrac12$，则 $-\dfrac{\sqrt2}{2}\leqslant m\leqslant0$. 其中正确命题的个数为\mbox{（\quad）}

\keepwithprev
A. $3$\qquad\qquad B. $2$\qquad\qquad C. $1$\qquad\qquad D. $0$

\ifanswer
\solbegin
\step{由定义设非空集合 $S=\{x\mid m\leqslant x\leqslant n\}$ 满足当 $x\in S$ 时，有 $x^2\in S$，}
\step{符合定义的参数 $m$ 的值一定大于等于 1 或小于等于 0，}
\step{惟如此才能保证 $m\in S$ 时，有 $m^2\in S$ 即 $m^2\geqslant m$，}
\step{符合条件的 $n$ 的值一定大于等于 0，小于等于 1，}
\step{惟如此才能保证 $n\in S$ 时，有 $n^2\in S$ 即 $n^2\leqslant n$，对各个命题进行判断：}
\step{对于① $m=1$，$m^2=1\in S$，故必有 $\begin{cases}n^2\leqslant n\\n\geqslant1\end{cases}$，得 $n=1$，$S=\{1\}$，}
\step{② $m=-\dfrac12$，$m^2=\dfrac14\in S$，则 $\begin{cases}n^2\leqslant n\\\dfrac14\leqslant n\end{cases}$，解得 $\dfrac14\leqslant n\leqslant1$；}
\step{对于③，若 $n=\dfrac12$，则 $\begin{cases}\dfrac12\geqslant m\\m^2\geqslant m\\\dfrac12\geqslant m^2\end{cases}$，解得 $-\dfrac{\sqrt2}{2}\leqslant m\leqslant0$，}
\step{所以正确命题有 3 个. 故选 $A$.}
\solend
\fi

\end{enumerate}

\bigsec{三、解答题}

\begin{enumerate}
\setcounter{enumi}{16}

\item 已知 $a$、$b$ 是实数.

（1）求证：$a^2+b^2\geqslant2a-2b-2$，并指出等号成立的条件；

（2）若 $ab=1$，求 $a^2+4b^2$ 的最小值.

\ifanswer
\solbegin
\step{（1）因为 $a^2+b^2-(2a-2b-2)=a^2+b^2-2a+2b+2$\dots\dots2 分}
\step{$=(a-1)^2+(b+1)^2\geqslant0$\dots\dots4 分}
\step{所以 $a^2+b^2\geqslant2a-2b-2$.}
\step{当且仅当 $a=1,b=-1$ 时，不等式中等号成立.\dots\dots7 分}
\step{（2）$a^2+4b^2=a^2+(2b)^2\geqslant2\cdot a\cdot(2b)=4ab=4$，\dots\dots5 分}
\step{当且仅当 $a=2b$，即 $\begin{cases}a=\sqrt2\\b=\dfrac{\sqrt2}{2}\end{cases}$ 或 $\begin{cases}a=-\sqrt2\\b=-\dfrac{\sqrt2}{2}\end{cases}$ 时，不等式中等号成立，}
\step{所以 $a^2+4b^2$ 的最小值为 4.\dots\dots7 分}
\solend
\fi

\item 已知关于 $x$ 的方程 $x^2+(m-3)x+2=0$，其中 $m\in\R$，有两个实数根 $\alpha$、$\beta$.

（1）若 $m=8$，求 $\dfrac1\alpha+\dfrac1\beta$ 的值；

（2）若 $\alpha^2+\beta^2=\alpha\beta+3$，求实数 $m$ 的值.

\ifanswer
\solbegin
\step{（1）若 $m=8$，则关于 $x$ 的方程为 $x^2+5x+2=0$，}
\step{则 $\alpha+\beta=-5$，$\alpha\beta=2$，所以 $\dfrac1\alpha+\dfrac1\beta=\dfrac{\alpha+\beta}{\alpha\beta}=-\dfrac52$；}
\step{（2）由题意得 $\alpha+\beta=-(m-3)=3-m$，$\alpha\beta=2$，}
\step{又 $\alpha^2+\beta^2=\alpha\beta+3$，则 $(\alpha+\beta)^2-2\alpha\beta=\alpha\beta+3$，}
\step{即 $(\alpha+\beta)^2=3\alpha\beta+3$，所以 $(3-m)^2=3\times2+3=9$，解得 $m=0$ 或 $6$.}
\solend
\fi

\item 设集合 $A=\{x\mid\lvert x-a\rvert<2\}$，$B=\left\{x\mid\dfrac{x+4}{x+2}>2\right\}$.

（1）若 $a=1$，试用区间表示集合 $A$、$B$；

（2）若 $A\cup\overline{B}=\R$，求实数 $a$ 的取值范围.

\ifanswer
\solbegin
\step{（1）当 $a=1$ 时，由 $\lvert x-1\rvert<2$ 得 $-2<x-1<2$，解得 $-1<x<3$，}
\step{所以 $A=(-1,3)$.}
\step{由 $\dfrac{x+4}{x+2}>2$ 得 $2-\dfrac{x+4}{x+2}=\dfrac{x}{x+2}<0$，则有 $x(x+2)<0$，解得 $-2<x<0$，}
\step{所以 $B=(-2,0)$；}
\step{（2）由 $\lvert x-a\rvert<2$ 得 $-2<x-a<2$，解得 $a-2<x<a+2$，}
\step{所以 $A=(a-2,a+2)$，}
\step{由（1）得 $B=(-2,0)$，则 $\overline{B}=\{x\mid x\geqslant0$ 或 $x\leqslant-2\}$，}
\step{由于 $A\cup\overline{B}=\R$，所以 $\begin{cases}a+2\geqslant0\\a-2\leqslant-2\end{cases}$，解得 $-2\leqslant a\leqslant0$，}
\step{所以实数 $a$ 的取值范围是 $[-2,0]$.}
\solend
\fi

\item 某服装公司生产的衬衫每件定价为 80 元，在某城市年销售 8 万件. 现该公司计划在该市招收代理商来销售衬衫，以降低管理和营销成本. 已知代理商收取的代理费为总销售金额的 $r\%$（其中 $0<r\leqslant10$）.

（1）若单件衬衫的利润保持不变，求招代理商后每件衬衫的定价为多少元（用 $r$ 表示）；

（2）已知将每件衬衫的价格提高到 $\dfrac{80}{1-r\%}$ 元后，每年的销售量会减少 $0.62r$ 万件. 若代理商每年收取的代理费不少于 16 万元，求 $r$ 的取值范围.

\ifanswer
\solbegin
\step{（1）设招代理商后每件衬衫的定价为 $x$ 元，代理商收取总销售金额的 $r\%$，}
\step{因此公司每件衬衫的实际净收入为 $x(1-r\%)$ 元，}
\step{单件利润保持不变，等价于净收入仍为 80 元，}
\step{于是 $x(1-r\%)=80$，解得 $x=\dfrac{80}{1-r\%}$，因为 $0<r\leqslant10$，有 $1-r\%\geqslant0.9>0$，}
\step{分母恒不为零，该定价有意义，所以招代理商后每件衬衫的定价为 $\dfrac{80}{1-r\%}$ 元；}
\step{（2）价格提高到 $\dfrac{80}{1-r\%}$ 元，年销售量减少 $0.62r$ 万件，}
\step{故年销售量为 $8-0.62r$ 万件，$8-0.62r>0$，}
\step{由题意得 $\dfrac{80}{1-r\%}\cdot(8-0.62r)\cdot r\%\geqslant16$，}
\step{即 $\dfrac{80}{1-\dfrac r{100}}\cdot(8-0.62r)\cdot\dfrac r{100}=\dfrac{80r(8-0.62r)}{100-r}\geqslant16$，}
\step{由 $0<r\leqslant10$ 得 $100-r>0$，整理得 $80r(8-0.62r)\geqslant16(100-r)$，}
\step{整理得 $31r^2-410r+1000\leqslant0$，解得 $\dfrac{100}{31}\leqslant r\leqslant10$，}
\step{当 $\dfrac{100}{31}\leqslant r\leqslant10$ 时，$8-0.62r>0$ 且 $1-r\%>0$，}
\step{均符合实际经营条件，所以 $r$ 的取值范围是 $\left[\dfrac{100}{31},10\right]$.}
\solend
\fi

\item 定义 $\min\{a_1,a_2,\cdots,a_n\}$ 为 $n$ 个实数 $a_1,a_2,\cdots,a_n$ 中的最小数，$\max\{a_1,a_2,\cdots,a_n\}$ 为 $n$ 个实数 $a_1,a_2,\cdots,a_n$ 中的最大数.

（1）设 $a,b$ 都是正实数，且 $a+b=1$，求 $\max\left\{ab,\dfrac14\right\}$；

（2）解不等式：$\min\{x+1,x^2+3,\lvert x-1\rvert\}>2x-3$；

（3）设 $a,b$ 都是正实数，求 $\max\left\{a+\dfrac1b,\dfrac2a+b\right\}$ 的最小值.

\ifanswer
\solbegin
\step{（1）$ab\leqslant\dfrac{(a+b)^2}{4}=\dfrac14$，所以 $\max\left\{ab,\dfrac14\right\}=\dfrac14$.}
\step{（2）因为 $x^2+3>x+1$，$x^2+3>\lvert x-1\rvert$，}
\step{所以 $\min\{x+1,x^2+3,\lvert x-1\rvert\}=\min\{x+1,\lvert x-1\rvert\}$.}
\step{当 $x\geqslant0$ 时，$\lvert x-1\rvert\leqslant x+1$，所以 $\lvert x-1\rvert>2x-3$，解得 $0\leqslant x<2$；}
\step{当 $x<0$ 时，$\lvert x-1\rvert>x+1$，所以 $x+1>2x-3$，解得 $x<0$.}
\step{综上，不等式的解集为 $(-\infty,2)$.}
\step{（3）设 $M=\max\left\{a+\dfrac1b,\dfrac2a+b\right\}$，则 $M\geqslant a+\dfrac1b$，$M\geqslant\dfrac2a+b$，}
\step{所以 $M^2\geqslant\left(a+\dfrac1b\right)\left(\dfrac2a+b\right)=3+ab+\dfrac2{ab}\geqslant3+2\sqrt2$，}
\step{所以 $M$ 的最小值为 $\sqrt2+1$，当 $a=\sqrt2$，$b=1$ 时取等号.}
\solend
\fi

\end{enumerate}

\end{document}
"
%
% 原始材料：微信公众号"魔都数学加油站"文章，作者破晓，2026-11-11 18:35 发布；
% 连载"27学年高一126"，原文标题"27学年高一126——2026-2027年上海市建青实验高一上期中"；
% 原文链接：https://mp.weixin.qq.com/s/HpxpBJu4d_wOnUCIbZSQiA。
% 正文为 8 张 PNG 页（01、05、08 为 2560x3620，02、03、04、06、07 为 4762x6735），
% 题目下印有【解析】或【答案】，为讲评版。题目、解析均据原图转录；
% 卷面粉红色水印"魔都数学加油站"及公众号页脚未录入。全卷无插图。
%
% 卷首标题：原卷印作"2026-2027 年建青实验高一上期中"（公众号标题中的"上海市"
% 卷面标题行没有，本稿照卷面），年份区间按全稿惯例排 en dash；原卷未印副标题，
% 本卷为期中卷、跨集合与常用逻辑用语／不等式／函数三章，按题目内容标为
% "集合与常用逻辑用语、不等式、函数"（与 125 同为期中卷，口径相同）：
% 集合与逻辑见 Q5、Q8、Q9、Q11、Q13、Q14、Q16、Q19，不等式见 Q6、Q7、Q10、
% Q17、Q20、Q21，函数（含指数对数、二次函数零点）见 Q3、Q4、Q12、Q15、Q18。
%
% 与原卷的排版差异：
%   ① 原文从第 3 题起，填空题编号为 3—12、选择题为 13—16、解答题为 17—21，本稿保留原题号；
%   ② 原卷每题下直接印【解析】／【答案】，本稿按全稿惯例置于答案版；第 4、5 题只印
%      【答案】，另提 \ans{}；第 13—16 题答案均在【解析】末句"故选 …"，题干括号内留空、
%      不另提 \ans{}；第 3 题印一行【解析】，按 \solbegin+\step 排；第 17 题解析保留
%      原卷评分点（照 高一54／64 成例排 \dots\dots N 分）；
%   ③ 原卷未印副标题，本稿按"知识模块"口径补出（见上"卷首标题"）；
%   ④ 原卷填空横线按全稿惯例排为 \blankline[4.4em]；选择题括号排 \mbox{（\quad）}、
%      选项段前按惯例加 \keepwithprev；第 14 题原卷括号另起一行，本稿并入题干行末；
%      解答题不留作答区（原卷为讲评版，与 109—125 同）；
%   ⑤ 原卷数集排斜体 $R$，按全稿惯例改 \R；绝对值排 \lvert…\rvert，集合描述竖线排 \mid；
%      补集记号照原卷上划线排 \overline{B}；不等号按全稿惯例排 \leqslant／\geqslant；
%      子题标号原卷一律用全角括号（1）（2）（3），本稿照录（\step 行内亦同，照 高一121 成例）。
%
% 判别式字形：第 10 题解析 $\Delta=a^2-4a<0$ 原卷排 $\Delta$（等腰、高），
%   照录 \Delta（与 高一124／125 同题排 $\triangle$（等边、扁）不同，逐处照本卷字形）。
%
% 子集符号：第 8 题题干与解析两处 $S\subset P$ 均印单线、无下划线，照录 \subset；
%   全卷未出现带下划线的 \subseteq。
%
% 原卷存疑照录：
%   第 10 题题干作 $ax^2-ax+1>0$，解析却作 $ax^2+ax+1>0$（一次项符号不一致），照录、未统一；
%   第 11 题解析"此时 $M=\dfrac32,1,\dfrac12$，满足"未用集合花括号（题干 $M$ 为集合），照录；
%   第 11 题解析"因为它们是集合 $M$ 中的元素，所以存在互异性"（"存在互异性"疑为
%     "满足互异性"之误），又"当 $x_1\neq x_2\neq x_3$ 时"连写两个不等号，均照录；
%   第 16 题解析两处"惟如此才能保证"（用"惟"字）及"符合定义的参数 $m$ 的值一定大于等于 1
%     或小于等于 0"的措辞，照录；
%   第 17 题 (1) 解析"当且仅当 $a=1,b=-1$ 时"用半角逗号、第 21 题 (3) 解析
%     "当 $a=\sqrt2$，$b=1$ 时取等号"用全角逗号，均照录；
%   第 19 题解析由 $\dfrac{x+4}{x+2}>2$ 得 "$2-\dfrac{x+4}{x+2}=\dfrac{x}{x+2}<0$"
%     （以 2 减左式，与常见的"左式减 2"次序相反），照录；
%   第 15 题选项 A、B 作"该方程…"、C、D 作"方程组…"，照录；
%   第 20 题题干作"招收代理商"，(1) 却作"招代理商后"，照录。
%
% 行数复核：逐题数 \step 条数，合计 91 条（Q3：1；Q4：0；Q5：0；Q6：3；Q7：2；
% Q8：5；Q9：5；Q10：4；Q11：8；Q12：3；Q13：3；Q14：2；Q15：3；Q16：9；
% Q17：7；Q18：5；Q19：9；Q20：13；Q21：9），与原卷印刷行逐题相等
% （cases 块作一个印刷行计）。第 4、5 题只印【答案】，提 \ans{}，不计入 \step 账；
% 四道选择题答案均在解析末句，括号内留空、不另提 \ans{}。
%
\documentclass[a4paper,11pt]{article}
\usepackage{common}

\begin{document}

\examtitlest[3]{2026--2027 年建青实验高一上期中}{集合与常用逻辑用语、不等式、函数}

\bigsec{一、填空题}

\begin{enumerate}
\setcounter{enumi}{2}

\item 已知 $\log_2 x=3$，则 $x=$\blankline[4.4em].

\ifanswer
\solbegin
\step{$x=2^3=8$.}
\solend
\fi

\item 将 $\sqrt{a^2\sqrt a}$（其中 $a>0$）化为有理数指数幂的形式为\blankline[4.4em].

\ifanswer
\ans{$a^{\frac54}$}
\fi

\item 用反证法证明命题：“已知 $m$、$n\in\R$，若 $m+n\leqslant5$，则 $m\leqslant2$ 或 $n\leqslant3$” 时应假设\blankline[4.4em].

\ifanswer
\ans{$m>2$ 且 $n>3$}
\fi

\item 若 $x\in\R$，则 $\lvert x+1\rvert+\lvert x-1\rvert$ 的最小值为\blankline[4.4em].

\ifanswer
\solbegin
\step{$\lvert x+1\rvert+\lvert x-1\rvert\geqslant\lvert(x+1)-(x-1)\rvert=2$，}
\step{当且仅当 $(x+1)(x-1)\leqslant0$，即 $-1\leqslant x\leqslant1$ 时取等号，}
\step{所以 $\lvert x+1\rvert+\lvert x-1\rvert$ 的最小值为 2.}
\solend
\fi

\item 已知 $x>1$，则 $x+\dfrac1{x-1}$ 的最小值是\blankline[4.4em].

\ifanswer
\solbegin
\step{$x-1>0$，则 $x+\dfrac1{x-1}=(x-1)+\dfrac1{x-1}+1\geqslant2\sqrt{(x-1)\cdot\dfrac1{x-1}}+1=3$，}
\step{当且仅当 $x=2$ 时取等号，故 $x+\dfrac1{x-1}$ 的最小值为 3.}
\solend
\fi

\item 已知集合 $P=\{x\mid x^2+x-6=0\}$，$S=\{x\mid ax+1=0\}$ 且 $S\subset P$，则 $a$ 的值为\blankline[4.4em].

\ifanswer
\solbegin
\step{集合 $P=\{-3,2\}$，集合 $S$ 中至多有一个元素，}
\step{若集合 $S$ 为空集，即 $a=0$ 时，显然满足条件 $S\subset P$，故 $a=0$.}
\step{若集合 $S$ 非空集，即 $a\neq0$，此时 $S=\left\{-\dfrac1a\right\}$，}
\step{若 $-\dfrac1a=-3$，则 $a=\dfrac13$，若 $-\dfrac1a=2$，则 $a=-\dfrac12$；}
\step{故 $a$ 的取值为 $0$、$\dfrac13$、$-\dfrac12$.}
\solend
\fi

\item 已知 $p:\begin{cases}a>4\\b>4\end{cases}$，$q:\begin{cases}a+b>8\\ab>16\end{cases}$，则 $p$ 是 $q$ 的\blankline[4.4em]条件.

\ifanswer
\solbegin
\step{若 $a>4$，$b>4$，则 $a+b>4+4=8$，$ab>4\times4=16$，}
\step{所以 $p$ 成立一定能推出 $q$ 成立，充分性成立，}
\step{取 $a=2$，$b=9$，$a+b=11>8$，$ab=18>16$，满足 $q$；}
\step{但是 $a=2<4$，不满足 $p$，所以由 $q$ 不能推出 $p$，必要性不成立，}
\step{综上所述，$p$ 是 $q$ 的充分不必要条件.}
\solend
\fi

\item 若关于 $x$ 的不等式 $ax^2-ax+1>0$ 对一切实数 $x$ 恒成立，则实数 $a$ 的取值范围是\blankline[4.4em].

\ifanswer
\solbegin
\step{当 $a=0$ 时，不等式即 $1>0$，满足条件.}
\step{当 $a\neq0$ 时，要使不等式 $ax^2+ax+1>0$ 对一切 $x\in\R$ 恒成立，}
\step{需 $\begin{cases}a>0\\\Delta=a^2-4a<0\end{cases}$，解得 $0<a<4$.}
\step{综上，实数 $a$ 的取值范围是 $[0,4)$.}
\solend
\fi

\item 已知集合 $M=\{x\mid (x-a)(x^2-ax+a-1)=0\}$ 各元素之和等于 3，则实数 $a=$\blankline[4.4em].

\ifanswer
\solbegin
\step{方程 $(x-a)(x^2-ax+a-1)=0$，化简得 $(x-a)(x-1)[x-(a-1)]=0$，}
\step{解得 $x_1=a$，$x_2=1$，$x_3=a-1$，}
\step{因为它们是集合 $M$ 中的元素，所以存在互异性.}
\step{当 $x_1\neq x_2\neq x_3$ 时，即 $a+1+a-1=3$，解得 $a=\dfrac32$，此时 $M=\dfrac32,1,\dfrac12$，满足.}
\step{当 $a=1$ 时，$1+0\neq3$，不合题意；}
\step{当 $a=a-1$ 时，不符合题意；}
\step{当 $1=a-1$ 时，$a=2$，此时 $M=\{2,1\}$ 满足题意.}
\step{故 $a$ 的取值为 $2$ 或 $\dfrac32$.}
\solend
\fi

\item 已知 $a>b>1$ 且 $a^2-6a-2m+1=0$，$b^2-6b-2m+1=0$，则实数 $m$ 的取值范围是\blankline[4.4em].

\ifanswer
\solbegin
\step{由题意得 $a$、$b$ 是函数 $f(x)=x^2-6x-2m+1$ 两个大于 1 的不同零点，}
\step{所以 $\begin{cases}f(1)>0\\f(3)<0\end{cases}$，即 $\begin{cases}-4-2m>0\\-8-2m<0\end{cases}$，解得 $-4<m<-2$，}
\step{故 $m$ 的取值范围为 $(-4,-2)$.}
\solend
\fi

\end{enumerate}

\bigsec{二、选择题}

\begin{enumerate}
\setcounter{enumi}{12}

\item “$x<2$” 是 “$x^2<4$” 的\mbox{（\quad）}条件

\keepwithprev
A. 充分不必要\qquad\qquad B. 必要不充分

C. 充要\qquad\qquad\qquad\quad D. 既不充分也不必要

\ifanswer
\solbegin
\step{解 $x^2<4$，得 $-2<x<2$，}
\step{由 $x<2$ 得不出 $-2<x<2$，而由 $-2<x<2$ 能得出 $x<2$，}
\step{所以 “$x<2$” 是 “$x^2<4$” 的必要不充分条件. 故选 $B$.}
\solend
\fi

\item 已知集合 $A=\{x\mid x<-2\}$，$B=\{x\mid x\geqslant a\}$，且 $A\cup B=\R$，则实数 $a$ 的取值范围是\mbox{（\quad）}

\keepwithprev
A. $a\leqslant-2$\qquad\qquad B. $a<-2$\qquad\qquad C. $a\geqslant-2$\qquad\qquad D. $a>-2$

\ifanswer
\solbegin
\step{集合 $A=\{x\mid x<-2\}$，$B=\{x\mid x\geqslant a\}$，且 $A\cup B=\R$，}
\step{则实数 $a$ 的取值范围是 $a\leqslant-2$. 故选 $A$.}
\solend
\fi

\item 下列关于 $x,y$ 的方程组 $\begin{cases}y=kx+1\\y=2kx+3\end{cases}$ 的解的说法中正确的是\mbox{（\quad）}

\keepwithprev
A. 该方程一定有唯一解\qquad\qquad B. 该方程没有解

C. $k\neq0$ 时，方程组有唯一解\qquad\qquad D. $k=0$ 时，方程组有无数组解

\ifanswer
\solbegin
\step{由题意得 $kx+1=2kx+3$，即 $kx=-2$，}
\step{当 $k=0$ 时，$0=-2$ 不成立，方程组无解；}
\step{当 $k\neq0$ 时，$x=-\dfrac2k$，方程组有唯一解. 故选 $C$.}
\solend
\fi

\item 设非空集合 $S=\{x\mid m\leqslant x\leqslant n\}$ 满足：当 $x\in S$ 时，有 $x^2\in S$. 给出如下三个结论：①若 $m=1$，则 $S=\{1\}$；②若 $m=-\dfrac12$，则 $\dfrac14\leqslant n\leqslant1$；③若 $n=\dfrac12$，则 $-\dfrac{\sqrt2}{2}\leqslant m\leqslant0$. 其中正确命题的个数为\mbox{（\quad）}

\keepwithprev
A. $3$\qquad\qquad B. $2$\qquad\qquad C. $1$\qquad\qquad D. $0$

\ifanswer
\solbegin
\step{由定义设非空集合 $S=\{x\mid m\leqslant x\leqslant n\}$ 满足当 $x\in S$ 时，有 $x^2\in S$，}
\step{符合定义的参数 $m$ 的值一定大于等于 1 或小于等于 0，}
\step{惟如此才能保证 $m\in S$ 时，有 $m^2\in S$ 即 $m^2\geqslant m$，}
\step{符合条件的 $n$ 的值一定大于等于 0，小于等于 1，}
\step{惟如此才能保证 $n\in S$ 时，有 $n^2\in S$ 即 $n^2\leqslant n$，对各个命题进行判断：}
\step{对于① $m=1$，$m^2=1\in S$，故必有 $\begin{cases}n^2\leqslant n\\n\geqslant1\end{cases}$，得 $n=1$，$S=\{1\}$，}
\step{② $m=-\dfrac12$，$m^2=\dfrac14\in S$，则 $\begin{cases}n^2\leqslant n\\\dfrac14\leqslant n\end{cases}$，解得 $\dfrac14\leqslant n\leqslant1$；}
\step{对于③，若 $n=\dfrac12$，则 $\begin{cases}\dfrac12\geqslant m\\m^2\geqslant m\\\dfrac12\geqslant m^2\end{cases}$，解得 $-\dfrac{\sqrt2}{2}\leqslant m\leqslant0$，}
\step{所以正确命题有 3 个. 故选 $A$.}
\solend
\fi

\end{enumerate}

\bigsec{三、解答题}

\begin{enumerate}
\setcounter{enumi}{16}

\item 已知 $a$、$b$ 是实数.

（1）求证：$a^2+b^2\geqslant2a-2b-2$，并指出等号成立的条件；

（2）若 $ab=1$，求 $a^2+4b^2$ 的最小值.

\ifanswer
\solbegin
\step{（1）因为 $a^2+b^2-(2a-2b-2)=a^2+b^2-2a+2b+2$\dots\dots2 分}
\step{$=(a-1)^2+(b+1)^2\geqslant0$\dots\dots4 分}
\step{所以 $a^2+b^2\geqslant2a-2b-2$.}
\step{当且仅当 $a=1,b=-1$ 时，不等式中等号成立.\dots\dots7 分}
\step{（2）$a^2+4b^2=a^2+(2b)^2\geqslant2\cdot a\cdot(2b)=4ab=4$，\dots\dots5 分}
\step{当且仅当 $a=2b$，即 $\begin{cases}a=\sqrt2\\b=\dfrac{\sqrt2}{2}\end{cases}$ 或 $\begin{cases}a=-\sqrt2\\b=-\dfrac{\sqrt2}{2}\end{cases}$ 时，不等式中等号成立，}
\step{所以 $a^2+4b^2$ 的最小值为 4.\dots\dots7 分}
\solend
\fi

\item 已知关于 $x$ 的方程 $x^2+(m-3)x+2=0$，其中 $m\in\R$，有两个实数根 $\alpha$、$\beta$.

（1）若 $m=8$，求 $\dfrac1\alpha+\dfrac1\beta$ 的值；

（2）若 $\alpha^2+\beta^2=\alpha\beta+3$，求实数 $m$ 的值.

\ifanswer
\solbegin
\step{（1）若 $m=8$，则关于 $x$ 的方程为 $x^2+5x+2=0$，}
\step{则 $\alpha+\beta=-5$，$\alpha\beta=2$，所以 $\dfrac1\alpha+\dfrac1\beta=\dfrac{\alpha+\beta}{\alpha\beta}=-\dfrac52$；}
\step{（2）由题意得 $\alpha+\beta=-(m-3)=3-m$，$\alpha\beta=2$，}
\step{又 $\alpha^2+\beta^2=\alpha\beta+3$，则 $(\alpha+\beta)^2-2\alpha\beta=\alpha\beta+3$，}
\step{即 $(\alpha+\beta)^2=3\alpha\beta+3$，所以 $(3-m)^2=3\times2+3=9$，解得 $m=0$ 或 $6$.}
\solend
\fi

\item 设集合 $A=\{x\mid\lvert x-a\rvert<2\}$，$B=\left\{x\mid\dfrac{x+4}{x+2}>2\right\}$.

（1）若 $a=1$，试用区间表示集合 $A$、$B$；

（2）若 $A\cup\overline{B}=\R$，求实数 $a$ 的取值范围.

\ifanswer
\solbegin
\step{（1）当 $a=1$ 时，由 $\lvert x-1\rvert<2$ 得 $-2<x-1<2$，解得 $-1<x<3$，}
\step{所以 $A=(-1,3)$.}
\step{由 $\dfrac{x+4}{x+2}>2$ 得 $2-\dfrac{x+4}{x+2}=\dfrac{x}{x+2}<0$，则有 $x(x+2)<0$，解得 $-2<x<0$，}
\step{所以 $B=(-2,0)$；}
\step{（2）由 $\lvert x-a\rvert<2$ 得 $-2<x-a<2$，解得 $a-2<x<a+2$，}
\step{所以 $A=(a-2,a+2)$，}
\step{由（1）得 $B=(-2,0)$，则 $\overline{B}=\{x\mid x\geqslant0$ 或 $x\leqslant-2\}$，}
\step{由于 $A\cup\overline{B}=\R$，所以 $\begin{cases}a+2\geqslant0\\a-2\leqslant-2\end{cases}$，解得 $-2\leqslant a\leqslant0$，}
\step{所以实数 $a$ 的取值范围是 $[-2,0]$.}
\solend
\fi

\item 某服装公司生产的衬衫每件定价为 80 元，在某城市年销售 8 万件. 现该公司计划在该市招收代理商来销售衬衫，以降低管理和营销成本. 已知代理商收取的代理费为总销售金额的 $r\%$（其中 $0<r\leqslant10$）.

（1）若单件衬衫的利润保持不变，求招代理商后每件衬衫的定价为多少元（用 $r$ 表示）；

（2）已知将每件衬衫的价格提高到 $\dfrac{80}{1-r\%}$ 元后，每年的销售量会减少 $0.62r$ 万件. 若代理商每年收取的代理费不少于 16 万元，求 $r$ 的取值范围.

\ifanswer
\solbegin
\step{（1）设招代理商后每件衬衫的定价为 $x$ 元，代理商收取总销售金额的 $r\%$，}
\step{因此公司每件衬衫的实际净收入为 $x(1-r\%)$ 元，}
\step{单件利润保持不变，等价于净收入仍为 80 元，}
\step{于是 $x(1-r\%)=80$，解得 $x=\dfrac{80}{1-r\%}$，因为 $0<r\leqslant10$，有 $1-r\%\geqslant0.9>0$，}
\step{分母恒不为零，该定价有意义，所以招代理商后每件衬衫的定价为 $\dfrac{80}{1-r\%}$ 元；}
\step{（2）价格提高到 $\dfrac{80}{1-r\%}$ 元，年销售量减少 $0.62r$ 万件，}
\step{故年销售量为 $8-0.62r$ 万件，$8-0.62r>0$，}
\step{由题意得 $\dfrac{80}{1-r\%}\cdot(8-0.62r)\cdot r\%\geqslant16$，}
\step{即 $\dfrac{80}{1-\dfrac r{100}}\cdot(8-0.62r)\cdot\dfrac r{100}=\dfrac{80r(8-0.62r)}{100-r}\geqslant16$，}
\step{由 $0<r\leqslant10$ 得 $100-r>0$，整理得 $80r(8-0.62r)\geqslant16(100-r)$，}
\step{整理得 $31r^2-410r+1000\leqslant0$，解得 $\dfrac{100}{31}\leqslant r\leqslant10$，}
\step{当 $\dfrac{100}{31}\leqslant r\leqslant10$ 时，$8-0.62r>0$ 且 $1-r\%>0$，}
\step{均符合实际经营条件，所以 $r$ 的取值范围是 $\left[\dfrac{100}{31},10\right]$.}
\solend
\fi

\item 定义 $\min\{a_1,a_2,\cdots,a_n\}$ 为 $n$ 个实数 $a_1,a_2,\cdots,a_n$ 中的最小数，$\max\{a_1,a_2,\cdots,a_n\}$ 为 $n$ 个实数 $a_1,a_2,\cdots,a_n$ 中的最大数.

（1）设 $a,b$ 都是正实数，且 $a+b=1$，求 $\max\left\{ab,\dfrac14\right\}$；

（2）解不等式：$\min\{x+1,x^2+3,\lvert x-1\rvert\}>2x-3$；

（3）设 $a,b$ 都是正实数，求 $\max\left\{a+\dfrac1b,\dfrac2a+b\right\}$ 的最小值.

\ifanswer
\solbegin
\step{（1）$ab\leqslant\dfrac{(a+b)^2}{4}=\dfrac14$，所以 $\max\left\{ab,\dfrac14\right\}=\dfrac14$.}
\step{（2）因为 $x^2+3>x+1$，$x^2+3>\lvert x-1\rvert$，}
\step{所以 $\min\{x+1,x^2+3,\lvert x-1\rvert\}=\min\{x+1,\lvert x-1\rvert\}$.}
\step{当 $x\geqslant0$ 时，$\lvert x-1\rvert\leqslant x+1$，所以 $\lvert x-1\rvert>2x-3$，解得 $0\leqslant x<2$；}
\step{当 $x<0$ 时，$\lvert x-1\rvert>x+1$，所以 $x+1>2x-3$，解得 $x<0$.}
\step{综上，不等式的解集为 $(-\infty,2)$.}
\step{（3）设 $M=\max\left\{a+\dfrac1b,\dfrac2a+b\right\}$，则 $M\geqslant a+\dfrac1b$，$M\geqslant\dfrac2a+b$，}
\step{所以 $M^2\geqslant\left(a+\dfrac1b\right)\left(\dfrac2a+b\right)=3+ab+\dfrac2{ab}\geqslant3+2\sqrt2$，}
\step{所以 $M$ 的最小值为 $\sqrt2+1$，当 $a=\sqrt2$，$b=1$ 时取等号.}
\solend
\fi

\end{enumerate}

\end{document}
"
% 紧凑版：xelatex "\def\iscompact{1}% 高一126_建青实验_期中.tex
%   —— 高一上数学"2026-2027 年建青实验高一上期中"（试卷版/答案版）
% 试卷版：xelatex 高一126_建青实验_期中.tex
% 答案版：xelatex "\def\isanswer{1}% 高一126_建青实验_期中.tex
%   —— 高一上数学"2026-2027 年建青实验高一上期中"（试卷版/答案版）
% 试卷版：xelatex 高一126_建青实验_期中.tex
% 答案版：xelatex "\def\isanswer{1}\input{高一126_建青实验_期中.tex}"
% 紧凑版：xelatex "\def\iscompact{1}\input{高一126_建青实验_期中.tex}"
%
% 原始材料：微信公众号"魔都数学加油站"文章，作者破晓，2026-11-11 18:35 发布；
% 连载"27学年高一126"，原文标题"27学年高一126——2026-2027年上海市建青实验高一上期中"；
% 原文链接：https://mp.weixin.qq.com/s/HpxpBJu4d_wOnUCIbZSQiA。
% 正文为 8 张 PNG 页（01、05、08 为 2560x3620，02、03、04、06、07 为 4762x6735），
% 题目下印有【解析】或【答案】，为讲评版。题目、解析均据原图转录；
% 卷面粉红色水印"魔都数学加油站"及公众号页脚未录入。全卷无插图。
%
% 卷首标题：原卷印作"2026-2027 年建青实验高一上期中"（公众号标题中的"上海市"
% 卷面标题行没有，本稿照卷面），年份区间按全稿惯例排 en dash；原卷未印副标题，
% 本卷为期中卷、跨集合与常用逻辑用语／不等式／函数三章，按题目内容标为
% "集合与常用逻辑用语、不等式、函数"（与 125 同为期中卷，口径相同）：
% 集合与逻辑见 Q5、Q8、Q9、Q11、Q13、Q14、Q16、Q19，不等式见 Q6、Q7、Q10、
% Q17、Q20、Q21，函数（含指数对数、二次函数零点）见 Q3、Q4、Q12、Q15、Q18。
%
% 与原卷的排版差异：
%   ① 原文从第 3 题起，填空题编号为 3—12、选择题为 13—16、解答题为 17—21，本稿保留原题号；
%   ② 原卷每题下直接印【解析】／【答案】，本稿按全稿惯例置于答案版；第 4、5 题只印
%      【答案】，另提 \ans{}；第 13—16 题答案均在【解析】末句"故选 …"，题干括号内留空、
%      不另提 \ans{}；第 3 题印一行【解析】，按 \solbegin+\step 排；第 17 题解析保留
%      原卷评分点（照 高一54／64 成例排 \dots\dots N 分）；
%   ③ 原卷未印副标题，本稿按"知识模块"口径补出（见上"卷首标题"）；
%   ④ 原卷填空横线按全稿惯例排为 \blankline[4.4em]；选择题括号排 \mbox{（\quad）}、
%      选项段前按惯例加 \keepwithprev；第 14 题原卷括号另起一行，本稿并入题干行末；
%      解答题不留作答区（原卷为讲评版，与 109—125 同）；
%   ⑤ 原卷数集排斜体 $R$，按全稿惯例改 \R；绝对值排 \lvert…\rvert，集合描述竖线排 \mid；
%      补集记号照原卷上划线排 \overline{B}；不等号按全稿惯例排 \leqslant／\geqslant；
%      子题标号原卷一律用全角括号（1）（2）（3），本稿照录（\step 行内亦同，照 高一121 成例）。
%
% 判别式字形：第 10 题解析 $\Delta=a^2-4a<0$ 原卷排 $\Delta$（等腰、高），
%   照录 \Delta（与 高一124／125 同题排 $\triangle$（等边、扁）不同，逐处照本卷字形）。
%
% 子集符号：第 8 题题干与解析两处 $S\subset P$ 均印单线、无下划线，照录 \subset；
%   全卷未出现带下划线的 \subseteq。
%
% 原卷存疑照录：
%   第 10 题题干作 $ax^2-ax+1>0$，解析却作 $ax^2+ax+1>0$（一次项符号不一致），照录、未统一；
%   第 11 题解析"此时 $M=\dfrac32,1,\dfrac12$，满足"未用集合花括号（题干 $M$ 为集合），照录；
%   第 11 题解析"因为它们是集合 $M$ 中的元素，所以存在互异性"（"存在互异性"疑为
%     "满足互异性"之误），又"当 $x_1\neq x_2\neq x_3$ 时"连写两个不等号，均照录；
%   第 16 题解析两处"惟如此才能保证"（用"惟"字）及"符合定义的参数 $m$ 的值一定大于等于 1
%     或小于等于 0"的措辞，照录；
%   第 17 题 (1) 解析"当且仅当 $a=1,b=-1$ 时"用半角逗号、第 21 题 (3) 解析
%     "当 $a=\sqrt2$，$b=1$ 时取等号"用全角逗号，均照录；
%   第 19 题解析由 $\dfrac{x+4}{x+2}>2$ 得 "$2-\dfrac{x+4}{x+2}=\dfrac{x}{x+2}<0$"
%     （以 2 减左式，与常见的"左式减 2"次序相反），照录；
%   第 15 题选项 A、B 作"该方程…"、C、D 作"方程组…"，照录；
%   第 20 题题干作"招收代理商"，(1) 却作"招代理商后"，照录。
%
% 行数复核：逐题数 \step 条数，合计 91 条（Q3：1；Q4：0；Q5：0；Q6：3；Q7：2；
% Q8：5；Q9：5；Q10：4；Q11：8；Q12：3；Q13：3；Q14：2；Q15：3；Q16：9；
% Q17：7；Q18：5；Q19：9；Q20：13；Q21：9），与原卷印刷行逐题相等
% （cases 块作一个印刷行计）。第 4、5 题只印【答案】，提 \ans{}，不计入 \step 账；
% 四道选择题答案均在解析末句，括号内留空、不另提 \ans{}。
%
\documentclass[a4paper,11pt]{article}
\usepackage{common}

\begin{document}

\examtitlest[3]{2026--2027 年建青实验高一上期中}{集合与常用逻辑用语、不等式、函数}

\bigsec{一、填空题}

\begin{enumerate}
\setcounter{enumi}{2}

\item 已知 $\log_2 x=3$，则 $x=$\blankline[4.4em].

\ifanswer
\solbegin
\step{$x=2^3=8$.}
\solend
\fi

\item 将 $\sqrt{a^2\sqrt a}$（其中 $a>0$）化为有理数指数幂的形式为\blankline[4.4em].

\ifanswer
\ans{$a^{\frac54}$}
\fi

\item 用反证法证明命题：“已知 $m$、$n\in\R$，若 $m+n\leqslant5$，则 $m\leqslant2$ 或 $n\leqslant3$” 时应假设\blankline[4.4em].

\ifanswer
\ans{$m>2$ 且 $n>3$}
\fi

\item 若 $x\in\R$，则 $\lvert x+1\rvert+\lvert x-1\rvert$ 的最小值为\blankline[4.4em].

\ifanswer
\solbegin
\step{$\lvert x+1\rvert+\lvert x-1\rvert\geqslant\lvert(x+1)-(x-1)\rvert=2$，}
\step{当且仅当 $(x+1)(x-1)\leqslant0$，即 $-1\leqslant x\leqslant1$ 时取等号，}
\step{所以 $\lvert x+1\rvert+\lvert x-1\rvert$ 的最小值为 2.}
\solend
\fi

\item 已知 $x>1$，则 $x+\dfrac1{x-1}$ 的最小值是\blankline[4.4em].

\ifanswer
\solbegin
\step{$x-1>0$，则 $x+\dfrac1{x-1}=(x-1)+\dfrac1{x-1}+1\geqslant2\sqrt{(x-1)\cdot\dfrac1{x-1}}+1=3$，}
\step{当且仅当 $x=2$ 时取等号，故 $x+\dfrac1{x-1}$ 的最小值为 3.}
\solend
\fi

\item 已知集合 $P=\{x\mid x^2+x-6=0\}$，$S=\{x\mid ax+1=0\}$ 且 $S\subset P$，则 $a$ 的值为\blankline[4.4em].

\ifanswer
\solbegin
\step{集合 $P=\{-3,2\}$，集合 $S$ 中至多有一个元素，}
\step{若集合 $S$ 为空集，即 $a=0$ 时，显然满足条件 $S\subset P$，故 $a=0$.}
\step{若集合 $S$ 非空集，即 $a\neq0$，此时 $S=\left\{-\dfrac1a\right\}$，}
\step{若 $-\dfrac1a=-3$，则 $a=\dfrac13$，若 $-\dfrac1a=2$，则 $a=-\dfrac12$；}
\step{故 $a$ 的取值为 $0$、$\dfrac13$、$-\dfrac12$.}
\solend
\fi

\item 已知 $p:\begin{cases}a>4\\b>4\end{cases}$，$q:\begin{cases}a+b>8\\ab>16\end{cases}$，则 $p$ 是 $q$ 的\blankline[4.4em]条件.

\ifanswer
\solbegin
\step{若 $a>4$，$b>4$，则 $a+b>4+4=8$，$ab>4\times4=16$，}
\step{所以 $p$ 成立一定能推出 $q$ 成立，充分性成立，}
\step{取 $a=2$，$b=9$，$a+b=11>8$，$ab=18>16$，满足 $q$；}
\step{但是 $a=2<4$，不满足 $p$，所以由 $q$ 不能推出 $p$，必要性不成立，}
\step{综上所述，$p$ 是 $q$ 的充分不必要条件.}
\solend
\fi

\item 若关于 $x$ 的不等式 $ax^2-ax+1>0$ 对一切实数 $x$ 恒成立，则实数 $a$ 的取值范围是\blankline[4.4em].

\ifanswer
\solbegin
\step{当 $a=0$ 时，不等式即 $1>0$，满足条件.}
\step{当 $a\neq0$ 时，要使不等式 $ax^2+ax+1>0$ 对一切 $x\in\R$ 恒成立，}
\step{需 $\begin{cases}a>0\\\Delta=a^2-4a<0\end{cases}$，解得 $0<a<4$.}
\step{综上，实数 $a$ 的取值范围是 $[0,4)$.}
\solend
\fi

\item 已知集合 $M=\{x\mid (x-a)(x^2-ax+a-1)=0\}$ 各元素之和等于 3，则实数 $a=$\blankline[4.4em].

\ifanswer
\solbegin
\step{方程 $(x-a)(x^2-ax+a-1)=0$，化简得 $(x-a)(x-1)[x-(a-1)]=0$，}
\step{解得 $x_1=a$，$x_2=1$，$x_3=a-1$，}
\step{因为它们是集合 $M$ 中的元素，所以存在互异性.}
\step{当 $x_1\neq x_2\neq x_3$ 时，即 $a+1+a-1=3$，解得 $a=\dfrac32$，此时 $M=\dfrac32,1,\dfrac12$，满足.}
\step{当 $a=1$ 时，$1+0\neq3$，不合题意；}
\step{当 $a=a-1$ 时，不符合题意；}
\step{当 $1=a-1$ 时，$a=2$，此时 $M=\{2,1\}$ 满足题意.}
\step{故 $a$ 的取值为 $2$ 或 $\dfrac32$.}
\solend
\fi

\item 已知 $a>b>1$ 且 $a^2-6a-2m+1=0$，$b^2-6b-2m+1=0$，则实数 $m$ 的取值范围是\blankline[4.4em].

\ifanswer
\solbegin
\step{由题意得 $a$、$b$ 是函数 $f(x)=x^2-6x-2m+1$ 两个大于 1 的不同零点，}
\step{所以 $\begin{cases}f(1)>0\\f(3)<0\end{cases}$，即 $\begin{cases}-4-2m>0\\-8-2m<0\end{cases}$，解得 $-4<m<-2$，}
\step{故 $m$ 的取值范围为 $(-4,-2)$.}
\solend
\fi

\end{enumerate}

\bigsec{二、选择题}

\begin{enumerate}
\setcounter{enumi}{12}

\item “$x<2$” 是 “$x^2<4$” 的\mbox{（\quad）}条件

\keepwithprev
A. 充分不必要\qquad\qquad B. 必要不充分

C. 充要\qquad\qquad\qquad\quad D. 既不充分也不必要

\ifanswer
\solbegin
\step{解 $x^2<4$，得 $-2<x<2$，}
\step{由 $x<2$ 得不出 $-2<x<2$，而由 $-2<x<2$ 能得出 $x<2$，}
\step{所以 “$x<2$” 是 “$x^2<4$” 的必要不充分条件. 故选 $B$.}
\solend
\fi

\item 已知集合 $A=\{x\mid x<-2\}$，$B=\{x\mid x\geqslant a\}$，且 $A\cup B=\R$，则实数 $a$ 的取值范围是\mbox{（\quad）}

\keepwithprev
A. $a\leqslant-2$\qquad\qquad B. $a<-2$\qquad\qquad C. $a\geqslant-2$\qquad\qquad D. $a>-2$

\ifanswer
\solbegin
\step{集合 $A=\{x\mid x<-2\}$，$B=\{x\mid x\geqslant a\}$，且 $A\cup B=\R$，}
\step{则实数 $a$ 的取值范围是 $a\leqslant-2$. 故选 $A$.}
\solend
\fi

\item 下列关于 $x,y$ 的方程组 $\begin{cases}y=kx+1\\y=2kx+3\end{cases}$ 的解的说法中正确的是\mbox{（\quad）}

\keepwithprev
A. 该方程一定有唯一解\qquad\qquad B. 该方程没有解

C. $k\neq0$ 时，方程组有唯一解\qquad\qquad D. $k=0$ 时，方程组有无数组解

\ifanswer
\solbegin
\step{由题意得 $kx+1=2kx+3$，即 $kx=-2$，}
\step{当 $k=0$ 时，$0=-2$ 不成立，方程组无解；}
\step{当 $k\neq0$ 时，$x=-\dfrac2k$，方程组有唯一解. 故选 $C$.}
\solend
\fi

\item 设非空集合 $S=\{x\mid m\leqslant x\leqslant n\}$ 满足：当 $x\in S$ 时，有 $x^2\in S$. 给出如下三个结论：①若 $m=1$，则 $S=\{1\}$；②若 $m=-\dfrac12$，则 $\dfrac14\leqslant n\leqslant1$；③若 $n=\dfrac12$，则 $-\dfrac{\sqrt2}{2}\leqslant m\leqslant0$. 其中正确命题的个数为\mbox{（\quad）}

\keepwithprev
A. $3$\qquad\qquad B. $2$\qquad\qquad C. $1$\qquad\qquad D. $0$

\ifanswer
\solbegin
\step{由定义设非空集合 $S=\{x\mid m\leqslant x\leqslant n\}$ 满足当 $x\in S$ 时，有 $x^2\in S$，}
\step{符合定义的参数 $m$ 的值一定大于等于 1 或小于等于 0，}
\step{惟如此才能保证 $m\in S$ 时，有 $m^2\in S$ 即 $m^2\geqslant m$，}
\step{符合条件的 $n$ 的值一定大于等于 0，小于等于 1，}
\step{惟如此才能保证 $n\in S$ 时，有 $n^2\in S$ 即 $n^2\leqslant n$，对各个命题进行判断：}
\step{对于① $m=1$，$m^2=1\in S$，故必有 $\begin{cases}n^2\leqslant n\\n\geqslant1\end{cases}$，得 $n=1$，$S=\{1\}$，}
\step{② $m=-\dfrac12$，$m^2=\dfrac14\in S$，则 $\begin{cases}n^2\leqslant n\\\dfrac14\leqslant n\end{cases}$，解得 $\dfrac14\leqslant n\leqslant1$；}
\step{对于③，若 $n=\dfrac12$，则 $\begin{cases}\dfrac12\geqslant m\\m^2\geqslant m\\\dfrac12\geqslant m^2\end{cases}$，解得 $-\dfrac{\sqrt2}{2}\leqslant m\leqslant0$，}
\step{所以正确命题有 3 个. 故选 $A$.}
\solend
\fi

\end{enumerate}

\bigsec{三、解答题}

\begin{enumerate}
\setcounter{enumi}{16}

\item 已知 $a$、$b$ 是实数.

（1）求证：$a^2+b^2\geqslant2a-2b-2$，并指出等号成立的条件；

（2）若 $ab=1$，求 $a^2+4b^2$ 的最小值.

\ifanswer
\solbegin
\step{（1）因为 $a^2+b^2-(2a-2b-2)=a^2+b^2-2a+2b+2$\dots\dots2 分}
\step{$=(a-1)^2+(b+1)^2\geqslant0$\dots\dots4 分}
\step{所以 $a^2+b^2\geqslant2a-2b-2$.}
\step{当且仅当 $a=1,b=-1$ 时，不等式中等号成立.\dots\dots7 分}
\step{（2）$a^2+4b^2=a^2+(2b)^2\geqslant2\cdot a\cdot(2b)=4ab=4$，\dots\dots5 分}
\step{当且仅当 $a=2b$，即 $\begin{cases}a=\sqrt2\\b=\dfrac{\sqrt2}{2}\end{cases}$ 或 $\begin{cases}a=-\sqrt2\\b=-\dfrac{\sqrt2}{2}\end{cases}$ 时，不等式中等号成立，}
\step{所以 $a^2+4b^2$ 的最小值为 4.\dots\dots7 分}
\solend
\fi

\item 已知关于 $x$ 的方程 $x^2+(m-3)x+2=0$，其中 $m\in\R$，有两个实数根 $\alpha$、$\beta$.

（1）若 $m=8$，求 $\dfrac1\alpha+\dfrac1\beta$ 的值；

（2）若 $\alpha^2+\beta^2=\alpha\beta+3$，求实数 $m$ 的值.

\ifanswer
\solbegin
\step{（1）若 $m=8$，则关于 $x$ 的方程为 $x^2+5x+2=0$，}
\step{则 $\alpha+\beta=-5$，$\alpha\beta=2$，所以 $\dfrac1\alpha+\dfrac1\beta=\dfrac{\alpha+\beta}{\alpha\beta}=-\dfrac52$；}
\step{（2）由题意得 $\alpha+\beta=-(m-3)=3-m$，$\alpha\beta=2$，}
\step{又 $\alpha^2+\beta^2=\alpha\beta+3$，则 $(\alpha+\beta)^2-2\alpha\beta=\alpha\beta+3$，}
\step{即 $(\alpha+\beta)^2=3\alpha\beta+3$，所以 $(3-m)^2=3\times2+3=9$，解得 $m=0$ 或 $6$.}
\solend
\fi

\item 设集合 $A=\{x\mid\lvert x-a\rvert<2\}$，$B=\left\{x\mid\dfrac{x+4}{x+2}>2\right\}$.

（1）若 $a=1$，试用区间表示集合 $A$、$B$；

（2）若 $A\cup\overline{B}=\R$，求实数 $a$ 的取值范围.

\ifanswer
\solbegin
\step{（1）当 $a=1$ 时，由 $\lvert x-1\rvert<2$ 得 $-2<x-1<2$，解得 $-1<x<3$，}
\step{所以 $A=(-1,3)$.}
\step{由 $\dfrac{x+4}{x+2}>2$ 得 $2-\dfrac{x+4}{x+2}=\dfrac{x}{x+2}<0$，则有 $x(x+2)<0$，解得 $-2<x<0$，}
\step{所以 $B=(-2,0)$；}
\step{（2）由 $\lvert x-a\rvert<2$ 得 $-2<x-a<2$，解得 $a-2<x<a+2$，}
\step{所以 $A=(a-2,a+2)$，}
\step{由（1）得 $B=(-2,0)$，则 $\overline{B}=\{x\mid x\geqslant0$ 或 $x\leqslant-2\}$，}
\step{由于 $A\cup\overline{B}=\R$，所以 $\begin{cases}a+2\geqslant0\\a-2\leqslant-2\end{cases}$，解得 $-2\leqslant a\leqslant0$，}
\step{所以实数 $a$ 的取值范围是 $[-2,0]$.}
\solend
\fi

\item 某服装公司生产的衬衫每件定价为 80 元，在某城市年销售 8 万件. 现该公司计划在该市招收代理商来销售衬衫，以降低管理和营销成本. 已知代理商收取的代理费为总销售金额的 $r\%$（其中 $0<r\leqslant10$）.

（1）若单件衬衫的利润保持不变，求招代理商后每件衬衫的定价为多少元（用 $r$ 表示）；

（2）已知将每件衬衫的价格提高到 $\dfrac{80}{1-r\%}$ 元后，每年的销售量会减少 $0.62r$ 万件. 若代理商每年收取的代理费不少于 16 万元，求 $r$ 的取值范围.

\ifanswer
\solbegin
\step{（1）设招代理商后每件衬衫的定价为 $x$ 元，代理商收取总销售金额的 $r\%$，}
\step{因此公司每件衬衫的实际净收入为 $x(1-r\%)$ 元，}
\step{单件利润保持不变，等价于净收入仍为 80 元，}
\step{于是 $x(1-r\%)=80$，解得 $x=\dfrac{80}{1-r\%}$，因为 $0<r\leqslant10$，有 $1-r\%\geqslant0.9>0$，}
\step{分母恒不为零，该定价有意义，所以招代理商后每件衬衫的定价为 $\dfrac{80}{1-r\%}$ 元；}
\step{（2）价格提高到 $\dfrac{80}{1-r\%}$ 元，年销售量减少 $0.62r$ 万件，}
\step{故年销售量为 $8-0.62r$ 万件，$8-0.62r>0$，}
\step{由题意得 $\dfrac{80}{1-r\%}\cdot(8-0.62r)\cdot r\%\geqslant16$，}
\step{即 $\dfrac{80}{1-\dfrac r{100}}\cdot(8-0.62r)\cdot\dfrac r{100}=\dfrac{80r(8-0.62r)}{100-r}\geqslant16$，}
\step{由 $0<r\leqslant10$ 得 $100-r>0$，整理得 $80r(8-0.62r)\geqslant16(100-r)$，}
\step{整理得 $31r^2-410r+1000\leqslant0$，解得 $\dfrac{100}{31}\leqslant r\leqslant10$，}
\step{当 $\dfrac{100}{31}\leqslant r\leqslant10$ 时，$8-0.62r>0$ 且 $1-r\%>0$，}
\step{均符合实际经营条件，所以 $r$ 的取值范围是 $\left[\dfrac{100}{31},10\right]$.}
\solend
\fi

\item 定义 $\min\{a_1,a_2,\cdots,a_n\}$ 为 $n$ 个实数 $a_1,a_2,\cdots,a_n$ 中的最小数，$\max\{a_1,a_2,\cdots,a_n\}$ 为 $n$ 个实数 $a_1,a_2,\cdots,a_n$ 中的最大数.

（1）设 $a,b$ 都是正实数，且 $a+b=1$，求 $\max\left\{ab,\dfrac14\right\}$；

（2）解不等式：$\min\{x+1,x^2+3,\lvert x-1\rvert\}>2x-3$；

（3）设 $a,b$ 都是正实数，求 $\max\left\{a+\dfrac1b,\dfrac2a+b\right\}$ 的最小值.

\ifanswer
\solbegin
\step{（1）$ab\leqslant\dfrac{(a+b)^2}{4}=\dfrac14$，所以 $\max\left\{ab,\dfrac14\right\}=\dfrac14$.}
\step{（2）因为 $x^2+3>x+1$，$x^2+3>\lvert x-1\rvert$，}
\step{所以 $\min\{x+1,x^2+3,\lvert x-1\rvert\}=\min\{x+1,\lvert x-1\rvert\}$.}
\step{当 $x\geqslant0$ 时，$\lvert x-1\rvert\leqslant x+1$，所以 $\lvert x-1\rvert>2x-3$，解得 $0\leqslant x<2$；}
\step{当 $x<0$ 时，$\lvert x-1\rvert>x+1$，所以 $x+1>2x-3$，解得 $x<0$.}
\step{综上，不等式的解集为 $(-\infty,2)$.}
\step{（3）设 $M=\max\left\{a+\dfrac1b,\dfrac2a+b\right\}$，则 $M\geqslant a+\dfrac1b$，$M\geqslant\dfrac2a+b$，}
\step{所以 $M^2\geqslant\left(a+\dfrac1b\right)\left(\dfrac2a+b\right)=3+ab+\dfrac2{ab}\geqslant3+2\sqrt2$，}
\step{所以 $M$ 的最小值为 $\sqrt2+1$，当 $a=\sqrt2$，$b=1$ 时取等号.}
\solend
\fi

\end{enumerate}

\end{document}
"
% 紧凑版：xelatex "\def\iscompact{1}% 高一126_建青实验_期中.tex
%   —— 高一上数学"2026-2027 年建青实验高一上期中"（试卷版/答案版）
% 试卷版：xelatex 高一126_建青实验_期中.tex
% 答案版：xelatex "\def\isanswer{1}\input{高一126_建青实验_期中.tex}"
% 紧凑版：xelatex "\def\iscompact{1}\input{高一126_建青实验_期中.tex}"
%
% 原始材料：微信公众号"魔都数学加油站"文章，作者破晓，2026-11-11 18:35 发布；
% 连载"27学年高一126"，原文标题"27学年高一126——2026-2027年上海市建青实验高一上期中"；
% 原文链接：https://mp.weixin.qq.com/s/HpxpBJu4d_wOnUCIbZSQiA。
% 正文为 8 张 PNG 页（01、05、08 为 2560x3620，02、03、04、06、07 为 4762x6735），
% 题目下印有【解析】或【答案】，为讲评版。题目、解析均据原图转录；
% 卷面粉红色水印"魔都数学加油站"及公众号页脚未录入。全卷无插图。
%
% 卷首标题：原卷印作"2026-2027 年建青实验高一上期中"（公众号标题中的"上海市"
% 卷面标题行没有，本稿照卷面），年份区间按全稿惯例排 en dash；原卷未印副标题，
% 本卷为期中卷、跨集合与常用逻辑用语／不等式／函数三章，按题目内容标为
% "集合与常用逻辑用语、不等式、函数"（与 125 同为期中卷，口径相同）：
% 集合与逻辑见 Q5、Q8、Q9、Q11、Q13、Q14、Q16、Q19，不等式见 Q6、Q7、Q10、
% Q17、Q20、Q21，函数（含指数对数、二次函数零点）见 Q3、Q4、Q12、Q15、Q18。
%
% 与原卷的排版差异：
%   ① 原文从第 3 题起，填空题编号为 3—12、选择题为 13—16、解答题为 17—21，本稿保留原题号；
%   ② 原卷每题下直接印【解析】／【答案】，本稿按全稿惯例置于答案版；第 4、5 题只印
%      【答案】，另提 \ans{}；第 13—16 题答案均在【解析】末句"故选 …"，题干括号内留空、
%      不另提 \ans{}；第 3 题印一行【解析】，按 \solbegin+\step 排；第 17 题解析保留
%      原卷评分点（照 高一54／64 成例排 \dots\dots N 分）；
%   ③ 原卷未印副标题，本稿按"知识模块"口径补出（见上"卷首标题"）；
%   ④ 原卷填空横线按全稿惯例排为 \blankline[4.4em]；选择题括号排 \mbox{（\quad）}、
%      选项段前按惯例加 \keepwithprev；第 14 题原卷括号另起一行，本稿并入题干行末；
%      解答题不留作答区（原卷为讲评版，与 109—125 同）；
%   ⑤ 原卷数集排斜体 $R$，按全稿惯例改 \R；绝对值排 \lvert…\rvert，集合描述竖线排 \mid；
%      补集记号照原卷上划线排 \overline{B}；不等号按全稿惯例排 \leqslant／\geqslant；
%      子题标号原卷一律用全角括号（1）（2）（3），本稿照录（\step 行内亦同，照 高一121 成例）。
%
% 判别式字形：第 10 题解析 $\Delta=a^2-4a<0$ 原卷排 $\Delta$（等腰、高），
%   照录 \Delta（与 高一124／125 同题排 $\triangle$（等边、扁）不同，逐处照本卷字形）。
%
% 子集符号：第 8 题题干与解析两处 $S\subset P$ 均印单线、无下划线，照录 \subset；
%   全卷未出现带下划线的 \subseteq。
%
% 原卷存疑照录：
%   第 10 题题干作 $ax^2-ax+1>0$，解析却作 $ax^2+ax+1>0$（一次项符号不一致），照录、未统一；
%   第 11 题解析"此时 $M=\dfrac32,1,\dfrac12$，满足"未用集合花括号（题干 $M$ 为集合），照录；
%   第 11 题解析"因为它们是集合 $M$ 中的元素，所以存在互异性"（"存在互异性"疑为
%     "满足互异性"之误），又"当 $x_1\neq x_2\neq x_3$ 时"连写两个不等号，均照录；
%   第 16 题解析两处"惟如此才能保证"（用"惟"字）及"符合定义的参数 $m$ 的值一定大于等于 1
%     或小于等于 0"的措辞，照录；
%   第 17 题 (1) 解析"当且仅当 $a=1,b=-1$ 时"用半角逗号、第 21 题 (3) 解析
%     "当 $a=\sqrt2$，$b=1$ 时取等号"用全角逗号，均照录；
%   第 19 题解析由 $\dfrac{x+4}{x+2}>2$ 得 "$2-\dfrac{x+4}{x+2}=\dfrac{x}{x+2}<0$"
%     （以 2 减左式，与常见的"左式减 2"次序相反），照录；
%   第 15 题选项 A、B 作"该方程…"、C、D 作"方程组…"，照录；
%   第 20 题题干作"招收代理商"，(1) 却作"招代理商后"，照录。
%
% 行数复核：逐题数 \step 条数，合计 91 条（Q3：1；Q4：0；Q5：0；Q6：3；Q7：2；
% Q8：5；Q9：5；Q10：4；Q11：8；Q12：3；Q13：3；Q14：2；Q15：3；Q16：9；
% Q17：7；Q18：5；Q19：9；Q20：13；Q21：9），与原卷印刷行逐题相等
% （cases 块作一个印刷行计）。第 4、5 题只印【答案】，提 \ans{}，不计入 \step 账；
% 四道选择题答案均在解析末句，括号内留空、不另提 \ans{}。
%
\documentclass[a4paper,11pt]{article}
\usepackage{common}

\begin{document}

\examtitlest[3]{2026--2027 年建青实验高一上期中}{集合与常用逻辑用语、不等式、函数}

\bigsec{一、填空题}

\begin{enumerate}
\setcounter{enumi}{2}

\item 已知 $\log_2 x=3$，则 $x=$\blankline[4.4em].

\ifanswer
\solbegin
\step{$x=2^3=8$.}
\solend
\fi

\item 将 $\sqrt{a^2\sqrt a}$（其中 $a>0$）化为有理数指数幂的形式为\blankline[4.4em].

\ifanswer
\ans{$a^{\frac54}$}
\fi

\item 用反证法证明命题：“已知 $m$、$n\in\R$，若 $m+n\leqslant5$，则 $m\leqslant2$ 或 $n\leqslant3$” 时应假设\blankline[4.4em].

\ifanswer
\ans{$m>2$ 且 $n>3$}
\fi

\item 若 $x\in\R$，则 $\lvert x+1\rvert+\lvert x-1\rvert$ 的最小值为\blankline[4.4em].

\ifanswer
\solbegin
\step{$\lvert x+1\rvert+\lvert x-1\rvert\geqslant\lvert(x+1)-(x-1)\rvert=2$，}
\step{当且仅当 $(x+1)(x-1)\leqslant0$，即 $-1\leqslant x\leqslant1$ 时取等号，}
\step{所以 $\lvert x+1\rvert+\lvert x-1\rvert$ 的最小值为 2.}
\solend
\fi

\item 已知 $x>1$，则 $x+\dfrac1{x-1}$ 的最小值是\blankline[4.4em].

\ifanswer
\solbegin
\step{$x-1>0$，则 $x+\dfrac1{x-1}=(x-1)+\dfrac1{x-1}+1\geqslant2\sqrt{(x-1)\cdot\dfrac1{x-1}}+1=3$，}
\step{当且仅当 $x=2$ 时取等号，故 $x+\dfrac1{x-1}$ 的最小值为 3.}
\solend
\fi

\item 已知集合 $P=\{x\mid x^2+x-6=0\}$，$S=\{x\mid ax+1=0\}$ 且 $S\subset P$，则 $a$ 的值为\blankline[4.4em].

\ifanswer
\solbegin
\step{集合 $P=\{-3,2\}$，集合 $S$ 中至多有一个元素，}
\step{若集合 $S$ 为空集，即 $a=0$ 时，显然满足条件 $S\subset P$，故 $a=0$.}
\step{若集合 $S$ 非空集，即 $a\neq0$，此时 $S=\left\{-\dfrac1a\right\}$，}
\step{若 $-\dfrac1a=-3$，则 $a=\dfrac13$，若 $-\dfrac1a=2$，则 $a=-\dfrac12$；}
\step{故 $a$ 的取值为 $0$、$\dfrac13$、$-\dfrac12$.}
\solend
\fi

\item 已知 $p:\begin{cases}a>4\\b>4\end{cases}$，$q:\begin{cases}a+b>8\\ab>16\end{cases}$，则 $p$ 是 $q$ 的\blankline[4.4em]条件.

\ifanswer
\solbegin
\step{若 $a>4$，$b>4$，则 $a+b>4+4=8$，$ab>4\times4=16$，}
\step{所以 $p$ 成立一定能推出 $q$ 成立，充分性成立，}
\step{取 $a=2$，$b=9$，$a+b=11>8$，$ab=18>16$，满足 $q$；}
\step{但是 $a=2<4$，不满足 $p$，所以由 $q$ 不能推出 $p$，必要性不成立，}
\step{综上所述，$p$ 是 $q$ 的充分不必要条件.}
\solend
\fi

\item 若关于 $x$ 的不等式 $ax^2-ax+1>0$ 对一切实数 $x$ 恒成立，则实数 $a$ 的取值范围是\blankline[4.4em].

\ifanswer
\solbegin
\step{当 $a=0$ 时，不等式即 $1>0$，满足条件.}
\step{当 $a\neq0$ 时，要使不等式 $ax^2+ax+1>0$ 对一切 $x\in\R$ 恒成立，}
\step{需 $\begin{cases}a>0\\\Delta=a^2-4a<0\end{cases}$，解得 $0<a<4$.}
\step{综上，实数 $a$ 的取值范围是 $[0,4)$.}
\solend
\fi

\item 已知集合 $M=\{x\mid (x-a)(x^2-ax+a-1)=0\}$ 各元素之和等于 3，则实数 $a=$\blankline[4.4em].

\ifanswer
\solbegin
\step{方程 $(x-a)(x^2-ax+a-1)=0$，化简得 $(x-a)(x-1)[x-(a-1)]=0$，}
\step{解得 $x_1=a$，$x_2=1$，$x_3=a-1$，}
\step{因为它们是集合 $M$ 中的元素，所以存在互异性.}
\step{当 $x_1\neq x_2\neq x_3$ 时，即 $a+1+a-1=3$，解得 $a=\dfrac32$，此时 $M=\dfrac32,1,\dfrac12$，满足.}
\step{当 $a=1$ 时，$1+0\neq3$，不合题意；}
\step{当 $a=a-1$ 时，不符合题意；}
\step{当 $1=a-1$ 时，$a=2$，此时 $M=\{2,1\}$ 满足题意.}
\step{故 $a$ 的取值为 $2$ 或 $\dfrac32$.}
\solend
\fi

\item 已知 $a>b>1$ 且 $a^2-6a-2m+1=0$，$b^2-6b-2m+1=0$，则实数 $m$ 的取值范围是\blankline[4.4em].

\ifanswer
\solbegin
\step{由题意得 $a$、$b$ 是函数 $f(x)=x^2-6x-2m+1$ 两个大于 1 的不同零点，}
\step{所以 $\begin{cases}f(1)>0\\f(3)<0\end{cases}$，即 $\begin{cases}-4-2m>0\\-8-2m<0\end{cases}$，解得 $-4<m<-2$，}
\step{故 $m$ 的取值范围为 $(-4,-2)$.}
\solend
\fi

\end{enumerate}

\bigsec{二、选择题}

\begin{enumerate}
\setcounter{enumi}{12}

\item “$x<2$” 是 “$x^2<4$” 的\mbox{（\quad）}条件

\keepwithprev
A. 充分不必要\qquad\qquad B. 必要不充分

C. 充要\qquad\qquad\qquad\quad D. 既不充分也不必要

\ifanswer
\solbegin
\step{解 $x^2<4$，得 $-2<x<2$，}
\step{由 $x<2$ 得不出 $-2<x<2$，而由 $-2<x<2$ 能得出 $x<2$，}
\step{所以 “$x<2$” 是 “$x^2<4$” 的必要不充分条件. 故选 $B$.}
\solend
\fi

\item 已知集合 $A=\{x\mid x<-2\}$，$B=\{x\mid x\geqslant a\}$，且 $A\cup B=\R$，则实数 $a$ 的取值范围是\mbox{（\quad）}

\keepwithprev
A. $a\leqslant-2$\qquad\qquad B. $a<-2$\qquad\qquad C. $a\geqslant-2$\qquad\qquad D. $a>-2$

\ifanswer
\solbegin
\step{集合 $A=\{x\mid x<-2\}$，$B=\{x\mid x\geqslant a\}$，且 $A\cup B=\R$，}
\step{则实数 $a$ 的取值范围是 $a\leqslant-2$. 故选 $A$.}
\solend
\fi

\item 下列关于 $x,y$ 的方程组 $\begin{cases}y=kx+1\\y=2kx+3\end{cases}$ 的解的说法中正确的是\mbox{（\quad）}

\keepwithprev
A. 该方程一定有唯一解\qquad\qquad B. 该方程没有解

C. $k\neq0$ 时，方程组有唯一解\qquad\qquad D. $k=0$ 时，方程组有无数组解

\ifanswer
\solbegin
\step{由题意得 $kx+1=2kx+3$，即 $kx=-2$，}
\step{当 $k=0$ 时，$0=-2$ 不成立，方程组无解；}
\step{当 $k\neq0$ 时，$x=-\dfrac2k$，方程组有唯一解. 故选 $C$.}
\solend
\fi

\item 设非空集合 $S=\{x\mid m\leqslant x\leqslant n\}$ 满足：当 $x\in S$ 时，有 $x^2\in S$. 给出如下三个结论：①若 $m=1$，则 $S=\{1\}$；②若 $m=-\dfrac12$，则 $\dfrac14\leqslant n\leqslant1$；③若 $n=\dfrac12$，则 $-\dfrac{\sqrt2}{2}\leqslant m\leqslant0$. 其中正确命题的个数为\mbox{（\quad）}

\keepwithprev
A. $3$\qquad\qquad B. $2$\qquad\qquad C. $1$\qquad\qquad D. $0$

\ifanswer
\solbegin
\step{由定义设非空集合 $S=\{x\mid m\leqslant x\leqslant n\}$ 满足当 $x\in S$ 时，有 $x^2\in S$，}
\step{符合定义的参数 $m$ 的值一定大于等于 1 或小于等于 0，}
\step{惟如此才能保证 $m\in S$ 时，有 $m^2\in S$ 即 $m^2\geqslant m$，}
\step{符合条件的 $n$ 的值一定大于等于 0，小于等于 1，}
\step{惟如此才能保证 $n\in S$ 时，有 $n^2\in S$ 即 $n^2\leqslant n$，对各个命题进行判断：}
\step{对于① $m=1$，$m^2=1\in S$，故必有 $\begin{cases}n^2\leqslant n\\n\geqslant1\end{cases}$，得 $n=1$，$S=\{1\}$，}
\step{② $m=-\dfrac12$，$m^2=\dfrac14\in S$，则 $\begin{cases}n^2\leqslant n\\\dfrac14\leqslant n\end{cases}$，解得 $\dfrac14\leqslant n\leqslant1$；}
\step{对于③，若 $n=\dfrac12$，则 $\begin{cases}\dfrac12\geqslant m\\m^2\geqslant m\\\dfrac12\geqslant m^2\end{cases}$，解得 $-\dfrac{\sqrt2}{2}\leqslant m\leqslant0$，}
\step{所以正确命题有 3 个. 故选 $A$.}
\solend
\fi

\end{enumerate}

\bigsec{三、解答题}

\begin{enumerate}
\setcounter{enumi}{16}

\item 已知 $a$、$b$ 是实数.

（1）求证：$a^2+b^2\geqslant2a-2b-2$，并指出等号成立的条件；

（2）若 $ab=1$，求 $a^2+4b^2$ 的最小值.

\ifanswer
\solbegin
\step{（1）因为 $a^2+b^2-(2a-2b-2)=a^2+b^2-2a+2b+2$\dots\dots2 分}
\step{$=(a-1)^2+(b+1)^2\geqslant0$\dots\dots4 分}
\step{所以 $a^2+b^2\geqslant2a-2b-2$.}
\step{当且仅当 $a=1,b=-1$ 时，不等式中等号成立.\dots\dots7 分}
\step{（2）$a^2+4b^2=a^2+(2b)^2\geqslant2\cdot a\cdot(2b)=4ab=4$，\dots\dots5 分}
\step{当且仅当 $a=2b$，即 $\begin{cases}a=\sqrt2\\b=\dfrac{\sqrt2}{2}\end{cases}$ 或 $\begin{cases}a=-\sqrt2\\b=-\dfrac{\sqrt2}{2}\end{cases}$ 时，不等式中等号成立，}
\step{所以 $a^2+4b^2$ 的最小值为 4.\dots\dots7 分}
\solend
\fi

\item 已知关于 $x$ 的方程 $x^2+(m-3)x+2=0$，其中 $m\in\R$，有两个实数根 $\alpha$、$\beta$.

（1）若 $m=8$，求 $\dfrac1\alpha+\dfrac1\beta$ 的值；

（2）若 $\alpha^2+\beta^2=\alpha\beta+3$，求实数 $m$ 的值.

\ifanswer
\solbegin
\step{（1）若 $m=8$，则关于 $x$ 的方程为 $x^2+5x+2=0$，}
\step{则 $\alpha+\beta=-5$，$\alpha\beta=2$，所以 $\dfrac1\alpha+\dfrac1\beta=\dfrac{\alpha+\beta}{\alpha\beta}=-\dfrac52$；}
\step{（2）由题意得 $\alpha+\beta=-(m-3)=3-m$，$\alpha\beta=2$，}
\step{又 $\alpha^2+\beta^2=\alpha\beta+3$，则 $(\alpha+\beta)^2-2\alpha\beta=\alpha\beta+3$，}
\step{即 $(\alpha+\beta)^2=3\alpha\beta+3$，所以 $(3-m)^2=3\times2+3=9$，解得 $m=0$ 或 $6$.}
\solend
\fi

\item 设集合 $A=\{x\mid\lvert x-a\rvert<2\}$，$B=\left\{x\mid\dfrac{x+4}{x+2}>2\right\}$.

（1）若 $a=1$，试用区间表示集合 $A$、$B$；

（2）若 $A\cup\overline{B}=\R$，求实数 $a$ 的取值范围.

\ifanswer
\solbegin
\step{（1）当 $a=1$ 时，由 $\lvert x-1\rvert<2$ 得 $-2<x-1<2$，解得 $-1<x<3$，}
\step{所以 $A=(-1,3)$.}
\step{由 $\dfrac{x+4}{x+2}>2$ 得 $2-\dfrac{x+4}{x+2}=\dfrac{x}{x+2}<0$，则有 $x(x+2)<0$，解得 $-2<x<0$，}
\step{所以 $B=(-2,0)$；}
\step{（2）由 $\lvert x-a\rvert<2$ 得 $-2<x-a<2$，解得 $a-2<x<a+2$，}
\step{所以 $A=(a-2,a+2)$，}
\step{由（1）得 $B=(-2,0)$，则 $\overline{B}=\{x\mid x\geqslant0$ 或 $x\leqslant-2\}$，}
\step{由于 $A\cup\overline{B}=\R$，所以 $\begin{cases}a+2\geqslant0\\a-2\leqslant-2\end{cases}$，解得 $-2\leqslant a\leqslant0$，}
\step{所以实数 $a$ 的取值范围是 $[-2,0]$.}
\solend
\fi

\item 某服装公司生产的衬衫每件定价为 80 元，在某城市年销售 8 万件. 现该公司计划在该市招收代理商来销售衬衫，以降低管理和营销成本. 已知代理商收取的代理费为总销售金额的 $r\%$（其中 $0<r\leqslant10$）.

（1）若单件衬衫的利润保持不变，求招代理商后每件衬衫的定价为多少元（用 $r$ 表示）；

（2）已知将每件衬衫的价格提高到 $\dfrac{80}{1-r\%}$ 元后，每年的销售量会减少 $0.62r$ 万件. 若代理商每年收取的代理费不少于 16 万元，求 $r$ 的取值范围.

\ifanswer
\solbegin
\step{（1）设招代理商后每件衬衫的定价为 $x$ 元，代理商收取总销售金额的 $r\%$，}
\step{因此公司每件衬衫的实际净收入为 $x(1-r\%)$ 元，}
\step{单件利润保持不变，等价于净收入仍为 80 元，}
\step{于是 $x(1-r\%)=80$，解得 $x=\dfrac{80}{1-r\%}$，因为 $0<r\leqslant10$，有 $1-r\%\geqslant0.9>0$，}
\step{分母恒不为零，该定价有意义，所以招代理商后每件衬衫的定价为 $\dfrac{80}{1-r\%}$ 元；}
\step{（2）价格提高到 $\dfrac{80}{1-r\%}$ 元，年销售量减少 $0.62r$ 万件，}
\step{故年销售量为 $8-0.62r$ 万件，$8-0.62r>0$，}
\step{由题意得 $\dfrac{80}{1-r\%}\cdot(8-0.62r)\cdot r\%\geqslant16$，}
\step{即 $\dfrac{80}{1-\dfrac r{100}}\cdot(8-0.62r)\cdot\dfrac r{100}=\dfrac{80r(8-0.62r)}{100-r}\geqslant16$，}
\step{由 $0<r\leqslant10$ 得 $100-r>0$，整理得 $80r(8-0.62r)\geqslant16(100-r)$，}
\step{整理得 $31r^2-410r+1000\leqslant0$，解得 $\dfrac{100}{31}\leqslant r\leqslant10$，}
\step{当 $\dfrac{100}{31}\leqslant r\leqslant10$ 时，$8-0.62r>0$ 且 $1-r\%>0$，}
\step{均符合实际经营条件，所以 $r$ 的取值范围是 $\left[\dfrac{100}{31},10\right]$.}
\solend
\fi

\item 定义 $\min\{a_1,a_2,\cdots,a_n\}$ 为 $n$ 个实数 $a_1,a_2,\cdots,a_n$ 中的最小数，$\max\{a_1,a_2,\cdots,a_n\}$ 为 $n$ 个实数 $a_1,a_2,\cdots,a_n$ 中的最大数.

（1）设 $a,b$ 都是正实数，且 $a+b=1$，求 $\max\left\{ab,\dfrac14\right\}$；

（2）解不等式：$\min\{x+1,x^2+3,\lvert x-1\rvert\}>2x-3$；

（3）设 $a,b$ 都是正实数，求 $\max\left\{a+\dfrac1b,\dfrac2a+b\right\}$ 的最小值.

\ifanswer
\solbegin
\step{（1）$ab\leqslant\dfrac{(a+b)^2}{4}=\dfrac14$，所以 $\max\left\{ab,\dfrac14\right\}=\dfrac14$.}
\step{（2）因为 $x^2+3>x+1$，$x^2+3>\lvert x-1\rvert$，}
\step{所以 $\min\{x+1,x^2+3,\lvert x-1\rvert\}=\min\{x+1,\lvert x-1\rvert\}$.}
\step{当 $x\geqslant0$ 时，$\lvert x-1\rvert\leqslant x+1$，所以 $\lvert x-1\rvert>2x-3$，解得 $0\leqslant x<2$；}
\step{当 $x<0$ 时，$\lvert x-1\rvert>x+1$，所以 $x+1>2x-3$，解得 $x<0$.}
\step{综上，不等式的解集为 $(-\infty,2)$.}
\step{（3）设 $M=\max\left\{a+\dfrac1b,\dfrac2a+b\right\}$，则 $M\geqslant a+\dfrac1b$，$M\geqslant\dfrac2a+b$，}
\step{所以 $M^2\geqslant\left(a+\dfrac1b\right)\left(\dfrac2a+b\right)=3+ab+\dfrac2{ab}\geqslant3+2\sqrt2$，}
\step{所以 $M$ 的最小值为 $\sqrt2+1$，当 $a=\sqrt2$，$b=1$ 时取等号.}
\solend
\fi

\end{enumerate}

\end{document}
"
%
% 原始材料：微信公众号"魔都数学加油站"文章，作者破晓，2026-11-11 18:35 发布；
% 连载"27学年高一126"，原文标题"27学年高一126——2026-2027年上海市建青实验高一上期中"；
% 原文链接：https://mp.weixin.qq.com/s/HpxpBJu4d_wOnUCIbZSQiA。
% 正文为 8 张 PNG 页（01、05、08 为 2560x3620，02、03、04、06、07 为 4762x6735），
% 题目下印有【解析】或【答案】，为讲评版。题目、解析均据原图转录；
% 卷面粉红色水印"魔都数学加油站"及公众号页脚未录入。全卷无插图。
%
% 卷首标题：原卷印作"2026-2027 年建青实验高一上期中"（公众号标题中的"上海市"
% 卷面标题行没有，本稿照卷面），年份区间按全稿惯例排 en dash；原卷未印副标题，
% 本卷为期中卷、跨集合与常用逻辑用语／不等式／函数三章，按题目内容标为
% "集合与常用逻辑用语、不等式、函数"（与 125 同为期中卷，口径相同）：
% 集合与逻辑见 Q5、Q8、Q9、Q11、Q13、Q14、Q16、Q19，不等式见 Q6、Q7、Q10、
% Q17、Q20、Q21，函数（含指数对数、二次函数零点）见 Q3、Q4、Q12、Q15、Q18。
%
% 与原卷的排版差异：
%   ① 原文从第 3 题起，填空题编号为 3—12、选择题为 13—16、解答题为 17—21，本稿保留原题号；
%   ② 原卷每题下直接印【解析】／【答案】，本稿按全稿惯例置于答案版；第 4、5 题只印
%      【答案】，另提 \ans{}；第 13—16 题答案均在【解析】末句"故选 …"，题干括号内留空、
%      不另提 \ans{}；第 3 题印一行【解析】，按 \solbegin+\step 排；第 17 题解析保留
%      原卷评分点（照 高一54／64 成例排 \dots\dots N 分）；
%   ③ 原卷未印副标题，本稿按"知识模块"口径补出（见上"卷首标题"）；
%   ④ 原卷填空横线按全稿惯例排为 \blankline[4.4em]；选择题括号排 \mbox{（\quad）}、
%      选项段前按惯例加 \keepwithprev；第 14 题原卷括号另起一行，本稿并入题干行末；
%      解答题不留作答区（原卷为讲评版，与 109—125 同）；
%   ⑤ 原卷数集排斜体 $R$，按全稿惯例改 \R；绝对值排 \lvert…\rvert，集合描述竖线排 \mid；
%      补集记号照原卷上划线排 \overline{B}；不等号按全稿惯例排 \leqslant／\geqslant；
%      子题标号原卷一律用全角括号（1）（2）（3），本稿照录（\step 行内亦同，照 高一121 成例）。
%
% 判别式字形：第 10 题解析 $\Delta=a^2-4a<0$ 原卷排 $\Delta$（等腰、高），
%   照录 \Delta（与 高一124／125 同题排 $\triangle$（等边、扁）不同，逐处照本卷字形）。
%
% 子集符号：第 8 题题干与解析两处 $S\subset P$ 均印单线、无下划线，照录 \subset；
%   全卷未出现带下划线的 \subseteq。
%
% 原卷存疑照录：
%   第 10 题题干作 $ax^2-ax+1>0$，解析却作 $ax^2+ax+1>0$（一次项符号不一致），照录、未统一；
%   第 11 题解析"此时 $M=\dfrac32,1,\dfrac12$，满足"未用集合花括号（题干 $M$ 为集合），照录；
%   第 11 题解析"因为它们是集合 $M$ 中的元素，所以存在互异性"（"存在互异性"疑为
%     "满足互异性"之误），又"当 $x_1\neq x_2\neq x_3$ 时"连写两个不等号，均照录；
%   第 16 题解析两处"惟如此才能保证"（用"惟"字）及"符合定义的参数 $m$ 的值一定大于等于 1
%     或小于等于 0"的措辞，照录；
%   第 17 题 (1) 解析"当且仅当 $a=1,b=-1$ 时"用半角逗号、第 21 题 (3) 解析
%     "当 $a=\sqrt2$，$b=1$ 时取等号"用全角逗号，均照录；
%   第 19 题解析由 $\dfrac{x+4}{x+2}>2$ 得 "$2-\dfrac{x+4}{x+2}=\dfrac{x}{x+2}<0$"
%     （以 2 减左式，与常见的"左式减 2"次序相反），照录；
%   第 15 题选项 A、B 作"该方程…"、C、D 作"方程组…"，照录；
%   第 20 题题干作"招收代理商"，(1) 却作"招代理商后"，照录。
%
% 行数复核：逐题数 \step 条数，合计 91 条（Q3：1；Q4：0；Q5：0；Q6：3；Q7：2；
% Q8：5；Q9：5；Q10：4；Q11：8；Q12：3；Q13：3；Q14：2；Q15：3；Q16：9；
% Q17：7；Q18：5；Q19：9；Q20：13；Q21：9），与原卷印刷行逐题相等
% （cases 块作一个印刷行计）。第 4、5 题只印【答案】，提 \ans{}，不计入 \step 账；
% 四道选择题答案均在解析末句，括号内留空、不另提 \ans{}。
%
\documentclass[a4paper,11pt]{article}
\usepackage{common}

\begin{document}

\examtitlest[3]{2026--2027 年建青实验高一上期中}{集合与常用逻辑用语、不等式、函数}

\bigsec{一、填空题}

\begin{enumerate}
\setcounter{enumi}{2}

\item 已知 $\log_2 x=3$，则 $x=$\blankline[4.4em].

\ifanswer
\solbegin
\step{$x=2^3=8$.}
\solend
\fi

\item 将 $\sqrt{a^2\sqrt a}$（其中 $a>0$）化为有理数指数幂的形式为\blankline[4.4em].

\ifanswer
\ans{$a^{\frac54}$}
\fi

\item 用反证法证明命题：“已知 $m$、$n\in\R$，若 $m+n\leqslant5$，则 $m\leqslant2$ 或 $n\leqslant3$” 时应假设\blankline[4.4em].

\ifanswer
\ans{$m>2$ 且 $n>3$}
\fi

\item 若 $x\in\R$，则 $\lvert x+1\rvert+\lvert x-1\rvert$ 的最小值为\blankline[4.4em].

\ifanswer
\solbegin
\step{$\lvert x+1\rvert+\lvert x-1\rvert\geqslant\lvert(x+1)-(x-1)\rvert=2$，}
\step{当且仅当 $(x+1)(x-1)\leqslant0$，即 $-1\leqslant x\leqslant1$ 时取等号，}
\step{所以 $\lvert x+1\rvert+\lvert x-1\rvert$ 的最小值为 2.}
\solend
\fi

\item 已知 $x>1$，则 $x+\dfrac1{x-1}$ 的最小值是\blankline[4.4em].

\ifanswer
\solbegin
\step{$x-1>0$，则 $x+\dfrac1{x-1}=(x-1)+\dfrac1{x-1}+1\geqslant2\sqrt{(x-1)\cdot\dfrac1{x-1}}+1=3$，}
\step{当且仅当 $x=2$ 时取等号，故 $x+\dfrac1{x-1}$ 的最小值为 3.}
\solend
\fi

\item 已知集合 $P=\{x\mid x^2+x-6=0\}$，$S=\{x\mid ax+1=0\}$ 且 $S\subset P$，则 $a$ 的值为\blankline[4.4em].

\ifanswer
\solbegin
\step{集合 $P=\{-3,2\}$，集合 $S$ 中至多有一个元素，}
\step{若集合 $S$ 为空集，即 $a=0$ 时，显然满足条件 $S\subset P$，故 $a=0$.}
\step{若集合 $S$ 非空集，即 $a\neq0$，此时 $S=\left\{-\dfrac1a\right\}$，}
\step{若 $-\dfrac1a=-3$，则 $a=\dfrac13$，若 $-\dfrac1a=2$，则 $a=-\dfrac12$；}
\step{故 $a$ 的取值为 $0$、$\dfrac13$、$-\dfrac12$.}
\solend
\fi

\item 已知 $p:\begin{cases}a>4\\b>4\end{cases}$，$q:\begin{cases}a+b>8\\ab>16\end{cases}$，则 $p$ 是 $q$ 的\blankline[4.4em]条件.

\ifanswer
\solbegin
\step{若 $a>4$，$b>4$，则 $a+b>4+4=8$，$ab>4\times4=16$，}
\step{所以 $p$ 成立一定能推出 $q$ 成立，充分性成立，}
\step{取 $a=2$，$b=9$，$a+b=11>8$，$ab=18>16$，满足 $q$；}
\step{但是 $a=2<4$，不满足 $p$，所以由 $q$ 不能推出 $p$，必要性不成立，}
\step{综上所述，$p$ 是 $q$ 的充分不必要条件.}
\solend
\fi

\item 若关于 $x$ 的不等式 $ax^2-ax+1>0$ 对一切实数 $x$ 恒成立，则实数 $a$ 的取值范围是\blankline[4.4em].

\ifanswer
\solbegin
\step{当 $a=0$ 时，不等式即 $1>0$，满足条件.}
\step{当 $a\neq0$ 时，要使不等式 $ax^2+ax+1>0$ 对一切 $x\in\R$ 恒成立，}
\step{需 $\begin{cases}a>0\\\Delta=a^2-4a<0\end{cases}$，解得 $0<a<4$.}
\step{综上，实数 $a$ 的取值范围是 $[0,4)$.}
\solend
\fi

\item 已知集合 $M=\{x\mid (x-a)(x^2-ax+a-1)=0\}$ 各元素之和等于 3，则实数 $a=$\blankline[4.4em].

\ifanswer
\solbegin
\step{方程 $(x-a)(x^2-ax+a-1)=0$，化简得 $(x-a)(x-1)[x-(a-1)]=0$，}
\step{解得 $x_1=a$，$x_2=1$，$x_3=a-1$，}
\step{因为它们是集合 $M$ 中的元素，所以存在互异性.}
\step{当 $x_1\neq x_2\neq x_3$ 时，即 $a+1+a-1=3$，解得 $a=\dfrac32$，此时 $M=\dfrac32,1,\dfrac12$，满足.}
\step{当 $a=1$ 时，$1+0\neq3$，不合题意；}
\step{当 $a=a-1$ 时，不符合题意；}
\step{当 $1=a-1$ 时，$a=2$，此时 $M=\{2,1\}$ 满足题意.}
\step{故 $a$ 的取值为 $2$ 或 $\dfrac32$.}
\solend
\fi

\item 已知 $a>b>1$ 且 $a^2-6a-2m+1=0$，$b^2-6b-2m+1=0$，则实数 $m$ 的取值范围是\blankline[4.4em].

\ifanswer
\solbegin
\step{由题意得 $a$、$b$ 是函数 $f(x)=x^2-6x-2m+1$ 两个大于 1 的不同零点，}
\step{所以 $\begin{cases}f(1)>0\\f(3)<0\end{cases}$，即 $\begin{cases}-4-2m>0\\-8-2m<0\end{cases}$，解得 $-4<m<-2$，}
\step{故 $m$ 的取值范围为 $(-4,-2)$.}
\solend
\fi

\end{enumerate}

\bigsec{二、选择题}

\begin{enumerate}
\setcounter{enumi}{12}

\item “$x<2$” 是 “$x^2<4$” 的\mbox{（\quad）}条件

\keepwithprev
A. 充分不必要\qquad\qquad B. 必要不充分

C. 充要\qquad\qquad\qquad\quad D. 既不充分也不必要

\ifanswer
\solbegin
\step{解 $x^2<4$，得 $-2<x<2$，}
\step{由 $x<2$ 得不出 $-2<x<2$，而由 $-2<x<2$ 能得出 $x<2$，}
\step{所以 “$x<2$” 是 “$x^2<4$” 的必要不充分条件. 故选 $B$.}
\solend
\fi

\item 已知集合 $A=\{x\mid x<-2\}$，$B=\{x\mid x\geqslant a\}$，且 $A\cup B=\R$，则实数 $a$ 的取值范围是\mbox{（\quad）}

\keepwithprev
A. $a\leqslant-2$\qquad\qquad B. $a<-2$\qquad\qquad C. $a\geqslant-2$\qquad\qquad D. $a>-2$

\ifanswer
\solbegin
\step{集合 $A=\{x\mid x<-2\}$，$B=\{x\mid x\geqslant a\}$，且 $A\cup B=\R$，}
\step{则实数 $a$ 的取值范围是 $a\leqslant-2$. 故选 $A$.}
\solend
\fi

\item 下列关于 $x,y$ 的方程组 $\begin{cases}y=kx+1\\y=2kx+3\end{cases}$ 的解的说法中正确的是\mbox{（\quad）}

\keepwithprev
A. 该方程一定有唯一解\qquad\qquad B. 该方程没有解

C. $k\neq0$ 时，方程组有唯一解\qquad\qquad D. $k=0$ 时，方程组有无数组解

\ifanswer
\solbegin
\step{由题意得 $kx+1=2kx+3$，即 $kx=-2$，}
\step{当 $k=0$ 时，$0=-2$ 不成立，方程组无解；}
\step{当 $k\neq0$ 时，$x=-\dfrac2k$，方程组有唯一解. 故选 $C$.}
\solend
\fi

\item 设非空集合 $S=\{x\mid m\leqslant x\leqslant n\}$ 满足：当 $x\in S$ 时，有 $x^2\in S$. 给出如下三个结论：①若 $m=1$，则 $S=\{1\}$；②若 $m=-\dfrac12$，则 $\dfrac14\leqslant n\leqslant1$；③若 $n=\dfrac12$，则 $-\dfrac{\sqrt2}{2}\leqslant m\leqslant0$. 其中正确命题的个数为\mbox{（\quad）}

\keepwithprev
A. $3$\qquad\qquad B. $2$\qquad\qquad C. $1$\qquad\qquad D. $0$

\ifanswer
\solbegin
\step{由定义设非空集合 $S=\{x\mid m\leqslant x\leqslant n\}$ 满足当 $x\in S$ 时，有 $x^2\in S$，}
\step{符合定义的参数 $m$ 的值一定大于等于 1 或小于等于 0，}
\step{惟如此才能保证 $m\in S$ 时，有 $m^2\in S$ 即 $m^2\geqslant m$，}
\step{符合条件的 $n$ 的值一定大于等于 0，小于等于 1，}
\step{惟如此才能保证 $n\in S$ 时，有 $n^2\in S$ 即 $n^2\leqslant n$，对各个命题进行判断：}
\step{对于① $m=1$，$m^2=1\in S$，故必有 $\begin{cases}n^2\leqslant n\\n\geqslant1\end{cases}$，得 $n=1$，$S=\{1\}$，}
\step{② $m=-\dfrac12$，$m^2=\dfrac14\in S$，则 $\begin{cases}n^2\leqslant n\\\dfrac14\leqslant n\end{cases}$，解得 $\dfrac14\leqslant n\leqslant1$；}
\step{对于③，若 $n=\dfrac12$，则 $\begin{cases}\dfrac12\geqslant m\\m^2\geqslant m\\\dfrac12\geqslant m^2\end{cases}$，解得 $-\dfrac{\sqrt2}{2}\leqslant m\leqslant0$，}
\step{所以正确命题有 3 个. 故选 $A$.}
\solend
\fi

\end{enumerate}

\bigsec{三、解答题}

\begin{enumerate}
\setcounter{enumi}{16}

\item 已知 $a$、$b$ 是实数.

（1）求证：$a^2+b^2\geqslant2a-2b-2$，并指出等号成立的条件；

（2）若 $ab=1$，求 $a^2+4b^2$ 的最小值.

\ifanswer
\solbegin
\step{（1）因为 $a^2+b^2-(2a-2b-2)=a^2+b^2-2a+2b+2$\dots\dots2 分}
\step{$=(a-1)^2+(b+1)^2\geqslant0$\dots\dots4 分}
\step{所以 $a^2+b^2\geqslant2a-2b-2$.}
\step{当且仅当 $a=1,b=-1$ 时，不等式中等号成立.\dots\dots7 分}
\step{（2）$a^2+4b^2=a^2+(2b)^2\geqslant2\cdot a\cdot(2b)=4ab=4$，\dots\dots5 分}
\step{当且仅当 $a=2b$，即 $\begin{cases}a=\sqrt2\\b=\dfrac{\sqrt2}{2}\end{cases}$ 或 $\begin{cases}a=-\sqrt2\\b=-\dfrac{\sqrt2}{2}\end{cases}$ 时，不等式中等号成立，}
\step{所以 $a^2+4b^2$ 的最小值为 4.\dots\dots7 分}
\solend
\fi

\item 已知关于 $x$ 的方程 $x^2+(m-3)x+2=0$，其中 $m\in\R$，有两个实数根 $\alpha$、$\beta$.

（1）若 $m=8$，求 $\dfrac1\alpha+\dfrac1\beta$ 的值；

（2）若 $\alpha^2+\beta^2=\alpha\beta+3$，求实数 $m$ 的值.

\ifanswer
\solbegin
\step{（1）若 $m=8$，则关于 $x$ 的方程为 $x^2+5x+2=0$，}
\step{则 $\alpha+\beta=-5$，$\alpha\beta=2$，所以 $\dfrac1\alpha+\dfrac1\beta=\dfrac{\alpha+\beta}{\alpha\beta}=-\dfrac52$；}
\step{（2）由题意得 $\alpha+\beta=-(m-3)=3-m$，$\alpha\beta=2$，}
\step{又 $\alpha^2+\beta^2=\alpha\beta+3$，则 $(\alpha+\beta)^2-2\alpha\beta=\alpha\beta+3$，}
\step{即 $(\alpha+\beta)^2=3\alpha\beta+3$，所以 $(3-m)^2=3\times2+3=9$，解得 $m=0$ 或 $6$.}
\solend
\fi

\item 设集合 $A=\{x\mid\lvert x-a\rvert<2\}$，$B=\left\{x\mid\dfrac{x+4}{x+2}>2\right\}$.

（1）若 $a=1$，试用区间表示集合 $A$、$B$；

（2）若 $A\cup\overline{B}=\R$，求实数 $a$ 的取值范围.

\ifanswer
\solbegin
\step{（1）当 $a=1$ 时，由 $\lvert x-1\rvert<2$ 得 $-2<x-1<2$，解得 $-1<x<3$，}
\step{所以 $A=(-1,3)$.}
\step{由 $\dfrac{x+4}{x+2}>2$ 得 $2-\dfrac{x+4}{x+2}=\dfrac{x}{x+2}<0$，则有 $x(x+2)<0$，解得 $-2<x<0$，}
\step{所以 $B=(-2,0)$；}
\step{（2）由 $\lvert x-a\rvert<2$ 得 $-2<x-a<2$，解得 $a-2<x<a+2$，}
\step{所以 $A=(a-2,a+2)$，}
\step{由（1）得 $B=(-2,0)$，则 $\overline{B}=\{x\mid x\geqslant0$ 或 $x\leqslant-2\}$，}
\step{由于 $A\cup\overline{B}=\R$，所以 $\begin{cases}a+2\geqslant0\\a-2\leqslant-2\end{cases}$，解得 $-2\leqslant a\leqslant0$，}
\step{所以实数 $a$ 的取值范围是 $[-2,0]$.}
\solend
\fi

\item 某服装公司生产的衬衫每件定价为 80 元，在某城市年销售 8 万件. 现该公司计划在该市招收代理商来销售衬衫，以降低管理和营销成本. 已知代理商收取的代理费为总销售金额的 $r\%$（其中 $0<r\leqslant10$）.

（1）若单件衬衫的利润保持不变，求招代理商后每件衬衫的定价为多少元（用 $r$ 表示）；

（2）已知将每件衬衫的价格提高到 $\dfrac{80}{1-r\%}$ 元后，每年的销售量会减少 $0.62r$ 万件. 若代理商每年收取的代理费不少于 16 万元，求 $r$ 的取值范围.

\ifanswer
\solbegin
\step{（1）设招代理商后每件衬衫的定价为 $x$ 元，代理商收取总销售金额的 $r\%$，}
\step{因此公司每件衬衫的实际净收入为 $x(1-r\%)$ 元，}
\step{单件利润保持不变，等价于净收入仍为 80 元，}
\step{于是 $x(1-r\%)=80$，解得 $x=\dfrac{80}{1-r\%}$，因为 $0<r\leqslant10$，有 $1-r\%\geqslant0.9>0$，}
\step{分母恒不为零，该定价有意义，所以招代理商后每件衬衫的定价为 $\dfrac{80}{1-r\%}$ 元；}
\step{（2）价格提高到 $\dfrac{80}{1-r\%}$ 元，年销售量减少 $0.62r$ 万件，}
\step{故年销售量为 $8-0.62r$ 万件，$8-0.62r>0$，}
\step{由题意得 $\dfrac{80}{1-r\%}\cdot(8-0.62r)\cdot r\%\geqslant16$，}
\step{即 $\dfrac{80}{1-\dfrac r{100}}\cdot(8-0.62r)\cdot\dfrac r{100}=\dfrac{80r(8-0.62r)}{100-r}\geqslant16$，}
\step{由 $0<r\leqslant10$ 得 $100-r>0$，整理得 $80r(8-0.62r)\geqslant16(100-r)$，}
\step{整理得 $31r^2-410r+1000\leqslant0$，解得 $\dfrac{100}{31}\leqslant r\leqslant10$，}
\step{当 $\dfrac{100}{31}\leqslant r\leqslant10$ 时，$8-0.62r>0$ 且 $1-r\%>0$，}
\step{均符合实际经营条件，所以 $r$ 的取值范围是 $\left[\dfrac{100}{31},10\right]$.}
\solend
\fi

\item 定义 $\min\{a_1,a_2,\cdots,a_n\}$ 为 $n$ 个实数 $a_1,a_2,\cdots,a_n$ 中的最小数，$\max\{a_1,a_2,\cdots,a_n\}$ 为 $n$ 个实数 $a_1,a_2,\cdots,a_n$ 中的最大数.

（1）设 $a,b$ 都是正实数，且 $a+b=1$，求 $\max\left\{ab,\dfrac14\right\}$；

（2）解不等式：$\min\{x+1,x^2+3,\lvert x-1\rvert\}>2x-3$；

（3）设 $a,b$ 都是正实数，求 $\max\left\{a+\dfrac1b,\dfrac2a+b\right\}$ 的最小值.

\ifanswer
\solbegin
\step{（1）$ab\leqslant\dfrac{(a+b)^2}{4}=\dfrac14$，所以 $\max\left\{ab,\dfrac14\right\}=\dfrac14$.}
\step{（2）因为 $x^2+3>x+1$，$x^2+3>\lvert x-1\rvert$，}
\step{所以 $\min\{x+1,x^2+3,\lvert x-1\rvert\}=\min\{x+1,\lvert x-1\rvert\}$.}
\step{当 $x\geqslant0$ 时，$\lvert x-1\rvert\leqslant x+1$，所以 $\lvert x-1\rvert>2x-3$，解得 $0\leqslant x<2$；}
\step{当 $x<0$ 时，$\lvert x-1\rvert>x+1$，所以 $x+1>2x-3$，解得 $x<0$.}
\step{综上，不等式的解集为 $(-\infty,2)$.}
\step{（3）设 $M=\max\left\{a+\dfrac1b,\dfrac2a+b\right\}$，则 $M\geqslant a+\dfrac1b$，$M\geqslant\dfrac2a+b$，}
\step{所以 $M^2\geqslant\left(a+\dfrac1b\right)\left(\dfrac2a+b\right)=3+ab+\dfrac2{ab}\geqslant3+2\sqrt2$，}
\step{所以 $M$ 的最小值为 $\sqrt2+1$，当 $a=\sqrt2$，$b=1$ 时取等号.}
\solend
\fi

\end{enumerate}

\end{document}
"
%
% 原始材料：微信公众号"魔都数学加油站"文章，作者破晓，2026-11-11 18:35 发布；
% 连载"27学年高一126"，原文标题"27学年高一126——2026-2027年上海市建青实验高一上期中"；
% 原文链接：https://mp.weixin.qq.com/s/HpxpBJu4d_wOnUCIbZSQiA。
% 正文为 8 张 PNG 页（01、05、08 为 2560x3620，02、03、04、06、07 为 4762x6735），
% 题目下印有【解析】或【答案】，为讲评版。题目、解析均据原图转录；
% 卷面粉红色水印"魔都数学加油站"及公众号页脚未录入。全卷无插图。
%
% 卷首标题：原卷印作"2026-2027 年建青实验高一上期中"（公众号标题中的"上海市"
% 卷面标题行没有，本稿照卷面），年份区间按全稿惯例排 en dash；原卷未印副标题，
% 本卷为期中卷、跨集合与常用逻辑用语／不等式／函数三章，按题目内容标为
% "集合与常用逻辑用语、不等式、函数"（与 125 同为期中卷，口径相同）：
% 集合与逻辑见 Q5、Q8、Q9、Q11、Q13、Q14、Q16、Q19，不等式见 Q6、Q7、Q10、
% Q17、Q20、Q21，函数（含指数对数、二次函数零点）见 Q3、Q4、Q12、Q15、Q18。
%
% 与原卷的排版差异：
%   ① 原文从第 3 题起，填空题编号为 3—12、选择题为 13—16、解答题为 17—21，本稿保留原题号；
%   ② 原卷每题下直接印【解析】／【答案】，本稿按全稿惯例置于答案版；第 4、5 题只印
%      【答案】，另提 \ans{}；第 13—16 题答案均在【解析】末句"故选 …"，题干括号内留空、
%      不另提 \ans{}；第 3 题印一行【解析】，按 \solbegin+\step 排；第 17 题解析保留
%      原卷评分点（照 高一54／64 成例排 \dots\dots N 分）；
%   ③ 原卷未印副标题，本稿按"知识模块"口径补出（见上"卷首标题"）；
%   ④ 原卷填空横线按全稿惯例排为 \blankline[4.4em]；选择题括号排 \mbox{（\quad）}、
%      选项段前按惯例加 \keepwithprev；第 14 题原卷括号另起一行，本稿并入题干行末；
%      解答题不留作答区（原卷为讲评版，与 109—125 同）；
%   ⑤ 原卷数集排斜体 $R$，按全稿惯例改 \R；绝对值排 \lvert…\rvert，集合描述竖线排 \mid；
%      补集记号照原卷上划线排 \overline{B}；不等号按全稿惯例排 \leqslant／\geqslant；
%      子题标号原卷一律用全角括号（1）（2）（3），本稿照录（\step 行内亦同，照 高一121 成例）。
%
% 判别式字形：第 10 题解析 $\Delta=a^2-4a<0$ 原卷排 $\Delta$（等腰、高），
%   照录 \Delta（与 高一124／125 同题排 $\triangle$（等边、扁）不同，逐处照本卷字形）。
%
% 子集符号：第 8 题题干与解析两处 $S\subset P$ 均印单线、无下划线，照录 \subset；
%   全卷未出现带下划线的 \subseteq。
%
% 原卷存疑照录：
%   第 10 题题干作 $ax^2-ax+1>0$，解析却作 $ax^2+ax+1>0$（一次项符号不一致），照录、未统一；
%   第 11 题解析"此时 $M=\dfrac32,1,\dfrac12$，满足"未用集合花括号（题干 $M$ 为集合），照录；
%   第 11 题解析"因为它们是集合 $M$ 中的元素，所以存在互异性"（"存在互异性"疑为
%     "满足互异性"之误），又"当 $x_1\neq x_2\neq x_3$ 时"连写两个不等号，均照录；
%   第 16 题解析两处"惟如此才能保证"（用"惟"字）及"符合定义的参数 $m$ 的值一定大于等于 1
%     或小于等于 0"的措辞，照录；
%   第 17 题 (1) 解析"当且仅当 $a=1,b=-1$ 时"用半角逗号、第 21 题 (3) 解析
%     "当 $a=\sqrt2$，$b=1$ 时取等号"用全角逗号，均照录；
%   第 19 题解析由 $\dfrac{x+4}{x+2}>2$ 得 "$2-\dfrac{x+4}{x+2}=\dfrac{x}{x+2}<0$"
%     （以 2 减左式，与常见的"左式减 2"次序相反），照录；
%   第 15 题选项 A、B 作"该方程…"、C、D 作"方程组…"，照录；
%   第 20 题题干作"招收代理商"，(1) 却作"招代理商后"，照录。
%
% 行数复核：逐题数 \step 条数，合计 91 条（Q3：1；Q4：0；Q5：0；Q6：3；Q7：2；
% Q8：5；Q9：5；Q10：4；Q11：8；Q12：3；Q13：3；Q14：2；Q15：3；Q16：9；
% Q17：7；Q18：5；Q19：9；Q20：13；Q21：9），与原卷印刷行逐题相等
% （cases 块作一个印刷行计）。第 4、5 题只印【答案】，提 \ans{}，不计入 \step 账；
% 四道选择题答案均在解析末句，括号内留空、不另提 \ans{}。
%
\documentclass[a4paper,11pt]{article}
\usepackage{common}

\begin{document}

\examtitlest[3]{2026--2027 年建青实验高一上期中}{集合与常用逻辑用语、不等式、函数}

\bigsec{一、填空题}

\begin{enumerate}
\setcounter{enumi}{2}

\item 已知 $\log_2 x=3$，则 $x=$\blankline[4.4em].

\ifanswer
\solbegin
\step{$x=2^3=8$.}
\solend
\fi

\item 将 $\sqrt{a^2\sqrt a}$（其中 $a>0$）化为有理数指数幂的形式为\blankline[4.4em].

\ifanswer
\ans{$a^{\frac54}$}
\fi

\item 用反证法证明命题：“已知 $m$、$n\in\R$，若 $m+n\leqslant5$，则 $m\leqslant2$ 或 $n\leqslant3$” 时应假设\blankline[4.4em].

\ifanswer
\ans{$m>2$ 且 $n>3$}
\fi

\item 若 $x\in\R$，则 $\lvert x+1\rvert+\lvert x-1\rvert$ 的最小值为\blankline[4.4em].

\ifanswer
\solbegin
\step{$\lvert x+1\rvert+\lvert x-1\rvert\geqslant\lvert(x+1)-(x-1)\rvert=2$，}
\step{当且仅当 $(x+1)(x-1)\leqslant0$，即 $-1\leqslant x\leqslant1$ 时取等号，}
\step{所以 $\lvert x+1\rvert+\lvert x-1\rvert$ 的最小值为 2.}
\solend
\fi

\item 已知 $x>1$，则 $x+\dfrac1{x-1}$ 的最小值是\blankline[4.4em].

\ifanswer
\solbegin
\step{$x-1>0$，则 $x+\dfrac1{x-1}=(x-1)+\dfrac1{x-1}+1\geqslant2\sqrt{(x-1)\cdot\dfrac1{x-1}}+1=3$，}
\step{当且仅当 $x=2$ 时取等号，故 $x+\dfrac1{x-1}$ 的最小值为 3.}
\solend
\fi

\item 已知集合 $P=\{x\mid x^2+x-6=0\}$，$S=\{x\mid ax+1=0\}$ 且 $S\subset P$，则 $a$ 的值为\blankline[4.4em].

\ifanswer
\solbegin
\step{集合 $P=\{-3,2\}$，集合 $S$ 中至多有一个元素，}
\step{若集合 $S$ 为空集，即 $a=0$ 时，显然满足条件 $S\subset P$，故 $a=0$.}
\step{若集合 $S$ 非空集，即 $a\neq0$，此时 $S=\left\{-\dfrac1a\right\}$，}
\step{若 $-\dfrac1a=-3$，则 $a=\dfrac13$，若 $-\dfrac1a=2$，则 $a=-\dfrac12$；}
\step{故 $a$ 的取值为 $0$、$\dfrac13$、$-\dfrac12$.}
\solend
\fi

\item 已知 $p:\begin{cases}a>4\\b>4\end{cases}$，$q:\begin{cases}a+b>8\\ab>16\end{cases}$，则 $p$ 是 $q$ 的\blankline[4.4em]条件.

\ifanswer
\solbegin
\step{若 $a>4$，$b>4$，则 $a+b>4+4=8$，$ab>4\times4=16$，}
\step{所以 $p$ 成立一定能推出 $q$ 成立，充分性成立，}
\step{取 $a=2$，$b=9$，$a+b=11>8$，$ab=18>16$，满足 $q$；}
\step{但是 $a=2<4$，不满足 $p$，所以由 $q$ 不能推出 $p$，必要性不成立，}
\step{综上所述，$p$ 是 $q$ 的充分不必要条件.}
\solend
\fi

\item 若关于 $x$ 的不等式 $ax^2-ax+1>0$ 对一切实数 $x$ 恒成立，则实数 $a$ 的取值范围是\blankline[4.4em].

\ifanswer
\solbegin
\step{当 $a=0$ 时，不等式即 $1>0$，满足条件.}
\step{当 $a\neq0$ 时，要使不等式 $ax^2+ax+1>0$ 对一切 $x\in\R$ 恒成立，}
\step{需 $\begin{cases}a>0\\\Delta=a^2-4a<0\end{cases}$，解得 $0<a<4$.}
\step{综上，实数 $a$ 的取值范围是 $[0,4)$.}
\solend
\fi

\item 已知集合 $M=\{x\mid (x-a)(x^2-ax+a-1)=0\}$ 各元素之和等于 3，则实数 $a=$\blankline[4.4em].

\ifanswer
\solbegin
\step{方程 $(x-a)(x^2-ax+a-1)=0$，化简得 $(x-a)(x-1)[x-(a-1)]=0$，}
\step{解得 $x_1=a$，$x_2=1$，$x_3=a-1$，}
\step{因为它们是集合 $M$ 中的元素，所以存在互异性.}
\step{当 $x_1\neq x_2\neq x_3$ 时，即 $a+1+a-1=3$，解得 $a=\dfrac32$，此时 $M=\dfrac32,1,\dfrac12$，满足.}
\step{当 $a=1$ 时，$1+0\neq3$，不合题意；}
\step{当 $a=a-1$ 时，不符合题意；}
\step{当 $1=a-1$ 时，$a=2$，此时 $M=\{2,1\}$ 满足题意.}
\step{故 $a$ 的取值为 $2$ 或 $\dfrac32$.}
\solend
\fi

\item 已知 $a>b>1$ 且 $a^2-6a-2m+1=0$，$b^2-6b-2m+1=0$，则实数 $m$ 的取值范围是\blankline[4.4em].

\ifanswer
\solbegin
\step{由题意得 $a$、$b$ 是函数 $f(x)=x^2-6x-2m+1$ 两个大于 1 的不同零点，}
\step{所以 $\begin{cases}f(1)>0\\f(3)<0\end{cases}$，即 $\begin{cases}-4-2m>0\\-8-2m<0\end{cases}$，解得 $-4<m<-2$，}
\step{故 $m$ 的取值范围为 $(-4,-2)$.}
\solend
\fi

\end{enumerate}

\bigsec{二、选择题}

\begin{enumerate}
\setcounter{enumi}{12}

\item “$x<2$” 是 “$x^2<4$” 的\mbox{（\quad）}条件

\keepwithprev
A. 充分不必要\qquad\qquad B. 必要不充分

C. 充要\qquad\qquad\qquad\quad D. 既不充分也不必要

\ifanswer
\solbegin
\step{解 $x^2<4$，得 $-2<x<2$，}
\step{由 $x<2$ 得不出 $-2<x<2$，而由 $-2<x<2$ 能得出 $x<2$，}
\step{所以 “$x<2$” 是 “$x^2<4$” 的必要不充分条件. 故选 $B$.}
\solend
\fi

\item 已知集合 $A=\{x\mid x<-2\}$，$B=\{x\mid x\geqslant a\}$，且 $A\cup B=\R$，则实数 $a$ 的取值范围是\mbox{（\quad）}

\keepwithprev
A. $a\leqslant-2$\qquad\qquad B. $a<-2$\qquad\qquad C. $a\geqslant-2$\qquad\qquad D. $a>-2$

\ifanswer
\solbegin
\step{集合 $A=\{x\mid x<-2\}$，$B=\{x\mid x\geqslant a\}$，且 $A\cup B=\R$，}
\step{则实数 $a$ 的取值范围是 $a\leqslant-2$. 故选 $A$.}
\solend
\fi

\item 下列关于 $x,y$ 的方程组 $\begin{cases}y=kx+1\\y=2kx+3\end{cases}$ 的解的说法中正确的是\mbox{（\quad）}

\keepwithprev
A. 该方程一定有唯一解\qquad\qquad B. 该方程没有解

C. $k\neq0$ 时，方程组有唯一解\qquad\qquad D. $k=0$ 时，方程组有无数组解

\ifanswer
\solbegin
\step{由题意得 $kx+1=2kx+3$，即 $kx=-2$，}
\step{当 $k=0$ 时，$0=-2$ 不成立，方程组无解；}
\step{当 $k\neq0$ 时，$x=-\dfrac2k$，方程组有唯一解. 故选 $C$.}
\solend
\fi

\item 设非空集合 $S=\{x\mid m\leqslant x\leqslant n\}$ 满足：当 $x\in S$ 时，有 $x^2\in S$. 给出如下三个结论：①若 $m=1$，则 $S=\{1\}$；②若 $m=-\dfrac12$，则 $\dfrac14\leqslant n\leqslant1$；③若 $n=\dfrac12$，则 $-\dfrac{\sqrt2}{2}\leqslant m\leqslant0$. 其中正确命题的个数为\mbox{（\quad）}

\keepwithprev
A. $3$\qquad\qquad B. $2$\qquad\qquad C. $1$\qquad\qquad D. $0$

\ifanswer
\solbegin
\step{由定义设非空集合 $S=\{x\mid m\leqslant x\leqslant n\}$ 满足当 $x\in S$ 时，有 $x^2\in S$，}
\step{符合定义的参数 $m$ 的值一定大于等于 1 或小于等于 0，}
\step{惟如此才能保证 $m\in S$ 时，有 $m^2\in S$ 即 $m^2\geqslant m$，}
\step{符合条件的 $n$ 的值一定大于等于 0，小于等于 1，}
\step{惟如此才能保证 $n\in S$ 时，有 $n^2\in S$ 即 $n^2\leqslant n$，对各个命题进行判断：}
\step{对于① $m=1$，$m^2=1\in S$，故必有 $\begin{cases}n^2\leqslant n\\n\geqslant1\end{cases}$，得 $n=1$，$S=\{1\}$，}
\step{② $m=-\dfrac12$，$m^2=\dfrac14\in S$，则 $\begin{cases}n^2\leqslant n\\\dfrac14\leqslant n\end{cases}$，解得 $\dfrac14\leqslant n\leqslant1$；}
\step{对于③，若 $n=\dfrac12$，则 $\begin{cases}\dfrac12\geqslant m\\m^2\geqslant m\\\dfrac12\geqslant m^2\end{cases}$，解得 $-\dfrac{\sqrt2}{2}\leqslant m\leqslant0$，}
\step{所以正确命题有 3 个. 故选 $A$.}
\solend
\fi

\end{enumerate}

\bigsec{三、解答题}

\begin{enumerate}
\setcounter{enumi}{16}

\item 已知 $a$、$b$ 是实数.

（1）求证：$a^2+b^2\geqslant2a-2b-2$，并指出等号成立的条件；

（2）若 $ab=1$，求 $a^2+4b^2$ 的最小值.

\ifanswer
\solbegin
\step{（1）因为 $a^2+b^2-(2a-2b-2)=a^2+b^2-2a+2b+2$\dots\dots2 分}
\step{$=(a-1)^2+(b+1)^2\geqslant0$\dots\dots4 分}
\step{所以 $a^2+b^2\geqslant2a-2b-2$.}
\step{当且仅当 $a=1,b=-1$ 时，不等式中等号成立.\dots\dots7 分}
\step{（2）$a^2+4b^2=a^2+(2b)^2\geqslant2\cdot a\cdot(2b)=4ab=4$，\dots\dots5 分}
\step{当且仅当 $a=2b$，即 $\begin{cases}a=\sqrt2\\b=\dfrac{\sqrt2}{2}\end{cases}$ 或 $\begin{cases}a=-\sqrt2\\b=-\dfrac{\sqrt2}{2}\end{cases}$ 时，不等式中等号成立，}
\step{所以 $a^2+4b^2$ 的最小值为 4.\dots\dots7 分}
\solend
\fi

\item 已知关于 $x$ 的方程 $x^2+(m-3)x+2=0$，其中 $m\in\R$，有两个实数根 $\alpha$、$\beta$.

（1）若 $m=8$，求 $\dfrac1\alpha+\dfrac1\beta$ 的值；

（2）若 $\alpha^2+\beta^2=\alpha\beta+3$，求实数 $m$ 的值.

\ifanswer
\solbegin
\step{（1）若 $m=8$，则关于 $x$ 的方程为 $x^2+5x+2=0$，}
\step{则 $\alpha+\beta=-5$，$\alpha\beta=2$，所以 $\dfrac1\alpha+\dfrac1\beta=\dfrac{\alpha+\beta}{\alpha\beta}=-\dfrac52$；}
\step{（2）由题意得 $\alpha+\beta=-(m-3)=3-m$，$\alpha\beta=2$，}
\step{又 $\alpha^2+\beta^2=\alpha\beta+3$，则 $(\alpha+\beta)^2-2\alpha\beta=\alpha\beta+3$，}
\step{即 $(\alpha+\beta)^2=3\alpha\beta+3$，所以 $(3-m)^2=3\times2+3=9$，解得 $m=0$ 或 $6$.}
\solend
\fi

\item 设集合 $A=\{x\mid\lvert x-a\rvert<2\}$，$B=\left\{x\mid\dfrac{x+4}{x+2}>2\right\}$.

（1）若 $a=1$，试用区间表示集合 $A$、$B$；

（2）若 $A\cup\overline{B}=\R$，求实数 $a$ 的取值范围.

\ifanswer
\solbegin
\step{（1）当 $a=1$ 时，由 $\lvert x-1\rvert<2$ 得 $-2<x-1<2$，解得 $-1<x<3$，}
\step{所以 $A=(-1,3)$.}
\step{由 $\dfrac{x+4}{x+2}>2$ 得 $2-\dfrac{x+4}{x+2}=\dfrac{x}{x+2}<0$，则有 $x(x+2)<0$，解得 $-2<x<0$，}
\step{所以 $B=(-2,0)$；}
\step{（2）由 $\lvert x-a\rvert<2$ 得 $-2<x-a<2$，解得 $a-2<x<a+2$，}
\step{所以 $A=(a-2,a+2)$，}
\step{由（1）得 $B=(-2,0)$，则 $\overline{B}=\{x\mid x\geqslant0$ 或 $x\leqslant-2\}$，}
\step{由于 $A\cup\overline{B}=\R$，所以 $\begin{cases}a+2\geqslant0\\a-2\leqslant-2\end{cases}$，解得 $-2\leqslant a\leqslant0$，}
\step{所以实数 $a$ 的取值范围是 $[-2,0]$.}
\solend
\fi

\item 某服装公司生产的衬衫每件定价为 80 元，在某城市年销售 8 万件. 现该公司计划在该市招收代理商来销售衬衫，以降低管理和营销成本. 已知代理商收取的代理费为总销售金额的 $r\%$（其中 $0<r\leqslant10$）.

（1）若单件衬衫的利润保持不变，求招代理商后每件衬衫的定价为多少元（用 $r$ 表示）；

（2）已知将每件衬衫的价格提高到 $\dfrac{80}{1-r\%}$ 元后，每年的销售量会减少 $0.62r$ 万件. 若代理商每年收取的代理费不少于 16 万元，求 $r$ 的取值范围.

\ifanswer
\solbegin
\step{（1）设招代理商后每件衬衫的定价为 $x$ 元，代理商收取总销售金额的 $r\%$，}
\step{因此公司每件衬衫的实际净收入为 $x(1-r\%)$ 元，}
\step{单件利润保持不变，等价于净收入仍为 80 元，}
\step{于是 $x(1-r\%)=80$，解得 $x=\dfrac{80}{1-r\%}$，因为 $0<r\leqslant10$，有 $1-r\%\geqslant0.9>0$，}
\step{分母恒不为零，该定价有意义，所以招代理商后每件衬衫的定价为 $\dfrac{80}{1-r\%}$ 元；}
\step{（2）价格提高到 $\dfrac{80}{1-r\%}$ 元，年销售量减少 $0.62r$ 万件，}
\step{故年销售量为 $8-0.62r$ 万件，$8-0.62r>0$，}
\step{由题意得 $\dfrac{80}{1-r\%}\cdot(8-0.62r)\cdot r\%\geqslant16$，}
\step{即 $\dfrac{80}{1-\dfrac r{100}}\cdot(8-0.62r)\cdot\dfrac r{100}=\dfrac{80r(8-0.62r)}{100-r}\geqslant16$，}
\step{由 $0<r\leqslant10$ 得 $100-r>0$，整理得 $80r(8-0.62r)\geqslant16(100-r)$，}
\step{整理得 $31r^2-410r+1000\leqslant0$，解得 $\dfrac{100}{31}\leqslant r\leqslant10$，}
\step{当 $\dfrac{100}{31}\leqslant r\leqslant10$ 时，$8-0.62r>0$ 且 $1-r\%>0$，}
\step{均符合实际经营条件，所以 $r$ 的取值范围是 $\left[\dfrac{100}{31},10\right]$.}
\solend
\fi

\item 定义 $\min\{a_1,a_2,\cdots,a_n\}$ 为 $n$ 个实数 $a_1,a_2,\cdots,a_n$ 中的最小数，$\max\{a_1,a_2,\cdots,a_n\}$ 为 $n$ 个实数 $a_1,a_2,\cdots,a_n$ 中的最大数.

（1）设 $a,b$ 都是正实数，且 $a+b=1$，求 $\max\left\{ab,\dfrac14\right\}$；

（2）解不等式：$\min\{x+1,x^2+3,\lvert x-1\rvert\}>2x-3$；

（3）设 $a,b$ 都是正实数，求 $\max\left\{a+\dfrac1b,\dfrac2a+b\right\}$ 的最小值.

\ifanswer
\solbegin
\step{（1）$ab\leqslant\dfrac{(a+b)^2}{4}=\dfrac14$，所以 $\max\left\{ab,\dfrac14\right\}=\dfrac14$.}
\step{（2）因为 $x^2+3>x+1$，$x^2+3>\lvert x-1\rvert$，}
\step{所以 $\min\{x+1,x^2+3,\lvert x-1\rvert\}=\min\{x+1,\lvert x-1\rvert\}$.}
\step{当 $x\geqslant0$ 时，$\lvert x-1\rvert\leqslant x+1$，所以 $\lvert x-1\rvert>2x-3$，解得 $0\leqslant x<2$；}
\step{当 $x<0$ 时，$\lvert x-1\rvert>x+1$，所以 $x+1>2x-3$，解得 $x<0$.}
\step{综上，不等式的解集为 $(-\infty,2)$.}
\step{（3）设 $M=\max\left\{a+\dfrac1b,\dfrac2a+b\right\}$，则 $M\geqslant a+\dfrac1b$，$M\geqslant\dfrac2a+b$，}
\step{所以 $M^2\geqslant\left(a+\dfrac1b\right)\left(\dfrac2a+b\right)=3+ab+\dfrac2{ab}\geqslant3+2\sqrt2$，}
\step{所以 $M$ 的最小值为 $\sqrt2+1$，当 $a=\sqrt2$，$b=1$ 时取等号.}
\solend
\fi

\end{enumerate}

\end{document}
"
%
% 原始材料：微信公众号"魔都数学加油站"文章，作者破晓，2026-11-11 18:35 发布；
% 连载"27学年高一126"，原文标题"27学年高一126——2026-2027年上海市建青实验高一上期中"；
% 原文链接：https://mp.weixin.qq.com/s/HpxpBJu4d_wOnUCIbZSQiA。
% 正文为 8 张 PNG 页（01、05、08 为 2560x3620，02、03、04、06、07 为 4762x6735），
% 题目下印有【解析】或【答案】，为讲评版。题目、解析均据原图转录；
% 卷面粉红色水印"魔都数学加油站"及公众号页脚未录入。全卷无插图。
%
% 卷首标题：原卷印作"2026-2027 年建青实验高一上期中"（公众号标题中的"上海市"
% 卷面标题行没有，本稿照卷面），年份区间按全稿惯例排 en dash；原卷未印副标题，
% 本卷为期中卷、跨集合与常用逻辑用语／不等式／函数三章，按题目内容标为
% "集合与常用逻辑用语、不等式、函数"（与 125 同为期中卷，口径相同）：
% 集合与逻辑见 Q5、Q8、Q9、Q11、Q13、Q14、Q16、Q19，不等式见 Q6、Q7、Q10、
% Q17、Q20、Q21，函数（含指数对数、二次函数零点）见 Q3、Q4、Q12、Q15、Q18。
%
% 与原卷的排版差异：
%   ① 原文从第 3 题起，填空题编号为 3—12、选择题为 13—16、解答题为 17—21，本稿保留原题号；
%   ② 原卷每题下直接印【解析】／【答案】，本稿按全稿惯例置于答案版；第 4、5 题只印
%      【答案】，另提 \ans{}；第 13—16 题答案均在【解析】末句"故选 …"，题干括号内留空、
%      不另提 \ans{}；第 3 题印一行【解析】，按 \solbegin+\step 排；第 17 题解析保留
%      原卷评分点（照 高一54／64 成例排 \dots\dots N 分）；
%   ③ 原卷未印副标题，本稿按"知识模块"口径补出（见上"卷首标题"）；
%   ④ 原卷填空横线按全稿惯例排为 \blankline[4.4em]；选择题括号排 \mbox{（\quad）}、
%      选项段前按惯例加 \keepwithprev；第 14 题原卷括号另起一行，本稿并入题干行末；
%      解答题不留作答区（原卷为讲评版，与 109—125 同）；
%   ⑤ 原卷数集排斜体 $R$，按全稿惯例改 \R；绝对值排 \lvert…\rvert，集合描述竖线排 \mid；
%      补集记号照原卷上划线排 \overline{B}；不等号按全稿惯例排 \leqslant／\geqslant；
%      子题标号原卷一律用全角括号（1）（2）（3），本稿照录（\step 行内亦同，照 高一121 成例）。
%
% 判别式字形：第 10 题解析 $\Delta=a^2-4a<0$ 原卷排 $\Delta$（等腰、高），
%   照录 \Delta（与 高一124／125 同题排 $\triangle$（等边、扁）不同，逐处照本卷字形）。
%
% 子集符号：第 8 题题干与解析两处 $S\subset P$ 均印单线、无下划线，照录 \subset；
%   全卷未出现带下划线的 \subseteq。
%
% 原卷存疑照录：
%   第 10 题题干作 $ax^2-ax+1>0$，解析却作 $ax^2+ax+1>0$（一次项符号不一致），照录、未统一；
%   第 11 题解析"此时 $M=\dfrac32,1,\dfrac12$，满足"未用集合花括号（题干 $M$ 为集合），照录；
%   第 11 题解析"因为它们是集合 $M$ 中的元素，所以存在互异性"（"存在互异性"疑为
%     "满足互异性"之误），又"当 $x_1\neq x_2\neq x_3$ 时"连写两个不等号，均照录；
%   第 16 题解析两处"惟如此才能保证"（用"惟"字）及"符合定义的参数 $m$ 的值一定大于等于 1
%     或小于等于 0"的措辞，照录；
%   第 17 题 (1) 解析"当且仅当 $a=1,b=-1$ 时"用半角逗号、第 21 题 (3) 解析
%     "当 $a=\sqrt2$，$b=1$ 时取等号"用全角逗号，均照录；
%   第 19 题解析由 $\dfrac{x+4}{x+2}>2$ 得 "$2-\dfrac{x+4}{x+2}=\dfrac{x}{x+2}<0$"
%     （以 2 减左式，与常见的"左式减 2"次序相反），照录；
%   第 15 题选项 A、B 作"该方程…"、C、D 作"方程组…"，照录；
%   第 20 题题干作"招收代理商"，(1) 却作"招代理商后"，照录。
%
% 行数复核：逐题数 \step 条数，合计 91 条（Q3：1；Q4：0；Q5：0；Q6：3；Q7：2；
% Q8：5；Q9：5；Q10：4；Q11：8；Q12：3；Q13：3；Q14：2；Q15：3；Q16：9；
% Q17：7；Q18：5；Q19：9；Q20：13；Q21：9），与原卷印刷行逐题相等
% （cases 块作一个印刷行计）。第 4、5 题只印【答案】，提 \ans{}，不计入 \step 账；
% 四道选择题答案均在解析末句，括号内留空、不另提 \ans{}。
%
\documentclass[a4paper,11pt]{article}
\usepackage{common}

\begin{document}

\examtitlest[3]{2026--2027 年建青实验高一上期中}{集合与常用逻辑用语、不等式、函数}

\bigsec{一、填空题}

\begin{enumerate}
\setcounter{enumi}{2}

\item 已知 $\log_2 x=3$，则 $x=$\blankline[4.4em].

\ifanswer
\solbegin
\step{$x=2^3=8$.}
\solend
\fi

\item 将 $\sqrt{a^2\sqrt a}$（其中 $a>0$）化为有理数指数幂的形式为\blankline[4.4em].

\ifanswer
\ans{$a^{\frac54}$}
\fi

\item 用反证法证明命题：“已知 $m$、$n\in\R$，若 $m+n\leqslant5$，则 $m\leqslant2$ 或 $n\leqslant3$” 时应假设\blankline[4.4em].

\ifanswer
\ans{$m>2$ 且 $n>3$}
\fi

\item 若 $x\in\R$，则 $\lvert x+1\rvert+\lvert x-1\rvert$ 的最小值为\blankline[4.4em].

\ifanswer
\solbegin
\step{$\lvert x+1\rvert+\lvert x-1\rvert\geqslant\lvert(x+1)-(x-1)\rvert=2$，}
\step{当且仅当 $(x+1)(x-1)\leqslant0$，即 $-1\leqslant x\leqslant1$ 时取等号，}
\step{所以 $\lvert x+1\rvert+\lvert x-1\rvert$ 的最小值为 2.}
\solend
\fi

\item 已知 $x>1$，则 $x+\dfrac1{x-1}$ 的最小值是\blankline[4.4em].

\ifanswer
\solbegin
\step{$x-1>0$，则 $x+\dfrac1{x-1}=(x-1)+\dfrac1{x-1}+1\geqslant2\sqrt{(x-1)\cdot\dfrac1{x-1}}+1=3$，}
\step{当且仅当 $x=2$ 时取等号，故 $x+\dfrac1{x-1}$ 的最小值为 3.}
\solend
\fi

\item 已知集合 $P=\{x\mid x^2+x-6=0\}$，$S=\{x\mid ax+1=0\}$ 且 $S\subset P$，则 $a$ 的值为\blankline[4.4em].

\ifanswer
\solbegin
\step{集合 $P=\{-3,2\}$，集合 $S$ 中至多有一个元素，}
\step{若集合 $S$ 为空集，即 $a=0$ 时，显然满足条件 $S\subset P$，故 $a=0$.}
\step{若集合 $S$ 非空集，即 $a\neq0$，此时 $S=\left\{-\dfrac1a\right\}$，}
\step{若 $-\dfrac1a=-3$，则 $a=\dfrac13$，若 $-\dfrac1a=2$，则 $a=-\dfrac12$；}
\step{故 $a$ 的取值为 $0$、$\dfrac13$、$-\dfrac12$.}
\solend
\fi

\item 已知 $p:\begin{cases}a>4\\b>4\end{cases}$，$q:\begin{cases}a+b>8\\ab>16\end{cases}$，则 $p$ 是 $q$ 的\blankline[4.4em]条件.

\ifanswer
\solbegin
\step{若 $a>4$，$b>4$，则 $a+b>4+4=8$，$ab>4\times4=16$，}
\step{所以 $p$ 成立一定能推出 $q$ 成立，充分性成立，}
\step{取 $a=2$，$b=9$，$a+b=11>8$，$ab=18>16$，满足 $q$；}
\step{但是 $a=2<4$，不满足 $p$，所以由 $q$ 不能推出 $p$，必要性不成立，}
\step{综上所述，$p$ 是 $q$ 的充分不必要条件.}
\solend
\fi

\item 若关于 $x$ 的不等式 $ax^2-ax+1>0$ 对一切实数 $x$ 恒成立，则实数 $a$ 的取值范围是\blankline[4.4em].

\ifanswer
\solbegin
\step{当 $a=0$ 时，不等式即 $1>0$，满足条件.}
\step{当 $a\neq0$ 时，要使不等式 $ax^2+ax+1>0$ 对一切 $x\in\R$ 恒成立，}
\step{需 $\begin{cases}a>0\\\Delta=a^2-4a<0\end{cases}$，解得 $0<a<4$.}
\step{综上，实数 $a$ 的取值范围是 $[0,4)$.}
\solend
\fi

\item 已知集合 $M=\{x\mid (x-a)(x^2-ax+a-1)=0\}$ 各元素之和等于 3，则实数 $a=$\blankline[4.4em].

\ifanswer
\solbegin
\step{方程 $(x-a)(x^2-ax+a-1)=0$，化简得 $(x-a)(x-1)[x-(a-1)]=0$，}
\step{解得 $x_1=a$，$x_2=1$，$x_3=a-1$，}
\step{因为它们是集合 $M$ 中的元素，所以存在互异性.}
\step{当 $x_1\neq x_2\neq x_3$ 时，即 $a+1+a-1=3$，解得 $a=\dfrac32$，此时 $M=\dfrac32,1,\dfrac12$，满足.}
\step{当 $a=1$ 时，$1+0\neq3$，不合题意；}
\step{当 $a=a-1$ 时，不符合题意；}
\step{当 $1=a-1$ 时，$a=2$，此时 $M=\{2,1\}$ 满足题意.}
\step{故 $a$ 的取值为 $2$ 或 $\dfrac32$.}
\solend
\fi

\item 已知 $a>b>1$ 且 $a^2-6a-2m+1=0$，$b^2-6b-2m+1=0$，则实数 $m$ 的取值范围是\blankline[4.4em].

\ifanswer
\solbegin
\step{由题意得 $a$、$b$ 是函数 $f(x)=x^2-6x-2m+1$ 两个大于 1 的不同零点，}
\step{所以 $\begin{cases}f(1)>0\\f(3)<0\end{cases}$，即 $\begin{cases}-4-2m>0\\-8-2m<0\end{cases}$，解得 $-4<m<-2$，}
\step{故 $m$ 的取值范围为 $(-4,-2)$.}
\solend
\fi

\end{enumerate}

\bigsec{二、选择题}

\begin{enumerate}
\setcounter{enumi}{12}

\item “$x<2$” 是 “$x^2<4$” 的\mbox{（\quad）}条件

\keepwithprev
A. 充分不必要\qquad\qquad B. 必要不充分

C. 充要\qquad\qquad\qquad\quad D. 既不充分也不必要

\ifanswer
\solbegin
\step{解 $x^2<4$，得 $-2<x<2$，}
\step{由 $x<2$ 得不出 $-2<x<2$，而由 $-2<x<2$ 能得出 $x<2$，}
\step{所以 “$x<2$” 是 “$x^2<4$” 的必要不充分条件. 故选 $B$.}
\solend
\fi

\item 已知集合 $A=\{x\mid x<-2\}$，$B=\{x\mid x\geqslant a\}$，且 $A\cup B=\R$，则实数 $a$ 的取值范围是\mbox{（\quad）}

\keepwithprev
A. $a\leqslant-2$\qquad\qquad B. $a<-2$\qquad\qquad C. $a\geqslant-2$\qquad\qquad D. $a>-2$

\ifanswer
\solbegin
\step{集合 $A=\{x\mid x<-2\}$，$B=\{x\mid x\geqslant a\}$，且 $A\cup B=\R$，}
\step{则实数 $a$ 的取值范围是 $a\leqslant-2$. 故选 $A$.}
\solend
\fi

\item 下列关于 $x,y$ 的方程组 $\begin{cases}y=kx+1\\y=2kx+3\end{cases}$ 的解的说法中正确的是\mbox{（\quad）}

\keepwithprev
A. 该方程一定有唯一解\qquad\qquad B. 该方程没有解

C. $k\neq0$ 时，方程组有唯一解\qquad\qquad D. $k=0$ 时，方程组有无数组解

\ifanswer
\solbegin
\step{由题意得 $kx+1=2kx+3$，即 $kx=-2$，}
\step{当 $k=0$ 时，$0=-2$ 不成立，方程组无解；}
\step{当 $k\neq0$ 时，$x=-\dfrac2k$，方程组有唯一解. 故选 $C$.}
\solend
\fi

\item 设非空集合 $S=\{x\mid m\leqslant x\leqslant n\}$ 满足：当 $x\in S$ 时，有 $x^2\in S$. 给出如下三个结论：①若 $m=1$，则 $S=\{1\}$；②若 $m=-\dfrac12$，则 $\dfrac14\leqslant n\leqslant1$；③若 $n=\dfrac12$，则 $-\dfrac{\sqrt2}{2}\leqslant m\leqslant0$. 其中正确命题的个数为\mbox{（\quad）}

\keepwithprev
A. $3$\qquad\qquad B. $2$\qquad\qquad C. $1$\qquad\qquad D. $0$

\ifanswer
\solbegin
\step{由定义设非空集合 $S=\{x\mid m\leqslant x\leqslant n\}$ 满足当 $x\in S$ 时，有 $x^2\in S$，}
\step{符合定义的参数 $m$ 的值一定大于等于 1 或小于等于 0，}
\step{惟如此才能保证 $m\in S$ 时，有 $m^2\in S$ 即 $m^2\geqslant m$，}
\step{符合条件的 $n$ 的值一定大于等于 0，小于等于 1，}
\step{惟如此才能保证 $n\in S$ 时，有 $n^2\in S$ 即 $n^2\leqslant n$，对各个命题进行判断：}
\step{对于① $m=1$，$m^2=1\in S$，故必有 $\begin{cases}n^2\leqslant n\\n\geqslant1\end{cases}$，得 $n=1$，$S=\{1\}$，}
\step{② $m=-\dfrac12$，$m^2=\dfrac14\in S$，则 $\begin{cases}n^2\leqslant n\\\dfrac14\leqslant n\end{cases}$，解得 $\dfrac14\leqslant n\leqslant1$；}
\step{对于③，若 $n=\dfrac12$，则 $\begin{cases}\dfrac12\geqslant m\\m^2\geqslant m\\\dfrac12\geqslant m^2\end{cases}$，解得 $-\dfrac{\sqrt2}{2}\leqslant m\leqslant0$，}
\step{所以正确命题有 3 个. 故选 $A$.}
\solend
\fi

\end{enumerate}

\bigsec{三、解答题}

\begin{enumerate}
\setcounter{enumi}{16}

\item 已知 $a$、$b$ 是实数.

（1）求证：$a^2+b^2\geqslant2a-2b-2$，并指出等号成立的条件；

（2）若 $ab=1$，求 $a^2+4b^2$ 的最小值.

\ifanswer
\solbegin
\step{（1）因为 $a^2+b^2-(2a-2b-2)=a^2+b^2-2a+2b+2$\dots\dots2 分}
\step{$=(a-1)^2+(b+1)^2\geqslant0$\dots\dots4 分}
\step{所以 $a^2+b^2\geqslant2a-2b-2$.}
\step{当且仅当 $a=1,b=-1$ 时，不等式中等号成立.\dots\dots7 分}
\step{（2）$a^2+4b^2=a^2+(2b)^2\geqslant2\cdot a\cdot(2b)=4ab=4$，\dots\dots5 分}
\step{当且仅当 $a=2b$，即 $\begin{cases}a=\sqrt2\\b=\dfrac{\sqrt2}{2}\end{cases}$ 或 $\begin{cases}a=-\sqrt2\\b=-\dfrac{\sqrt2}{2}\end{cases}$ 时，不等式中等号成立，}
\step{所以 $a^2+4b^2$ 的最小值为 4.\dots\dots7 分}
\solend
\fi

\item 已知关于 $x$ 的方程 $x^2+(m-3)x+2=0$，其中 $m\in\R$，有两个实数根 $\alpha$、$\beta$.

（1）若 $m=8$，求 $\dfrac1\alpha+\dfrac1\beta$ 的值；

（2）若 $\alpha^2+\beta^2=\alpha\beta+3$，求实数 $m$ 的值.

\ifanswer
\solbegin
\step{（1）若 $m=8$，则关于 $x$ 的方程为 $x^2+5x+2=0$，}
\step{则 $\alpha+\beta=-5$，$\alpha\beta=2$，所以 $\dfrac1\alpha+\dfrac1\beta=\dfrac{\alpha+\beta}{\alpha\beta}=-\dfrac52$；}
\step{（2）由题意得 $\alpha+\beta=-(m-3)=3-m$，$\alpha\beta=2$，}
\step{又 $\alpha^2+\beta^2=\alpha\beta+3$，则 $(\alpha+\beta)^2-2\alpha\beta=\alpha\beta+3$，}
\step{即 $(\alpha+\beta)^2=3\alpha\beta+3$，所以 $(3-m)^2=3\times2+3=9$，解得 $m=0$ 或 $6$.}
\solend
\fi

\item 设集合 $A=\{x\mid\lvert x-a\rvert<2\}$，$B=\left\{x\mid\dfrac{x+4}{x+2}>2\right\}$.

（1）若 $a=1$，试用区间表示集合 $A$、$B$；

（2）若 $A\cup\overline{B}=\R$，求实数 $a$ 的取值范围.

\ifanswer
\solbegin
\step{（1）当 $a=1$ 时，由 $\lvert x-1\rvert<2$ 得 $-2<x-1<2$，解得 $-1<x<3$，}
\step{所以 $A=(-1,3)$.}
\step{由 $\dfrac{x+4}{x+2}>2$ 得 $2-\dfrac{x+4}{x+2}=\dfrac{x}{x+2}<0$，则有 $x(x+2)<0$，解得 $-2<x<0$，}
\step{所以 $B=(-2,0)$；}
\step{（2）由 $\lvert x-a\rvert<2$ 得 $-2<x-a<2$，解得 $a-2<x<a+2$，}
\step{所以 $A=(a-2,a+2)$，}
\step{由（1）得 $B=(-2,0)$，则 $\overline{B}=\{x\mid x\geqslant0$ 或 $x\leqslant-2\}$，}
\step{由于 $A\cup\overline{B}=\R$，所以 $\begin{cases}a+2\geqslant0\\a-2\leqslant-2\end{cases}$，解得 $-2\leqslant a\leqslant0$，}
\step{所以实数 $a$ 的取值范围是 $[-2,0]$.}
\solend
\fi

\item 某服装公司生产的衬衫每件定价为 80 元，在某城市年销售 8 万件. 现该公司计划在该市招收代理商来销售衬衫，以降低管理和营销成本. 已知代理商收取的代理费为总销售金额的 $r\%$（其中 $0<r\leqslant10$）.

（1）若单件衬衫的利润保持不变，求招代理商后每件衬衫的定价为多少元（用 $r$ 表示）；

（2）已知将每件衬衫的价格提高到 $\dfrac{80}{1-r\%}$ 元后，每年的销售量会减少 $0.62r$ 万件. 若代理商每年收取的代理费不少于 16 万元，求 $r$ 的取值范围.

\ifanswer
\solbegin
\step{（1）设招代理商后每件衬衫的定价为 $x$ 元，代理商收取总销售金额的 $r\%$，}
\step{因此公司每件衬衫的实际净收入为 $x(1-r\%)$ 元，}
\step{单件利润保持不变，等价于净收入仍为 80 元，}
\step{于是 $x(1-r\%)=80$，解得 $x=\dfrac{80}{1-r\%}$，因为 $0<r\leqslant10$，有 $1-r\%\geqslant0.9>0$，}
\step{分母恒不为零，该定价有意义，所以招代理商后每件衬衫的定价为 $\dfrac{80}{1-r\%}$ 元；}
\step{（2）价格提高到 $\dfrac{80}{1-r\%}$ 元，年销售量减少 $0.62r$ 万件，}
\step{故年销售量为 $8-0.62r$ 万件，$8-0.62r>0$，}
\step{由题意得 $\dfrac{80}{1-r\%}\cdot(8-0.62r)\cdot r\%\geqslant16$，}
\step{即 $\dfrac{80}{1-\dfrac r{100}}\cdot(8-0.62r)\cdot\dfrac r{100}=\dfrac{80r(8-0.62r)}{100-r}\geqslant16$，}
\step{由 $0<r\leqslant10$ 得 $100-r>0$，整理得 $80r(8-0.62r)\geqslant16(100-r)$，}
\step{整理得 $31r^2-410r+1000\leqslant0$，解得 $\dfrac{100}{31}\leqslant r\leqslant10$，}
\step{当 $\dfrac{100}{31}\leqslant r\leqslant10$ 时，$8-0.62r>0$ 且 $1-r\%>0$，}
\step{均符合实际经营条件，所以 $r$ 的取值范围是 $\left[\dfrac{100}{31},10\right]$.}
\solend
\fi

\item 定义 $\min\{a_1,a_2,\cdots,a_n\}$ 为 $n$ 个实数 $a_1,a_2,\cdots,a_n$ 中的最小数，$\max\{a_1,a_2,\cdots,a_n\}$ 为 $n$ 个实数 $a_1,a_2,\cdots,a_n$ 中的最大数.

（1）设 $a,b$ 都是正实数，且 $a+b=1$，求 $\max\left\{ab,\dfrac14\right\}$；

（2）解不等式：$\min\{x+1,x^2+3,\lvert x-1\rvert\}>2x-3$；

（3）设 $a,b$ 都是正实数，求 $\max\left\{a+\dfrac1b,\dfrac2a+b\right\}$ 的最小值.

\ifanswer
\solbegin
\step{（1）$ab\leqslant\dfrac{(a+b)^2}{4}=\dfrac14$，所以 $\max\left\{ab,\dfrac14\right\}=\dfrac14$.}
\step{（2）因为 $x^2+3>x+1$，$x^2+3>\lvert x-1\rvert$，}
\step{所以 $\min\{x+1,x^2+3,\lvert x-1\rvert\}=\min\{x+1,\lvert x-1\rvert\}$.}
\step{当 $x\geqslant0$ 时，$\lvert x-1\rvert\leqslant x+1$，所以 $\lvert x-1\rvert>2x-3$，解得 $0\leqslant x<2$；}
\step{当 $x<0$ 时，$\lvert x-1\rvert>x+1$，所以 $x+1>2x-3$，解得 $x<0$.}
\step{综上，不等式的解集为 $(-\infty,2)$.}
\step{（3）设 $M=\max\left\{a+\dfrac1b,\dfrac2a+b\right\}$，则 $M\geqslant a+\dfrac1b$，$M\geqslant\dfrac2a+b$，}
\step{所以 $M^2\geqslant\left(a+\dfrac1b\right)\left(\dfrac2a+b\right)=3+ab+\dfrac2{ab}\geqslant3+2\sqrt2$，}
\step{所以 $M$ 的最小值为 $\sqrt2+1$，当 $a=\sqrt2$，$b=1$ 时取等号.}
\solend
\fi

\end{enumerate}

\end{document}
